\documentclass[11pt]{amsart}
\usepackage[margin=1.15in]{geometry}
\usepackage{amssymb,amsmath,amsthm,mathtools}
\usepackage{needspace,microtype,etoolbox}
\usepackage{listings,xcolor,textcomp}
\definecolor{LeanKeyword}{HTML}{254E91}
\definecolor{LeanType}{HTML}{006B65}
\definecolor{LeanComment}{HTML}{526B46}
\definecolor{LeanString}{HTML}{8A4727}
\lstdefinelanguage{Lean}{
  sensitive=true,
  morekeywords={abbrev,by,class,def,end,exact,fun,import,in,instance,let,
    namespace,noncomputable,open,scoped,section,theorem,variable,where},
  morekeywords=[2]{Prop,Type,Set,NontriviallyNormedField,NormedAddCommGroup,
    NormedSpace,TopologicalSpace,ChartedSpace,ModelWithCorners,IsOpen,
    ContMDiffOn,ContMDiffVectorBundle},
  morecomment=[l]{--},morecomment=[s]{/-}{-/},morestring=[b]"
}
% Render Lean's Unicode notation with both the editor and PDF export engines.
\lstnewenvironment{leancode}{\lstset{
  language=Lean,
  basicstyle=\small\ttfamily,columns=fullflexible,keepspaces=true,
  keywordstyle=\color{LeanKeyword}\bfseries,
  keywordstyle=[2]\color{LeanType},
  commentstyle=\color{LeanComment}\itshape,
  stringstyle=\color{LeanString},
  showstringspaces=false,upquote=true,
  literate={∀}{{\ensuremath{\forall}}}1
           {γ}{{\ensuremath{\gamma}}}1
           {→}{{\ensuremath{\to}}}1
           {↦}{{\ensuremath{\mapsto}}}1
           {ℕ}{{\ensuremath{\mathbb N}}}1
           {∞}{{\ensuremath{\infty}}}1
           {ω}{{\ensuremath{\omega}}}1
           {×}{{\ensuremath{\times}}}1
           {⋀}{{\ensuremath{\bigwedge}}}1
           {𝓘}{{\ensuremath{\mathcal I}}}1
           {₁}{{\ensuremath{{}_1}}}1
           {₂}{{\ensuremath{{}_2}}}1
}}{}
\usepackage[colorlinks=true,linkcolor=blue,urlcolor=blue,citecolor=blue]{hyperref}
\hypersetup{pdftitle={Alternating bundles and analytic descent},
 pdfauthor={Jack McCarthy},pdfsubject={Alternating bundles over nontrivially normed fields},
 pdfkeywords={alternating maps, vector bundles, analytic descent, nonarchimedean analysis, Lean}}
\newtheorem{theorem}{Theorem}[section]
\newtheorem{lemma}[theorem]{Lemma}
\newtheorem{proposition}[theorem]{Proposition}
\newtheorem{corollary}[theorem]{Corollary}
\theoremstyle{definition}
\newtheorem{definition}[theorem]{Definition}
\theoremstyle{remark}
\newtheorem{remark}[theorem]{Remark}
\numberwithin{equation}{section}
\newcommand{\N}{\mathbb N}
\newcommand{\Z}{\mathbb Z}
\newcommand{\F}{\mathbb F}
\newcommand{\cL}{\mathcal L}
\newcommand{\Alt}{\operatorname{Alt}}
\newcommand{\Mult}{\operatorname{Mult}}
\newcommand{\End}{\operatorname{End}}
\newcommand{\linf}{\ell^\infty}
\newcommand{\sgn}{\operatorname{sgn}}
\newcommand{\spn}{\operatorname{span}}
\newcommand{\id}{\mathrm{id}}
\newcommand{\abs}[1]{\lvert #1\rvert}
\newcommand{\norm}[1]{\lVert #1\rVert}
\newcommand{\cC}{\mathcal{C}}
\newcommand{\cD}{\mathcal{D}}
\newcommand{\coeff}{\operatorname{coeff}}
\newcommand{\ord}{\operatorname{ord}}
\newcommand{\rank}{\operatorname{rank}}
\newcommand{\type}{\operatorname{type}}
\newcommand{\sdim}{\operatorname{sdim}}
\newcommand{\incl}{\mathrm{incl}}
\newcommand{\Cop}{C^{\mathrm{op}}}
\newcommand{\Cvec}{C^{\mathrm{vec}}}
\newcommand{\Kh}{\widehat K}
\setlength{\emergencystretch}{1em}
\title{Alternating bundles and analytic descent}
\author{Jack McCarthy}
\thanks{The original counterexample and its initial write-up were developed with Anthropic's Claude.
Subsequent research and exposition used Claude and OpenAI's Codex; the organization of the paper was
prepared with Codex. AI assistants contributed to mathematical arguments and independent proof audits.
The original Lean formalization used Codex and Autoform. The remaining claims and the per-claim
ledger were formalized with Anthropic's Claude Code in October 2026. Proposition~\ref{prop:tangent-values}
was developed and written with Claude and is not formalized.
Section~\ref{sec:verification} records the verification scope.}
\makeatletter
% Keep the complete front matter on one page without changing the body layout.
\patchcmd{\@maketitle}{\global\topskip42\p@}{\global\topskip30\p@}{}%
  {\PackageError{charp}{Unable to adjust title spacing}{Check the amsart title definition.}}
\patchcmd{\@maketitle}{\dimen@34\p@}{\dimen@10\p@}{}%
  {\PackageError{charp}{Unable to adjust abstract spacing}{Check the amsart title definition.}}
\patchcmd{\@setauthors}{\@topsep30\p@}{\@topsep18\p@}{}%
  {\PackageError{charp}{Unable to adjust author spacing}{Check the amsart author definition.}}
\g@addto@macro\@setauthors{\vskip6pt{\centering\footnotesize October 2026\par}}
\makeatother
\begin{document}
\begin{abstract}
Regular functors on normed spaces induce functors on vector bundles, and natural transformations pass
fiberwise to bundles. We study this principle for bounded alternating maps over nontrivially normed
fields. Their inclusion into all multilinear maps reduces the analytic question to one problem:
lifting homogeneous coefficients through a closed linear subspace with controlled norms.
We give an exact descent criterion and two mechanisms for satisfying it, using retractions in the
fibers or coordinates on the parameter space.

The alternating bifunctor is always $C^\infty$ and is analytic on all hom spaces exactly when $k!\ne0$ in the field.
Nonarchimedean spherically complete targets give analyticity in every degree. Independently, bases
with open $K^d$ charts support analytic alternating bundles with arbitrary normed fibers in every
characteristic; summable infinite-dimensional parameter spaces give further cases.
Counterexamples show the limits of both general and scalar-valued constructions, and polynomial
shears realize the failures in bundle and tangent transitions.
The exterior-power representation identifies the universal action and its scalar bidual test.
Finally, over complete fields, ambient analytic differential forms retain pullback, wedge products
and exterior differentiation even when the natural alternating atlas is not analytic.
The counterexamples and supplementary classification results are proved in the appendices.
Every theorem, proposition, lemma and corollary except Proposition~\ref{prop:tangent-values} has a
Lean proof of a corresponding statement; Section~\ref{sec:verification} lists the main places where those statements
are narrower than the text.
\end{abstract}
\maketitle
\setcounter{tocdepth}{1}
\begingroup
\small
\linespacing=8pt
\tableofcontents\removelastskip
\endgroup
\clearpage

\section{Introduction}\label{sec:intro}

Linear algebra supplies differential geometry with dual bundles, tensor bundles, and bundles of
differential forms. These constructions reflect a general principle: a sufficiently regular functor
on normed spaces extends fiberwise to vector bundles and their morphisms, and natural transformations
extend with it. The analytic requirement concerns the action on linear maps: $C^n$ families must
give rise to $C^n$ families of induced maps. Every $C^n$ functor on normed spaces satisfies this
condition. Thus, for a fixed $C^n$ manifold $M$, fiberwise application defines a functor
\[
\begin{gathered}
 \{C^n\text{ functors on normed spaces}\}
 \longrightarrow
 \{\text{functors on }C^n\text{ vector bundles over }M\},\\
 \mathcal F\longmapsto\mathcal F_M.
\end{gathered}
\]
Its morphisms are natural transformations: $\eta\colon\mathcal F\Rightarrow\mathcal G$ gives
$\eta_M\colon\mathcal F_M\Rightarrow\mathcal G_M$. The same construction applies to several variables
and to contravariant variables. It is this passage, rather than the construction of fibers alone, that
makes regularity of a linear-algebraic functor a geometric question.

We study the alternating-map bifunctor
\[
 \Alt^k\colon\mathsf{Vec}_K^{\mathrm{op}}\times\mathsf{Vec}_K\longrightarrow\mathsf{Vec}_K,
 \qquad (E,F)\longmapsto\Alt^k(E;F),
\]
where $\mathsf{Vec}_K$ denotes normed spaces and bounded linear maps over a nontrivially normed field
$K$. The spaces need not be complete, and the degree $k$ is any nonnegative integer. Our main question
is when applying this bifunctor to analytic input bundles produces an analytic bundle, naturally in
the inputs and their morphisms.

The answer is governed by one local problem. Alternating maps form a closed linear subspace of all
bounded multilinear maps. The multilinear construction is always analytic. Its restriction to
alternating maps is analytic precisely when the homogeneous coefficients can be represented by bounded
multilinear maps with alternating values and suitable norm bounds. We call this passage
\emph{analytic descent through a closed linear inclusion}. Theorem~\ref{thm:coefficient-descent}
gives an exact criterion. It separates two ways to obtain alternating bundles: split the inclusion on
the fiber side, or construct the coefficients using the coordinates of the base.

\subsection*{The principal conclusions}
The bifunctor $\Alt^k$ is always $C^\infty$. It is analytic on every hom space if and only if $k!\ne0$
in $K$, and this classification is unchanged on the full subcategory $\mathsf{Ban}_K$ of complete
spaces. Thus every characteristic-zero field is covered in every degree. In characteristic $p>0$,
the restriction to degrees $k<p$ is sharp for arbitrary normed targets. In positive characteristic
there is also an analytic bifunctor in every degree
\[
 \Alt^k\colon\mathsf{Vec}_K^{\mathrm{op}}\times\mathsf{Vec}_K^\circ
 \longrightarrow\mathsf{Vec}_K^\circ,
\]
where $\mathsf{Vec}_K^\circ$ is the full subcategory of nonarchimedean, spherically complete spaces.
The last assertion includes closure of that target category under alternating-map spaces
(Theorem~\ref{thm:spherical}).

For bundles the condition is weaker. An analytic base with open $K^d$ charts tests only
finite-coordinate families of operators. Every such family is preserved by $\Alt^k$, in every
characteristic and degree, with arbitrary normed fibers. Consequently $\Alt^k$ induces a bifunctor
on analytic normed bundles over every such base, including their operator-valued analytic morphisms.
There are also infinite-dimensional parameter classes for which descent holds, with explicit
hypotheses on the final target fiber (Theorem~\ref{thm:summable-parameters}). A single bundle theorem,
Corollary~\ref{cor:bundle-cases}, collects these cases.

The distinction between whole hom spaces and the families tested by a base is essential. If $k!=0$
in $K$, analytic input bundles over a suitable Banach base can have a naturally induced alternating
atlas with a nowhere analytic transition map. An affine family of invertible operators realizes this
failure, so restricting to invertible transitions does not remove the obstruction. For scalar forms
there is an exact classification over complete nonarchimedean fields:
\[
 \Alt^k(-;K)\text{ is analytic on all normed sources}
 \quad\Longleftrightarrow\quad
 k!\ne0\ \text{or}\ K\text{ is spherically complete}.
\]
In the negative case, even the alternating cotangent atlas of an analytic Banach manifold can fail.
With Banach or tangent values the corresponding atlas can fail over every field of characteristic $p$ once $k\ge p$
(Proposition~\ref{prop:tangent-values}).
Section~\ref{sec:sharpness} proves the geometric reductions from the two operator obstructions.

An analytic calculus of forms nevertheless survives. Over every complete field, alternating-valued
analytic sections of the multilinear cotangent bundle admit pullback, shuffle products, and exterior
differentiation, and form a differential graded algebra. They agree with analytic sections of the
alternating bundle under the descent hypotheses stated in Section~\ref{sec:forms}. The distinction is
substantive: a polynomial pullback can preserve ambient analyticity while destroying intrinsic
alternating-valued analyticity.

\subsection*{Method and scope}
The categorical construction is proved once in Section~\ref{sec:bundles}. Section~\ref{sec:descent}
isolates its analytic obstruction, and Section~\ref{sec:positive} proves the mechanisms that remove it.
Section~\ref{sec:duality} identifies the representing exterior power and explains what scalar and
dual-valued forms detect. The geometric obstructions and the resulting calculus occupy
Sections~\ref{sec:sharpness} and~\ref{sec:forms}.
Section~\ref{sec:mathlib-interface} proposes a Mathlib interface with one alternating-bundle instance, backed by the family-preservation theorems.

The generality in the fibers matters. The ordinary projective exterior power represents alternating
maps with arbitrary normed targets; the nonarchimedean projective exterior power represents complete
nonarchimedean targets. Their analyticity properties can differ. Completeness assumptions below are
attached to the arguments that need them, rather than imposed on the category at the outset.

Our bundle construction follows Lang's functors of class $C^n$~\cite[Ch.~III, \S4]{La}.
Gou\"ezel's characteristic-zero alternating-map construction~\cite{G1,G2} provides the factorial case,
and closed-range reflection of smoothness is the argument formalized in~\cite{PR}.
Bourbaki states the alternating bundle under the hypothesis that the field has characteristic zero or that
the vector bundles have finite dimension~\cite[7.8]{Bou}, as quoted in~\cite{G1}; Gou\"ezel notes there that the
finite-dimensional case works in positive characteristic, which over a complete field is Proposition~\ref{prop:split-lift}(2).
The nonarchimedean extension mechanism belongs to the functional analysis of spherical completeness;
we use its norm-controlled form and relate it to~\cite{In,Schi,Sch,Yuan}.

The long constructions are proved in the appendices. Appendices~\ref{app:FF}--\ref{app:assembly}
establish the general operator obstruction, and Appendix~\ref{app:scalar} proves the scalar obstruction.
Appendix~\ref{app:reflection} characterizes universal analytic reflection by tensor projections;
Appendix~\ref{app:no-largest} proves that a largest full analytic domain need not exist.
Further coordinate methods and a nonspherical analytic target appear in
Appendix~\ref{app:coordinate-methods}. These results supplement the common descent framework without
interrupting the bundle arguments. Section~\ref{sec:verification} distinguishes the Lean-checked
statements from those with written proofs only; a proof in this article is not by itself a claim
of formal verification.

\subsection*{Conventions}
The field $K$ is nontrivially normed throughout. We write $\cL(E,F)$ for bounded linear maps with their
operator norm, $\Mult^k(E;F)$ for bounded $k$-linear maps, and $\Alt^k(E;F)$ for those vanishing whenever
two arguments coincide. This last condition, rather than antisymmetry alone, defines alternation in
characteristic two. The degree-zero space is $F$. Products of normed spaces carry the maximum norm.
The letters $E,D,F$ usually denote fibers; $P$ denotes the model space of a base or a parameter domain.
\emph{Finite coordinates} means a bounded linear isomorphism with $K^d$. Over a complete field this is
equivalent to finite algebraic dimension; over an incomplete field the coordinate condition is retained.

A norm is \emph{nonarchimedean} when $\norm{x+y}\le\max(\norm x,\norm y)$.
Every normed field of characteristic $p>0$ has this property: Frobenius gives
$\abs{x+y}^{p^m}\le2\max(\abs x,\abs y)^{p^m}$, and taking roots and letting $m\to\infty$ proves it.
For $n\le\infty$, $C^n$ refers to iterated Fr\'echet differentiability. Power-series analyticity means
local expansion in bounded multilinear coefficients with a positive norm-controlled radius.
For manifolds and $C^\omega$ bundles we use the analytic calculus in which the iterated derivatives are
analytic as well, as in Mathlib's $C^\omega$ convention~\cite{docs}. On open sets this agrees with
power-series analyticity at every point, even for incomplete ranges: the derivative series at each point
converges in the completion to the Fr\'echet derivative, which exists in the original operator space because every nearby point is again a point of analyticity. The two notions
differ only for germs at a single point.

\section{From linear constructions to bundles}\label{sec:bundles}

A vector bundle allows a linear construction to vary over a manifold. The construction must act on linear maps, so that it transports changes of coordinates; its regularity must also be strong enough to transport families of those maps. These are separate requirements. Functoriality supplies the compatibility of the new coordinates, while analysis determines whether they define a smooth or analytic bundle.

Throughout, $K$ is a nontrivially normed field and $\mathsf{Vec}_K$ is the category of normed $K$-spaces and bounded linear maps. No completeness assumption is imposed on the spaces. Write $\cL(E,E')$ for its hom spaces, equipped with the operator norm. For $n\in\N\cup\{\infty,\omega\}$, regularity of order $n$ means continuity when $n=0$, iterated Fr\'echet smoothness for finite $n$ and $n=\infty$, and the analytic manifold regularity class for $n=\omega$. When the range is incomplete, this last convention includes the derivative regularity required by that calculus.

\subsection{The bundle construction}

Fix a $C^n$ manifold $M$ modeled on a normed space $P$. The parameter space $P$ describes local coordinates on the base; it need not be related to the typical fibers $E,F$ of the bundles under consideration. Write $\mathsf{VB}^n_K(M)$ for the category of $C^n$ normed vector bundles over $M$, with morphisms over $\id_M$ whose local \emph{operator-valued} expressions are $C^n$. Thus a morphism from a bundle with fiber $E$ to one with fiber $F$ is locally a $C^n$ map
\[
 x\longmapsto T(x)\in\cL(E,F).
\]
This convention records the regularity needed here; merely requiring regularity of the associated map of total spaces would give a different definition. Bounded bilinearity of composition makes the convention independent of the trivializations.

Let $\epsilon=(\epsilon_1,\ldots,\epsilon_r)$ specify a variance in each variable, and put
\[
 \mathbf C^\epsilon=\mathsf C_1\times\cdots\times\mathsf C_r,
 \qquad
 \mathsf C_a=\mathsf{Vec}_K\ \text{or}\ \mathsf{Vec}_K^{\mathrm{op}}.
\]
A functor $\mathcal F\colon\mathbf C^\epsilon\to\mathsf{Vec}_K$ is $C^n$ if all its maps on products of hom spaces are jointly $C^n$. Products carry the maximum norm. Let $\mathbf B^\epsilon_M$ denote the corresponding product of bundle categories, with the same variances. Natural transformations between such functors have bounded linear components, since their codomain is $\mathsf{Vec}_K$.

\begin{theorem}[lifting linear constructions to bundles]\label{thm:bundle-lifting}
For a fixed $C^n$ manifold $M$, fiberwise application defines, up to canonical natural isomorphism, a functor
\[
\begin{aligned}
 \mathcal B_M^n\colon
 \operatorname{Fun}^{C^n}(\mathbf C^\epsilon,\mathsf{Vec}_K)
 &\longrightarrow
 \operatorname{Fun}(\mathbf B^\epsilon_M,\mathsf{VB}^n_K(M)),\\
 \mathcal F&\longmapsto\mathcal F_M,\qquad
 \eta\longmapsto\eta_M.
\end{aligned}
\]
Here the category on the left is the full subcategory of the ordinary functor category on the functors regular on hom spaces. The resulting bundles, bundle morphisms, and natural transformations are given fiberwise by $\mathcal F$, its action on maps, and the components of $\eta$, respectively.

The same conclusion holds if regularity on whole hom spaces is replaced by preservation of $C^n$ families parametrized by open subsets of the model space $P$.
\end{theorem}

\begin{proof}
Choose a common trivializing cover for the input bundles. For the $a$th bundle, let $g^a_{ij}(x)$ convert coordinates in chart $j$ to coordinates in chart $i$. These are $C^n$ operator-valued maps and satisfy
\[
 g^a_{ij}(x)g^a_{j\ell}(x)=g^a_{i\ell}(x),
 \qquad g^a_{ii}(x)=\id.
\]
In a covariant variable use $g^a_{ij}$; in a contravariant variable use $g^a_{ji}$. Denote the resulting tuple of arrows in $\mathbf C^\epsilon$ by $\mathbf g_{ij}$. Functoriality gives
\[
 \theta_{ij}(x)=\mathcal F(\mathbf g_{ij}(x)),\qquad
 \theta_{ij}(x)\theta_{j\ell}(x)=\theta_{i\ell}(x).
\]
The maps $\theta_{ij}$ are $C^n$ by composition. They therefore glue the products with typical fiber $\mathcal F(E_1,\ldots,E_r)$ into a $C^n$ vector bundle: the associated coordinate changes are $C^n$ because operator evaluation is bounded bilinear. Inverse input transitions are already present among the $g^a_{ij}$, so this argument needs no theorem about inversion on an operator algebra.

Apply $\mathcal F$ also to the local operator families of a tuple of bundle morphisms, reversing their directions in the contravariant variables. Their compatibility with transitions follows from functoriality, and their regularity follows from composition. This gives $\mathcal F_M$. The same argument applied to changes of trivialization identifies different choices canonically and preserves identities and composition.

For a natural transformation $\eta\colon\mathcal F\Rightarrow\mathcal G$, its component on the induced bundles is locally the constant operator $\eta_{(E_1,\ldots,E_r)}$. Naturality says precisely that this operator intertwines the two transition systems:
\[
 \eta_{\mathbf E}\,\mathcal F(\mathbf g_{ij}(x))
 =\mathcal G(\mathbf g_{ij}(x))\,\eta_{\mathbf E}.
\]
It consequently defines a $C^n$ bundle morphism. The same identity for arbitrary input morphisms proves naturality of $\eta_M$. Identities and compositions of transformations are preserved fiberwise. No extra regularity of $\eta$ is needed: its local expression is constant.

For the final assertion, write every transition and local morphism family in a base chart. Its domain is an open subset of $P$, and preservation of such families supplies exactly the regularity used above. The argument otherwise remains unchanged.
\end{proof}

The familywise hypothesis is a condition on $\mathcal F$ for \emph{all} families with the specified parameter space: for fixed objects and every $C^n$ map
\[
 \gamma\colon U\longrightarrow
 \prod_{a=1}^r\operatorname{Hom}_{\mathsf C_a}(E_a,E'_a),
 \qquad U\subseteq P\text{ open},
\]
the operator family $\mathcal F\circ\gamma$ must be $C^n$. It is weaker than regularity on whole hom spaces. In particular, preserving finite-coordinate analytic families can suffice for bundles on a finite-coordinate base even when the functor fails to be analytic on an infinite-dimensional hom space. The distinction between these two requirements is central to this paper.

The theorem also applies to restricted categories of fibers, provided the prescribed bundle morphisms take values in their hom spaces. It is the familiar transport of a linear construction to bundles, expressed so that both mixed variance and the required operator regularity are explicit; compare \cite{G2,La}.

\subsection{Alternating maps inside multilinear maps}

For normed spaces $E,F$, let $\Mult^k(E;F)$ be the normed space of bounded $k$-linear maps $E^k\to F$, and let $\Alt^k(E;F)$ be its alternating subspace. Alternating means vanishing whenever two arguments agree; this definition applies unchanged in characteristic two. The inclusion
\[
 j_{E,F}\colon\Alt^k(E;F)\longrightarrow\Mult^k(E;F)
\]
is a linear isometry with closed range. Indeed, alternation is an intersection of closed conditions under evaluation at fixed tuples. Both constructions are contravariant in $E$ and covariant in $F$, and the inclusions form a natural transformation. If $F$ is complete, both spaces are complete: operator-norm Cauchy families converge pointwise in $F$, and their boundedness and defining identities pass to the limit. Thus the alternating bifunctor also restricts to Banach spaces.

For $f\in\cL(E,E')$ and $g\in\cL(F,F')$, their joint action is
\[
 B(f,g)(m)(x_1,\ldots,x_k)
 =g\bigl(m(fx_1,\ldots,fx_k)\bigr).
\]
On alternating maps, write
\[
 A=A^k_{E,E';F}\colon\cL(E,E')\longrightarrow
 \cL(\Alt^k(E';F),\Alt^k(E;F)),\qquad A(f)=B(f,\id_F),
\]
for precomposition. The elementary bound
\[
 \norm{B(f,g)(m)}\le\norm g\,\norm m\,\norm f^k
\]
shows that the actions are bounded operators. Their dependence on $(f,g)$ requires more care.

\begin{proposition}[ambient polynomial action]\label{prop:ambient-polynomial}
The joint action on multilinear maps
\[
 B\colon\cL(E,E')\times\cL(F,F')
 \longrightarrow\cL\bigl(\Mult^k(E';F),\Mult^k(E;F')\bigr)
\]
is the diagonal of a bounded $(k+1)$-linear map, and hence is $C^\omega$. Its restriction to alternating inputs is analytic as a map into
\[
 \cL\bigl(\Alt^k(E';F),\Mult^k(E;F')\bigr),
\]
and all its values belong to the closed subspace of operators taking values in $\Alt^k(E;F')$.
\end{proposition}

\begin{proof}
For $v_i=(f_i,g_i)$ in the product of hom spaces, define
\[
 H(v_0,\ldots,v_k)(m)(x_1,\ldots,x_k)
 =g_0\bigl(m(f_1x_1,\ldots,f_kx_k)\bigr).
\]
This is multilinear in the $v_i$ and has norm at most $1$. Its diagonal is $B$. A diagonal of a bounded multilinear map has finite polynomial derivatives of every order, so it is $C^\omega$ even with incomplete range. Restriction of the input $m$ to alternating maps is bounded linear and preserves analyticity. On the diagonal, two equal $x_i$ give two equal $fx_i$, so the output remains alternating. Closedness of the operator subspace follows by applying the closedness of $j_{E,F'}$ pointwise. For $k=0$, the same formula is simply postcomposition by $g_0$.
\end{proof}

For bundles $\mathcal E,\mathcal H$ with transition families $g_{ij},h_{ij}$, the induced alternating transitions are therefore
\[
 m\longmapsto h_{ij}(x)\circ m\circ
       \bigl(g_{ji}(x),\ldots,g_{ji}(x)\bigr).
\]
Their images in the multilinear operator space are analytic whenever the input families are analytic. The unresolved step is whether their power-series coefficients can be chosen in the alternating operator space itself. The mixed coefficient $H$ above need not take alternating values: different $f_i$ do not preserve equality of arguments. Consequently, the closed inclusion alone does not settle analyticity. This is the descent problem studied in Section~\ref{sec:descent}. Once it is solved for a class of fibers or parameter spaces, Theorem~\ref{thm:bundle-lifting} supplies the corresponding bundle construction without further gluing arguments.

\section{Analytic descent through closed subspaces}\label{sec:descent}

The ambient polynomial of Proposition~\ref{prop:ambient-polynomial} puts every alternating construction into the same situation: a map takes values in a closed subspace, and becomes analytic after inclusion into the ambient space. The obstruction is not convergence of its values. It is the existence of multilinear coefficients in the smaller space. We first make this distinction exact, and then explain why ordinary smoothness has no corresponding obstruction.

Throughout this section, $K$ is nontrivially normed, $P,W,Z$ are normed $K$-spaces, and $j\colon W\to Z$ is a linear isometry with closed range. None of these spaces is assumed complete. We identify $W$ with its image when convenient. If $U\subseteq P$ is open and $f\colon U\to W$, an ambient expansion at $x_0\in U$ means
\begin{equation}\label{eq:ambient-series}
 jf(x_0+h)=\sum_{n\ge0}b_n(h,\ldots,h),\qquad
 \sup_n\norm{b_n}R^n<\infty
\end{equation}
for some $R>0$, with bounded $n$-linear $b_n\colon P^n\to Z$ and convergence to the displayed value for $\norm h<R$. Degree zero denotes a constant. Shrinking $R$ when necessary gives the usual local definition of analyticity. For nonarchimedean normed spaces over a complete nonarchimedean field, this use of multilinear representatives agrees with the homogeneous-polynomial lifting norms of Bertram--Gl\"ockner--Neeb~\cite[\S\S7.14--7.19]{BGN}; mere continuity of an algebraic polynomial is a different condition.

\begin{theorem}[the coefficient descent criterion]\label{thm:coefficient-descent}
Suppose \eqref{eq:ambient-series} holds. Then every diagonal $b_n(h,\ldots,h)$ belongs to $j(W)$. The map $f$ is analytic at $x_0$ if and only if there are bounded $n$-linear maps $q_n\colon P^n\to W$ and $r>0$ such that
\[
 jq_n(h,\ldots,h)=b_n(h,\ldots,h)\quad(n\ge0,\ h\in P),
 \qquad \sup_n\norm{q_n}r^n<\infty.
\]
Only the diagonals must agree; the original mixed values of $b_n$ need not lie in $j(W)$.
\end{theorem}

\begin{proof}
Fix $h\in P$ and restrict the expansion to the scalar line $x_0+th$. Its coefficients satisfy a geometric bound. Such a convergent one-variable series has unique coefficients over any nontrivially normed field: if $N$ is the first nonzero coefficient of a series representing zero, division by $t^N$ bounds its norm by a geometric tail tending to zero as $t\to0$, $t\ne0$. This argument uses only convergence of the given series, not completeness of its target.

Closedness also recovers each coefficient separately. Starting with $b_0=jf(x_0)$, suppose the earlier diagonals lie in $j(W)$. For sufficiently small nonzero $t$,
\[
 t^{-n}\left(jf(x_0+th)-\sum_{m<n}t^m b_m(h,\ldots,h)\right)\in j(W).
\]
The geometric tail estimate makes this expression converge to $b_n(h,\ldots,h)$. Hence that vector lies in the closed subspace.

If $f$ is analytic, apply $j$ to one of its expansions and compare on every scalar line. Uniqueness gives precisely the required diagonal identities. Conversely, shrink $r$ below $R$. The proposed partial sums in $W$ have images converging to $jf(x_0+h)$ term by term. The isometry implies their convergence to the already specified vector $f(x_0+h)$ in $W$, and the bound supplies a positive radius.
\end{proof}

This is a local power-series criterion at a point, not an assertion that \texttt{AnalyticAt} alone supplies the full $C^\omega$ derivative tower used in manifold calculus. When the target is incomplete, an expansion there need not provide expansions throughout a neighborhood; analyticity on an open set requires the criterion at every point. The theorem also explains why derivative jets do not settle the problem in positive characteristic: a nonzero Frobenius term can have all ordinary derivatives zero while remaining a nonzero homogeneous term.

\begin{lemma}[smooth descent]\label{lem:closed-smooth}
For $n\in\mathbb N\cup\{\infty\}$, a map $f\colon U\to W$ is $C^n$ if and only if $jf$ is $C^n$.
\end{lemma}
\begin{proof}
Continuity reflects through an isometry. If $jf$ has Fr\'echet derivative $T$ at $x$, then
\[
 Th=\lim_{t\to0,\ t\ne0}\frac{jf(x+th)-jf(x)}t\in j(W).
\]
Thus $T=jS$ for a unique bounded linear $S\colon P\to W$, with the same norm. The isometry transfers the derivative remainder estimate, so $Df=S$. Postcomposition by $j$ is itself a closed linear isometry
$\cL(P,W)\to\cL(P,Z)$: its image consists exactly of operators whose values lie in $j(W)$, a closed condition under each evaluation. Applying the same argument to derivative maps proves every finite order by induction, and hence $C^\infty$.
\end{proof}

In particular, the full alternating bifunctor is $C^\infty$ on its operator spaces: its ambient action is a bounded polynomial, and the inclusion of alternating maps is closed and isometric. This is the smooth descent argument used in~\cite{PR}. Theorem~\ref{thm:bundle-lifting} therefore constructs alternating bundles and their morphisms at every finite regularity and at $C^\infty$, over arbitrary normed model spaces and fields. Analyticity is the additional coefficient problem.

We call a finite sum of diagonals of bounded multilinear maps a \emph{continuous polynomial}.
For a homogeneous ambient polynomial, the descent problem has a particularly economical form.
\begin{proposition}[one lift, every base point]\label{prop:homogeneous-lift}
Let $k\ge1$, and suppose $a\colon H\to W$ satisfies
$ja(h)=B(h,\ldots,h)$ for a bounded $k$-linear map $B\colon H^k\to Z$. The following are equivalent:
\begin{enumerate}
\item $a$ is analytic at some point of $H$;
\item there is a bounded $k$-linear $Q\colon H^k\to W$ with $Q(h,\ldots,h)=a(h)$;
\item $a$ has a finite power-series expansion at every point, of infinite radius.
\end{enumerate}
Consequently each precomposition action $A^k_{E,E';F}$ is either a continuous polynomial on its whole domain or analytic nowhere.
\end{proposition}
\begin{proof}
An expansion at $h_0$, compared along $h_0+th$ with the polynomial $B(h_0+th,\ldots,h_0+th)$, has degree-$k$ diagonal $B(h,\ldots,h)$. Its degree-$k$ coefficient is the required $Q$. Conversely, expand $Q(h_0+h,\ldots,h_0+h)$ by multilinearity. The coefficient of degree $r$ is the sum over the $r$ occupied slots, with norm at most
$\binom kr\norm Q\norm{h_0}^{k-r}$. This is a finite expansion. The last assertion follows from Proposition~\ref{prop:ambient-polynomial}.
\end{proof}

\section{When alternating bundles are analytic}\label{sec:positive}

There are two ways to solve the coefficient problem. One can retract the ambient space onto the alternating subspace, thereby handling every parameter space at once. Alternatively, one can use coordinates on the parameter space to replace the ambient coefficients. The first mechanism concerns the fibers; the second concerns the base and the families of morphisms being tested.

\subsection{Retractions and lifts in the fibers}

\begin{proposition}[splitting the alternating inclusion]\label{prop:split-lift}
If the inclusion $j_{E,F}\colon\Alt^k(E;F)\hookrightarrow\Mult^k(E;F)$ has a bounded linear retraction $\rho$, then $A^k_{E,E';F}$ has a bounded $k$-linear lift of norm at most $\norm\rho$, for every normed $E'$. Such a lift also exists in each of the following cases:
\begin{enumerate}
\item $k!\ne0$ in $K$;
\item either $E$ or $E'$ has a finite algebraic basis with continuous coordinate functionals.
\end{enumerate}
No completeness hypothesis is required.
\end{proposition}
\begin{proof}
For a retraction, set
\[
 Q(u_1,\ldots,u_k)(m)
 =\rho\bigl(m\circ(u_1,\ldots,u_k)\bigr).
\]
The ambient mixed composition has norm at most $\norm m\prod_i\norm{u_i}$, and its diagonal is already alternating.

If $k!\ne0$, use instead
\[
 Q(u_1,\ldots,u_k)(m)(x)
 =\frac1{k!}\sum_{\sigma\in S_k}
 m(u_{\sigma(1)}x_1,\ldots,u_{\sigma(k)}x_k).
\]
When two $x$-arguments agree, pairing permutations that exchange their positions proves vanishing, in every characteristic. On the diagonal all $k!$ summands coincide. The operator norm is bounded by $k!/\abs{k!}$, where the numerator is the real number counting summands. Equivalently, normalized antisymmetrization retracts $\Mult^k(E;F)$ onto $\Alt^k(E;F)$ with this bound.

For continuous coordinates $(e_i,\varepsilon_i)_{1\le i\le d}$ on $E$, sum over increasing words $s_1<\cdots<s_k$:
\[
 Q(u_1,\ldots,u_k)(m)(x)
 =\sum_s\det\bigl(\varepsilon_{s_a}(x_b)\bigr)_{a,b}
       m(u_1e_{s_1},\ldots,u_ke_{s_k}).
\]
If the coordinates are instead on $E'$, replace the summand by
$\det(\varepsilon_{s_a}(u_ax_b))_{a,b}\,m(e_{s_1},\ldots,e_{s_k})$.
Equal columns give alternation; the exterior-basis expansion gives the diagonal identity. Each is a finite bounded sum, with norm at most
$k!\sum_s\prod_a\norm{\varepsilon_{s_a}}\norm{e_{s_a}}$.
If $k>d$, the relevant alternating space is zero and the sum is empty.
\end{proof}

Over a complete field, every finite-dimensional normed space has continuous coordinates. One can prove this by induction on dimension: a codimension-one kernel is complete and hence closed; its positive distance from a fixed vector outside it bounds the corresponding coordinate functional. The coordinate map and its inverse are then bounded, giving completeness for the induction. Over an incomplete field this condition must be retained: if $a\in\widehat K\setminus K$, the $a$-coordinate on $K+Ka\subseteq\widehat K$ is discontinuous. Finite dimension of $F$ alone does not provide a substitute; scalar-valued forms already exhibit the obstruction of Section~\ref{sec:sharpness}.

A nonarchimedean normed space is \emph{spherically complete} if every nonempty family of pairwise-intersecting
closed balls has a common point. In a nonarchimedean metric, intersecting balls are nested.

\begin{theorem}[spherically complete targets]\label{thm:spherical}
Suppose $K$ is nonarchimedean and $F$ is nonarchimedean and spherically complete. Then, for every normed $E$, the space $\Alt^k(E;F)$ is nonarchimedean and spherically complete, and its inclusion into $\Mult^k(E;F)$ has a contractive linear retraction. Thus $A^k_{E,E';F}$ is a continuous polynomial for every normed $E,E'$ and every $k$. Neither the field nor the source spaces need be complete.
\end{theorem}
\begin{proof}
We recall the extension argument with the simultaneous bounds needed here. Let $V$ be a nonarchimedean seminormed space and $G$ a nonarchimedean spherically complete space. Suppose a nonempty family of bounded operators $T_\alpha\colon V\to G$ and radii $r_\alpha\ge0$ satisfy
$\norm{T_\alpha-T_\beta}\le\max(r_\alpha,r_\beta)$, and a linear map $T_0$ on a subspace $D$ satisfies
$\norm{T_0d-T_\alpha d}\le r_\alpha\norm d$.
To extend over $x\notin D$, choose its value in all balls
\[
 \overline B\bigl(T_\alpha(d+x)-T_0d,\ r_\alpha\norm{d+x}\bigr),
 \qquad \alpha,\ d\in D.
\]
These balls meet pairwise: for $r_\alpha\le r_\beta$, the difference of their centers is bounded by
$\max(r_\alpha\norm{d+x},r_\beta\norm{e+x})$, using the hypotheses and the nonarchimedean inequality. Spherical completeness supplies a common value. Zorn's lemma then extends $T_0$ to all of $V$, preserving every bound. Taking one center zero and radius one gives norm-preserving extension of contractions. This is the nonarchimedean extension mechanism underlying Ingleton's theorem~\cite{In}.

Apply the simultaneous form to the free vector space on symbols $[x_1,\ldots,x_k]$, with seminorm
\[
 \left\lVert\sum_x c_x[x]\right\rVert
 =\max_x\abs{c_x}\prod_i\norm{x_i},
\]
and let $D$ be the span of the multilinearity and repeated-input relations. Alternating maps to $F$ correspond isometrically to bounded linear maps on this free space vanishing on $D$: one direction uses the nonarchimedean inequality in $F$, and the other evaluates on generators. A family of pairwise-meeting balls in $\Alt^k(E;F)$ therefore supplies centers $T_\alpha$ as above. Extending $T_0=0$ from $D$ gives a common point of these balls. This proves spherical completeness without any separation assumption on the continuous dual of $E$.

The spaces $\Mult^k(E;F)$ and $\Alt^k(E;F)$ have nonarchimedean operator norms. Extend the inverse of their isometric inclusion to the entire multilinear space, now using the latter alternating space as the spherically complete target. The resulting contraction is the required retraction. Proposition~\ref{prop:split-lift} finishes the proof. In degree zero the alternating space is simply $F$.
\end{proof}

Quotienting the free-space construction by its relations and then by zero seminorm gives the separated nonarchimedean exterior power; completion gives the Banach-target representation in Section~\ref{sec:duality}. Its maximum seminorm is essential here; the norm on $E$ itself need not satisfy a nonarchimedean inequality. Equivalent changes of norms preserve all analytic conclusions, so an equivalent nonarchimedean spherically complete norm on $F$ is sufficient.

\begin{corollary}[discretely valued fields]\label{cor:discrete-targets}
Suppose $K$ is nonarchimedean with $\abs{K^\times}=r^\Z$ for some $0<r<1$.
Every Banach target $F$ admitting an equivalent nonarchimedean norm makes $A^k_{E,E';F}$
analytic in every degree and for all normed $E,E'$. The field itself need not be complete.
\end{corollary}
\begin{proof}
Round the equivalent nonarchimedean norm upwards to $r^\Z$, putting the value at zero equal to zero.
The new norm lies between the old one and $r^{-1}$ times that norm; monotonicity and compatibility
with multiplication by $r^\Z$ preserve the norm axioms. It is complete and its nonzero distances
lie in $r^\Z$. Such a space is spherically complete. Indeed, a nested family of nonempty closed balls
either has radii tending to zero, when completeness gives a common point, or has a smallest effective
radius in the discrete value set, when all balls of that radius coincide and lie in the other balls.
Pairwise-intersecting balls in a nonarchimedean space are nested. Apply Theorem~\ref{thm:spherical}
and transport the resulting lift back through the equivalent norms.
\end{proof}

\subsection{Coordinates on the parameter space}

\begin{theorem}[finite-coordinate parameters]\label{thm:finite-parameters}
Let $j\colon W\to Z$ be a closed linear isometry. If $P$ is boundedly linearly isomorphic to $K^d$ and $U\subseteq P$ is open, then $f\colon U\to W$ is analytic if and only if $jf$ is analytic. For the maximum norm on $K^d$, ambient coefficients $b_n$ can be replaced by $W$-valued coefficients $q_n$ with
\[
 \norm{q_n}\le d^n\norm{b_n}\quad(n\ge1).
\]
If $Z$ is nonarchimedean, the factor $d^n$ can be replaced by $1$.
\end{theorem}
\begin{proof}
Theorem~\ref{thm:coefficient-descent} puts every homogeneous diagonal in $W$. Expand it in coordinates as
$b_n(h,\ldots,h)=\sum_{|\alpha|=n}c_\alpha h^\alpha$.
Every $c_\alpha$ belongs to $W$: pass to the algebraic quotient $Z/W$ and use polynomial uniqueness over the infinite field $K$. Only the finite-dimensional span of these coefficients is involved, so the argument does not require separation of points by a continuous dual.

For each $\alpha$, choose the nondecreasing coordinate word $i_1(\alpha),\ldots,i_n(\alpha)$ having these multiplicities, and define
\[
 q_n(h_1,\ldots,h_n)
 =\sum_{|\alpha|=n}c_\alpha\prod_{r=1}^n(h_r)_{i_r(\alpha)}.
\]
There are exactly $d^n$ ordered terms before grouping. Thus the triangle inequality gives the stated bound, and two maximum estimates give the nonarchimedean refinement. The diagonals agree. Their exponential bounds preserve a positive radius in Theorem~\ref{thm:coefficient-descent}. The case $d=0$ is immediate.
\end{proof}

Finite-coordinate bases therefore support alternating bundles in every characteristic and degree, regardless of fiber dimension or completeness. Finite algebraic dimension of the base is not enough over an incomplete field: for $K=\F_p(t)$ and $k\ge p$, the end of Appendix~\ref{app:no-largest} gives a base of dimension $p+1$ without continuous coordinates and a nonanalytic alternating atlas. Infinite coordinates require a summability mechanism. There are two distinct versions; their differing hypotheses matter.

\begin{theorem}[summable parameter spaces]\label{thm:summable-parameters}
Let $j\colon W\to Z$ be a closed linear isometry.
\begin{enumerate}
\item Suppose $K$ and $Z$ are nonarchimedean and $W$ is complete. For $U\subseteq c_0(I,K)$ open, analyticity of $jf$ implies analyticity of $f\colon U\to W$. Coefficients can be lifted without increasing their norms.
\item Suppose $Z$ is complete. Let $a\colon H\to W$ have an ambient bounded homogeneous polynomial representation of fixed degree $d$, that is, $ja(h)=B(h,\ldots,h)$ for a bounded $d$-linear $B\colon H^d\to Z$ on a normed space $H$. For every analytic $\gamma\colon U\to H$, with $U\subseteq\ell^1(I,K)$ open, the composite $a\gamma$ is analytic.
\end{enumerate}
Here $c_0$ carries the supremum norm and $\ell^1$ the ordinary sum norm. The field need not be complete. Each analyticity assertion remains valid for a parameter space $P$ that is a bounded linear retract of the indicated space $V$, with bounded linear $i\colon P\to V$, $r\colon V\to P$ and $ri=\id_P$; in~(1) the lifted coefficients then satisfy $\norm{q_n}\le(\norm r\norm i)^n\norm{b_n}$.
\end{theorem}
\begin{proof}
For $c_0$, order $I$ and group the coordinate coefficients of each ambient $b_n$ into nondecreasing words. Finite-coordinate uniqueness puts the grouped coefficients $c_{i_1,\ldots,i_n}$ in $W$, and the nonarchimedean inequality bounds their norms by $\norm{b_n}$. The formula
\[
 q_n(x_1,\ldots,x_n)
 =\sum_{i_1\le\cdots\le i_n}
 c_{i_1,\ldots,i_n}(x_1)_{i_1}\cdots(x_n)_{i_n}
\]
converges: outside finite sets of words some coordinate factor is arbitrarily small, while the others are bounded. Completeness of $W$ supplies the sum, with norm at most $\norm{b_n}\prod_r\norm{x_r}$. Equality of diagonals holds on finite supports and then by density. Apply Theorem~\ref{thm:coefficient-descent}.

For $\ell^1$, the same sorting argument in degree $d$ gives a lift of norm at most $d!\norm{B'}$ for every bounded $d$-linear $B'\colon\ell^1(J,K)^d\to Z$ whose diagonal lies in $W$. Each grouped coefficient has at most $d!$ terms, and the resulting series is absolutely summable, bounded by $d!\norm{B'}\prod_r\norm{x_r}_1$.

To apply this bound to a nonlinear $\gamma$, factor its expansion locally as $Tg$, where $T\colon\ell^1(J,K)\to\widehat H$ is bounded linear. Explicitly, write $\gamma(x_0+x)=\sum_na_n(x,\ldots,x)$ and choose $s\in K$, $0<\abs s$ within its radius. Set $J=\coprod_{n\ge0}I^n$ and
\[
 T(e_{n,i_1,\ldots,i_n})=s^na_n(e_{i_1},\ldots,e_{i_n}),
 \qquad
 g(x)_{n,i_1,\ldots,i_n}=s^{-n}x_{i_1}\cdots x_{i_n}.
\]
The first formula has uniformly bounded values and extends by absolute summation into $\widehat H$.
The degree-$n$ multilinear coefficient of $g$ has entries
$s^{-n}(x_1)_{i_1}\cdots(x_n)_{i_n}$ in its $n$th block and zero elsewhere, so its norm is at most
$\abs s^{-n}$. Its diagonal block has sum norm $(\norm{x}_1/\abs s)^n$, so $g$ has a convergent
multilinear expansion for $\norm{x}_1<\abs s$, without requiring completeness of $K$, and $T(g(x))=\gamma(x_0+x)$ there.
Extend the fixed outer multilinear map to $\widehat H$ using completeness of $Z$; its diagonal still
lies in $W$ by closedness. The sorted lift of its composite with $T$ is a bounded polynomial into $W$.
Composing with $g$ proves analyticity at $x_0$, and the argument applies at each point.
The construction uses one sorting bound in degree $d$,
independent of the degrees in the expansion of $\gamma$.

Finally, if $ri=\id_P$ for bounded maps $i\colon P\to V$, $r\colon V\to P$, apply either result to the family composed with $r$ on $r^{-1}(U)$, and then restrict along $i$.
\end{proof}

Since Theorems~\ref{thm:finite-parameters} and~\ref{thm:summable-parameters} give analyticity at every point of an open set, the $C^\omega$ conclusion follows (see Conventions).

The $\ell^1$ assertion uses one factorial bound for the fixed outer degree. Applying $n!$ separately to every Taylor coefficient of an arbitrary ambient analytic map would not yield an exponential bound, and does not prove general analytic descent. For the alternating joint action the outer degree is $k+1$, or $k$ for precomposition alone.

\subsection{The bundle theorem}

\begin{corollary}\label{cor:bundle-cases}
Let $M$ be an analytic manifold over $K$, with normed model space $P$. Fiberwise alternating maps give an analytic bifunctor on analytic normed vector bundles and their operator-valued analytic morphisms in each of the following settings:
\begin{center}
\begin{tabular}{p{0.26\linewidth}p{0.63\linewidth}}
\hline
Base model & Hypotheses on the typical fibers\\
\hline
Arbitrary $P$ & $k!\ne0$ in $K$; or $K$ is nonarchimedean and the target fibers admit equivalent nonarchimedean spherically complete norms; or the source fibers have finite continuous coordinates.\\[3pt]
Finite continuous coordinates & Arbitrary normed fibers.\\[3pt]
$c_0(I,K)$, or a bounded linear retract & $K$ nonarchimedean; complete nonarchimedean target fibers.\\[3pt]
$\ell^1(I,K)$, or a bounded linear retract & Complete target fibers.\\
\hline
\end{tabular}
\end{center}
The fiber hypotheses apply to both pairs of bundles at the endpoints of a morphism. Each alternative in the first row specifies a separate setting. At every finite smoothness order and at $C^\infty$, no additional hypothesis on $K,P$ or the fibers is needed.
\end{corollary}
\begin{proof}
The fiber results give a bounded $k$-linear lift $Q$ of $A$; then $(u_1,\ldots,u_k,g)\mapsto\bigl(m\mapsto g\circ Q(u_1,\ldots,u_k)(m)\bigr)$ is a bounded $(k+1)$-linear lift of the joint action $B$, which is therefore analytic. The parameter results instead apply directly to their ambient polynomial along transition families and local morphism families. When the final value-space fiber is complete, the ambient multilinear and operator spaces are complete; they are also nonarchimedean when that fiber is. Theorem~\ref{thm:bundle-lifting} supplies the bundle construction in every case. The smooth assertion follows from Lemma~\ref{lem:closed-smooth}.
\end{proof}

For both $\mathbb Q_p$ and $\mathbb C_p$, every finite degree belongs to the first row: these fields have characteristic zero, although their residue characteristic is $p$. The possibly large bound $k!/\abs{k!}_p$ remains finite and causes no convergence problem for a finite polynomial. Spherical completeness of $\mathbb Q_p$, and its failure for $\mathbb C_p$~\cite[\S1]{Sch}, are therefore irrelevant to this conclusion.

The regularity of the input remains part of the statement. A $C^\infty$ base and $C^\infty$ transitions produce a $C^\infty$ alternating bundle; they do not automatically produce an analytic atlas. For example, on the open unit ball of either field the function $f(0)=0$, $f(x)=p^{\nu(x)^2+1}$ for $x\ne0$, where $\nu(x)=\lfloor-\log_p\abs x\rfloor$, is locally constant off zero and satisfies $f(x)=o(\abs x^m)$ for every $m$. Hence it is $C^\infty$ with zero derivatives, but is not analytic at zero. For $k\ge1$, changing a frame of a trivial rank-$k$ bundle by $\operatorname{diag}(1+f,1,\ldots,1)$ gives smooth nonanalytic transitions on scalar top forms. The underlying bundle is still analytically trivial; it is the chosen smooth atlas that fails to be analytic. Section~\ref{sec:sharpness} addresses the stronger obstruction arising from analytic input families.

\section{Exterior powers and scalar detection}\label{sec:duality}

The closed-subspace description explains how analytic coefficients descend. The exterior-power
description explains which targets detect their failure. Both are useful: the first supplies proofs,
while the second distinguishes a universal alternating construction from its scalar observations.

For $k\ge1$ give the algebraic exterior power $\Lambda^k_KE$ the seminorm
\[
 \norm z_\pi=\inf_{z=\sum_jx_{j1}\wedge\cdots\wedge x_{jk}}
                 \sum_j\prod_{a=1}^k\norm{x_{ja}}.
\]
Its quotient by the zero-seminorm subspace is denoted $\Lambda^k_\pi E$. This is an uncompleted normed
space, with universal alternating map $\omega_E(x_1,\ldots,x_k)=[x_1\wedge\cdots\wedge x_k]$.
Put $\Lambda^0_\pi E=K$ and $\omega_E()=1$.

\begin{proposition}[representation]\label{prop:exterior}
There is a natural linear isometry
\[
 \Alt^k(E;F)\cong\cL(\Lambda^k_\pi E,F)
\]
for every normed $F$. Under it, precomposition by $u\colon E\to D$ is precomposition by
$\Lambda^k_\pi u\colon\Lambda^k_\pi E\to\Lambda^k_\pi D$.
\end{proposition}
\begin{proof}
The algebraic linearization $\widetilde m$ of a bounded alternating map satisfies
\[
 \norm{\widetilde m(z)}\le\norm m\sum_j\prod_a\norm{x_{ja}},
 \qquad \norm{\widetilde m(z)}\le\norm m\norm z_\pi.
\]
It therefore kills the null subspace and descends contractively. Conversely,
$\norm{\omega_E(x_1,\ldots,x_k)}\le\prod_a\norm{x_a}$, so composing a bounded linear map with
$\omega_E$ is contractive. The two operations are inverse and hence isometric. Applying $u$ to each
wedge factor proves the intertwining identity and the bound $\norm{\Lambda^k_\pi u}\le\norm u^k$
for $k\ge1$. In degree zero the action is the identity on $K$.
\end{proof}

The quotient removes the vectors invisible to bounded linearizations and makes the representing
object normed. If $F$ is complete, the same correspondence holds with the completion
$\widehat{\Lambda^k_\pi E}$. The uncompleted version is needed for arbitrary normed targets.

Write $Y^*=\cL(Y,K)$ and $J_Y(y)(\lambda)=\lambda(y)$ for the canonical map into the continuous
bidual. The map $J_Y$ is contractive; no point-separation or Hahn--Banach assumption is implicit.

\begin{theorem}[universal and scalar tests]\label{thm:duality}
Let $U$ be an open parameter domain and $W\colon U\to\cL(X,Y)$ a family. Define
\[
 R_F(t)(T)=T\circ W(t),\qquad
 R_F\colon U\to\cL(\cL(Y,F),\cL(X,F)).
\]
For power-series analyticity, or for any fixed class $C^n$, the following statements hold.
\begin{enumerate}
\item $W$ is regular if and only if $R_F$ is regular for every normed $F$; the single target $Y$ suffices.
\item $R_K$ is regular if and only if $J_YW$ is regular in $\cL(X,Y^{**})$, and if and only if
$R_{G^*}$ is regular for every normed $G$. Any one nonzero continuous dual $G^*$ suffices for the last test.
\item If $J_Y$ has a bounded linear left inverse, regularity of $R_K$ implies regularity of $W$.
If $J_F$ has such a left inverse, regularity of $R_K$ implies regularity of $R_F$.
\end{enumerate}
\end{theorem}
\begin{proof}
Precomposition depends bounded-linearly on $W$, and $W(t)=R_Y(t)(\id_Y)$, proving the first assertion.
Exchanging the arguments defines a linear isometry
\[
 \Theta\colon\cL(Y^*,X^*)\longrightarrow\cL(X,Y^{**}),
 \qquad (\Theta(B)x)(\lambda)=(B\lambda)(x).
\]
Both operator norms are the least $C$ bounding the displayed scalar by $C\norm x\norm\lambda$.
The inverse exchanges arguments again, and $\Theta(R_K(t))=J_YW(t)$.
Similarly, the isometry
\[
 \cL(X,G^*)\cong\cL(G,X^*),\qquad T\longmapsto[g\mapsto(x\mapsto T(x)(g))]
\]
identifies $R_{G^*}(t)$ with $S\mapsto R_K(t)\circ S$. This proves regularity for every dual target.
Conversely, choose $v\in G^*$ and $g\in G$ with $v(g)=1$. Injection $\alpha\mapsto v\alpha$ and
evaluation at $g$ compress $R_{G^*}$ to $R_K$ by fixed bounded linear maps.
Finally, if $P_YJ_Y=\id_Y$, then $W=P_Y\Theta(R_K)$. If $P_FJ_F=\id_F$, use
\[
 R_F(t)(T)=P_F\circ R_{F^{**}}(t)(J_F\circ T).
\]
All identities use bounded linear operations, which preserve each stated regularity class.
\end{proof}

Apply the theorem with $X=\Lambda^k_\pi E$, $Y=\Lambda^k_\pi D$ and $W(u)=\Lambda^k_\pi u$. Under the isometry of Proposition~\ref{prop:exterior}, $R_F$ is the precomposition action $A^k_{E,D;F}$, which is $C^\infty$ by Proposition~\ref{prop:ambient-polynomial} and Lemma~\ref{lem:closed-smooth}.
It follows that the exterior-power functor is $C^\infty$, and is analytic on all hom spaces exactly
when $\Alt^k$ is. Scalar alternating maps detect its operator-norm action after composition with $J_Y$;
this is norm analyticity in a new target, not a weak-star regularity condition. In particular every
nonzero continuous-dual target has exactly the same analytic action as the scalar target, even over
an incomplete field. The retraction criterion gives a sufficient route from this bidual information
back to the original target; its necessity is not asserted.

There is a useful way to see the remaining difficulty. For $k\ge1$ and a fixed
$z=\sum_j\omega_E(x_{j1},\ldots,x_{jk})$, the map $u\mapsto(\Lambda^k_\pi u)z$ has the lift
\[
 (u_1,\ldots,u_k)\longmapsto
       \sum_j\omega_D(u_1x_{j1},\ldots,u_kx_{jk}).
\]
For $z\ne0$ its norm can be made at most $2\norm z$ by the definition of the projective norm.
However, the chosen mixed coefficients need not depend linearly on $z$. Operator-valued analyticity
requires precisely such a bounded linear dependence. Thus even uniformly bounded choices of individual
orbit lifts do not finish the construction. Degree zero has constant action.

\paragraph{\textit{Changing the target category.}}
Over a nonarchimedean field replace the sum in the projective seminorm by a maximum, then separate
and complete. Write the result as $\widehat\Lambda^k_{\pi,\mathrm{na}}E$, with
$\widehat\Lambda^0_{\pi,\mathrm{na}}E=\widehat K$. The preceding universal-property proof, using the
nonarchimedean inequality, gives
\[
 \Alt^k(E;F)\cong\cL(\widehat\Lambda^k_{\pi,\mathrm{na}}E,F)
 \quad\text{for complete nonarchimedean }F.
\]
This is the exterior version of Serre's nonarchimedean projective tensor construction~\cite[\S4]{Se}.
The sum and maximum constructions represent different target categories.
Over a discretely valued field, the completed nonarchimedean exterior-power action is analytic in every
degree: use its own target $F=\widehat\Lambda^k_{\pi,\mathrm{na}}D$ in
Corollary~\ref{cor:discrete-targets}, and evaluate the action at the universal alternating map $D^k\to\widehat\Lambda^k_{\pi,\mathrm{na}}D$.
The general-target obstruction of Section~\ref{sec:sharpness} therefore concerns the sum construction;
its target admits no equivalent nonarchimedean norm. This distinction is necessary when interpreting
the result in the usual nonarchimedean Banach category.

\section{Obstructions and their geometric realization}\label{sec:sharpness}

Analytic descent can fail even for the homogeneous polynomial underlying precomposition. The following
two constructions establish sharpness. Their full proofs occupy
Appendices~\ref{app:FF}--\ref{app:scalar}; we first state their conclusions and then show how the
operator obstructions become failures of analytic bundle transitions. The scalar construction uses
van der Put's description of the continuous dual of a bounded sequence space~\cite{VdP}; its
multilinear extension and the counterexample itself are proved in Appendix~\ref{app:scalar}.

\begin{theorem}[operator obstructions]\label{thm:operator-obstructions}
Let $K$ be a nontrivially normed field of characteristic $p>0$, and let $k\ge p$.
\begin{enumerate}
\item There are Banach $K$\nobreakdash-spaces $E,F$ for which
\[
 A^k_{E,E;F}\colon\cL(E,E)\longrightarrow
 \cL\bigl(\Alt^k(E;F),\Alt^k(E;F)\bigr)
\]
is analytic at no point. The target $F$ admits no equivalent nonarchimedean norm.
\item If $K$ is complete and not spherically complete, there are nonarchimedean Banach spaces $E,D$ for which $A^k_{E,D;K}$ is analytic at no point. They can be chosen as closed subspaces of
\[
 \linf(\Lambda,K^{k+1})\quad\text{and}\quad\linf(\Lambda,K^k),
\]
respectively, containing the corresponding $c_0$ spaces, with $\Lambda$ countable.
\end{enumerate}
\end{theorem}

For part~(1), begin over a Laurent series field with a bounded-sequence source and a completed
projective exterior target. The finite-field multiplier theorem of Appendix~\ref{app:FF} rules out
the required homogeneous lift, as shown in Appendix~\ref{app:laurent}. Projective extension of
scalars and passage through the completion transfer the obstruction to every $K$ in the statement
(Appendix~\ref{app:ascent}); Appendix~\ref{app:assembly} assembles the proof and establishes the claim
about the target norm. Over the Laurent field itself, the same construction gives the sharper
$c_0$-base obstruction in
Corollary~\ref{lau:cor:diagonal-bundle}.

For part~(2), the source spaces are cut out by limit conditions along chains in a countable tree.
A finite-dimensional incompatibility and a multilinear tail theorem exclude a bounded lift;
Appendix~\ref{app:scalar} gives the entire construction. In both arguments,
Proposition~\ref{prop:homogeneous-lift} turns the absence of a degree-$k$ lift into nonanalyticity at
every point. Proposition~\ref{prop:scalar-countable} shows that the scalar witnesses must both fail
to be of countable type.

These inputs concern different categories. Part~(1) proves that the factorial hypothesis cannot be removed for arbitrary normed targets, including over spherically complete fields. Part~(2) places the entire obstruction inside nonarchimedean spaces, with the smallest nonzero target. Completeness of the witnesses is useful for sharpness; it imposes no additional restriction on the positive results for general normed fibers.

\begin{corollary}[regularity on whole hom spaces]\label{cor:classification}
For every nontrivially normed field $K$, the bifunctor
$\Alt^k\colon\mathsf{Vec}_K^{\mathrm{op}}\times\mathsf{Vec}_K\to\mathsf{Vec}_K$
is $C^\infty$, and is analytic if and only if $k!\ne0$ in $K$.
Both assertions remain true on the full subcategory of Banach spaces.
\end{corollary}
\begin{proof}
Smooth descent applies to the ambient polynomial action. An invertible factorial gives a finite
multilinear lift, and bounded bilinear postcomposition supplies the joint bifunctor action.
If the factorial vanishes, Theorem~\ref{thm:operator-obstructions}(1) excludes analyticity even on
Banach spaces, after fixing the output morphism to be the identity.
\end{proof}

\begin{theorem}[universal scalar classification]\label{thm:scalar-classification}
Let $K$ be a complete nonarchimedean, nontrivially normed field. For any $k\ge0$, the contravariant functor
\[
 \Alt^k(-;K)\colon\mathsf{Vec}_K^{\mathrm{op}}\longrightarrow\mathsf{Vec}_K
\]
is analytic on every hom space if and only if $k!\ne0$ in $K$ or $K$ is spherically complete. The same equivalence holds when the source category is restricted to Banach spaces, or to nonarchimedean Banach spaces.
\end{theorem}

\begin{proof}
An invertible factorial gives a multilinear lift by normalized antisymmetrization. If $K$ is spherically complete, Theorem~\ref{thm:spherical} applies with target $K$. Both arguments permit arbitrary normed sources. If neither condition holds, then $K$ has characteristic $p>0$ with $k\ge p$, and Theorem~\ref{thm:operator-obstructions}(2) supplies nonarchimedean Banach witnesses against all three assertions.
\end{proof}

For a nonzero finite-dimensional target $F$, choose a basis $v_i$ with bounded coordinate functionals
$\lambda_i$, using completeness of $K$. Then
\[
 A_F(u)(m)=\sum_i v_i A_K(u)(\lambda_i\circ m).
\]
Conversely, choosing $\lambda(v)=1$ compresses $A_F$ to $A_K$ by injection along $v$ and evaluation
by $\lambda$. Thus the same criterion holds for every nonzero finite-dimensional target.
Theorem~\ref{thm:duality} extends this equivalence to every nonzero continuous-dual target.

In particular, the scalar obstruction has no characteristic-zero analogue. It does not arise over either $\mathbb Q_p$ or $\mathbb C_p$, even though the latter is not spherically complete. The relevant characteristic is that of the field, not its residue field.

\subsection{From a bad operator family to changes of frame}

Transition functions take values in invertible operators. The following compression identity shows that this restriction cannot repair a failure on an unrestricted operator space.

\begin{proposition}[shear realization]\label{prop:shear}
Suppose $A=A^k_{E,D;F}$ is not analytic at $u_0\in\cL(E,D)$. Set $H=D\oplus E$, with the maximum norm, and
\[
 g(u)(d,e)=(d+ue,e).
\]
Both $g$ and its inverse family are affine analytic. Nevertheless, the family of pullback operators $g(u)^*$ on $\Alt^k(H;F)$ is not analytic at $u_0$.
\end{proposition}

\begin{proof}
The inverse family is $g(-u)$. Define bounded linear maps
\[
 J\colon\Alt^k(D;F)\longrightarrow\Alt^k(H;F),\qquad
 Jm=m\circ\pi_D^{\times k},
\]
and $R\colon\Alt^k(H;F)\to\Alt^k(E;F)$ by restriction along $i_E(e)=(0,e)$. Since $\pi_Dg(u)i_E=u$,
\begin{equation}\label{eq:shear-compression}
 R\circ g(u)^*\circ J=A(u).
\end{equation}
Bounded linear compression preserves analyticity, so analyticity of $g(u)^*$ at $u_0$ would contradict the hypothesis.
\end{proof}

Take the base $M=\cL(E,D)$ and the trivial bundle $M\times H$. Its product trivialization and the trivialization $(u,z)\mapsto(u,g(u)z)$ form an analytic bundle atlas. The associated alternating transitions are pullback by $g(u)$ and $g(-u)$. Equation~\eqref{eq:shear-compression} proves that this induced atlas fails to be analytic whenever $A$ does. Equivalently, the analytic bundle automorphism $(u,z)\mapsto(u,g(u)z)$ does not induce an analytic morphism on the alternating product bundle.

Thus every failure of the hom-space criterion has a bundle realization. Theorem~\ref{thm:operator-obstructions} supplies Banach bases on which this failure occurs at every point, including the identity operator $g(0)$. This is a statement about the natural construction from an analytic atlas, and about its action on morphisms. The underlying pointwise bundle is trivial and still has an analytic product structure.

\subsection{Tangent bundles and individual forms}

For scalar values the obstruction can be realized by derivatives of analytic coordinate changes, rather than by an independently chosen vector bundle.

\begin{proposition}[tangent and pullback obstructions]\label{prop:tangent-obstruction}
Suppose $K$ is complete, has characteristic $p>0$, is not spherically complete, and $k\ge p$.
\begin{enumerate}
\item There is an analytic Banach manifold with two global analytic charts whose induced transition on scalar alternating $k$\nobreakdash-forms is analytic at no point.
\item There are Banach spaces $N,Y$, a polynomial analytic map $h\colon N\to Y$, and an analytic map $\omega_0\colon Y\to\Alt^k(Y;K)$ such that $h^*\omega_0$, expressed in the product coordinates on $N$, is analytic at no point as an $\Alt^k(N;K)$\nobreakdash-valued map. After inclusion in $\Mult^k(N;K)$ it is analytic everywhere.
\end{enumerate}
\end{proposition}

\begin{proof}
Choose $E,D$ from Theorem~\ref{thm:operator-obstructions}(2), and write $A=A^k_{E,D;K}$. For~(1), put $X=\cL(E,D)\times D\times E$ and use the charts $\id_X$ and
\[
 \psi(u,d,e)=(u,d+ue,e).
\]
Operator evaluation is bounded bilinear, so $\psi$ and its inverse are polynomial analytic. Its derivative is
\[
 D\psi(u,d,e)(v,a,b)=(v,a+ve+ub,b).
\]
Extend forms from $D$ along the projection $\pi_D\colon X\to D$, and restrict forms on $X$ along $i_E(b)=(0,0,b)$. These fixed bounded maps, again denoted $J,R$, give
\[
 R\circ A^k_{X,X;K}(D\psi(u,d,e))\circ J=A(u).
\]
Analyticity of the transition at $(u_0,d_0,e_0)$ would make $A$ analytic at $u_0$ after restriction to the affine slice $(u,d_0,e_0)$. The reverse transition gives the same argument with $-u$.

For~(2), let $B=\Alt^k(D;K)$, $Y=B\times D$, and $N=\cL(E,D)\times B\times E$. Define
\[
 \omega_0(m,d)=m\circ\pi_D^{\times k},\qquad
 h(u,m,e)=(m,ue).
\]
The first coefficient map is bounded linear and the second map is polynomial. Since
\[
 Dh(u,m,e)(v,a,b)=(a,ve+ub),
\]
restriction of the pulled-back form to the fixed $E$\nobreakdash-directions gives $\Gamma(u,m)=A(u)(m)$. If $h^*\omega_0$ were alternating-valued analytic at a point, restriction of its values and of its parameter $e$ would make $\Gamma$ analytic there. Its partial derivative in $m$ is the operator $A(u)$. A power series into a Banach target differentiates on a smaller ball, giving an analytic operator-valued derivative there. Fixing $m$ would therefore make $A$ analytic, a contradiction. In the ambient multilinear space, the pullback formula is a bounded multilinear expression in $Dh$ and $\omega_0\circ h$, and hence is analytic.
\end{proof}

The second assertion distinguishes a failure of intrinsic analyticity for individual forms from a failure of operator-valued bundle transitions. It also suggests the appropriate calculus when analytic descent is unavailable.

The chart in Proposition~\ref{prop:tangent-obstruction}(1) realizes any nowhere analytic action $A^k_{E,D;F}$. With Theorem~\ref{thm:operator-obstructions}(1) it works over every field of characteristic $p$, and it extends to tangent-valued forms.

\begin{proposition}[tangent obstructions with other values]\label{prop:tangent-values}
Let $K$ be a nontrivially normed field of characteristic $p>0$, complete or not, and let $k\ge p$.
\begin{enumerate}
\item There are a Banach space $F$ and an analytic Banach manifold with two global analytic charts whose induced transition on $F$\nobreakdash-valued alternating $k$\nobreakdash-forms is analytic at no point.
\item There is an analytic Banach manifold with two global analytic charts whose induced transition on tangent-valued alternating $k$\nobreakdash-forms is analytic at no point.
\item If $K$ is complete and not spherically complete, a manifold as in~(2) exists with a nonarchimedean Banach model.
\end{enumerate}
The target $F$ constructed for~(1) admits no equivalent nonarchimedean norm; for $K=\F_q((u))$ as in Corollary~\ref{lau:cor:dichotomy}, the model space can be chosen nonarchimedean. The model space constructed for~(2) admits no equivalent nonarchimedean norm.
\end{proposition}

\begin{proof}
For~(1), take $E,F$ from Theorem~\ref{thm:operator-obstructions}(1). Repeat the proof of Proposition~\ref{prop:tangent-obstruction}(1) with $D=E$ and values in $F$. On $X=\cL(E,E)\times E\times E$ the same maps $J,R$ give
\[
 R\circ A^k_{X,X;F}(D\psi(w,d,e))\circ J=A^k_{E,E;F}(w),
\]
which is analytic at no point. The argument uses only bounded linear compression and affine slices, so completeness of $K$ is not needed. For $K=\F_q((u))$ with its $u$\nobreakdash-adic absolute value, Remark~\ref{rem:spaces} allows $E=\linf(\N,K)$ and $F=B$, the completed projective exterior power of Theorem~\ref{thm:laurent}, and $B$ has no equivalent nonarchimedean norm by Proposition~\ref{prop:nonultra}. Then $X$ is nonarchimedean, since operator norms into a nonarchimedean space and maximum norms of products are nonarchimedean.

For~(2) and~(3), suppose that $A^k_{E,D;F}$ is analytic at no point. Put $D'=D\times F$ and $X=\cL(E,D')\times D'\times E$, with the charts $\id_X$ and $\psi(u,d,e)=(u,d+ue,e)$ as before. Write $g=D\psi(u,d,e)$, so that $g^{-1}(v,a,b)=(v,a-ve-ub,b)$. As functions of the point $(u,d,e)$, the induced transitions on tangent-valued forms are $m\mapsto g\circ m\circ(g^{-1},\ldots,g^{-1})$ and $m\mapsto g^{-1}\circ m\circ(g,\ldots,g)$. Let $\pi_D\colon X\to D$ and $\pi_F\colon X\to F$ be the components of the middle factor, and put $i_F(y)=(0,(0,y),0)$ and $i_E(b)=(0,0,b)$. The bounded linear map
\[
 \Phi(T)(m)=\pi_F\circ T\bigl(i_F\circ m\circ\pi_D^{\times k}\bigr)\circ i_E^{\times k}
\]
sends operators on $\Alt^k(X;X)$ to operators $\Alt^k(D;F)\to\Alt^k(E;F)$. For $f,h\in\cL(X,X)$ it sends $m\mapsto h\circ m\circ(f,\ldots,f)$ to $m\mapsto(\pi_Fhi_F)\circ m\circ(\pi_Dfi_E,\ldots,\pi_Dfi_E)$. Here $\pi_Fg^{\pm1}i_F=\id_F$ and $\pi_Dg^{\pm1}i_E=\pm u_D$, where $u_D\in\cL(E,D)$ is the $D$\nobreakdash-component of $u$. So $\Phi$ sends the two transitions to $A^k_{E,D;F}(\mp u_D)=(\mp1)^kA^k_{E,D;F}(u_D)$. If either transition were analytic at $(u_0,d_0,e_0)$, restriction to the affine slice $w\mapsto(u_0+(w,0),d_0,e_0)$, $w\in\cL(E,D)$, would make $A^k_{E,D;F}$ analytic at the $D$\nobreakdash-component of $u_0$.

For~(2), take $E,F$ from Theorem~\ref{thm:operator-obstructions}(1) and $D=E$. The model $X$ contains the isometric copy $i_F(F)$ of $F$, so it admits no equivalent nonarchimedean norm. For~(3), take $E,D$ from Theorem~\ref{thm:operator-obstructions}(2) and $F=K$. Then $X$ is a nonarchimedean Banach space, as in~(1).
\end{proof}

\paragraph{\textit{Tangent bundles.}}
For $\Alt^k(TM;F)$ the source fiber is the model space $P$. Finite coordinates on the base are therefore finite coordinates on the source fiber, and Proposition~\ref{prop:split-lift}(2) already makes the action analytic on whole hom spaces. This is the last case of the first row of Corollary~\ref{cor:bundle-cases}. Over a finite-coordinate base, Theorem~\ref{thm:finite-parameters} and the second row add cases only when the source fibers lack finite coordinates. Examples are source bundles of infinite rank, and pulled-back tangent bundles $f^*TN$ along analytic maps $f\colon M\to N$ from a finite-coordinate manifold $M$ into an infinite-dimensional manifold $N$.

Over a discretely valued field, Corollary~\ref{cor:discrete-targets} makes $\Alt^k(TM;F)$ analytic on every analytic manifold for every Banach $F$ with an equivalent nonarchimedean norm, as noted after Corollary~\ref{lau:cor:dichotomy}. It also makes tangent-valued forms analytic when the model is such a space. So over $\F_q((u))$ the target in Proposition~\ref{prop:tangent-values}(1) cannot have an equivalent nonarchimedean norm, while the manifold can be modeled on a nonarchimedean Banach space. This is consistent with~(3): a complete discretely valued field is spherically complete (Lemma~\ref{lem:discdist}).

\section{Analytic differential forms}\label{sec:forms}

An alternating form can be analytic as a multilinear-valued function even when it has no analytic expansion with alternating-valued coefficients (Proposition~\ref{prop:tangent-obstruction}(2)). The ambient viewpoint retains exactly the regularity needed for the usual operations on differential forms.

Throughout this section $K$ is complete and nontrivially normed. Manifold models may be arbitrary normed spaces. We use $C^\omega$ analytic charts and maps, so their derivatives are analytic also for incomplete models. Coefficient spaces of scalar multilinear forms are Banach, where this agrees with power-series analyticity. No characteristic or degree restriction is imposed.

\begin{definition}
An \emph{ambient analytic scalar $k$\nobreakdash-form} on $M$ is a family $\omega_x\in\Alt^k(T_xM;K)$ whose representative in each chart is analytic after the inclusion
\[
 j\colon\Alt^k(P;K)\hookrightarrow\Mult^k(P;K).
\]
Write $\Omega^k_{\mathrm{amb}}(M)$ for the space of such forms. In degree zero this is the space of analytic scalar functions.
\end{definition}

This is a bundle definition: $\omega$ is an analytic section of the multilinear bundle with alternating values. The multilinear bundle is always analytic, since its operator-valued transitions have bounded multilinear lifts. Concretely, under a coordinate change $\psi$ the representatives satisfy
\[
 j\omega_\phi(y)=
 (j\omega_{\phi'}(\psi(y)))\circ(D\psi(y),\ldots,D\psi(y)).
\]
The right-hand side is analytic by bounded multilinear composition. It therefore suffices to check the definition on any analytic atlas, and adjoining further analytic charts does not change the space of forms.

\begin{theorem}[ambient differential calculus]\label{thm:ambient-forms}
Ambient analytic forms define a sheaf of graded algebras on every analytic manifold. Analytic maps induce pullbacks, and there is an exterior derivative
\[
 d\colon\Omega^k_{\mathrm{amb}}(M)\longrightarrow\Omega^{k+1}_{\mathrm{amb}}(M).
\]
The product is associative and graded commutative, $d^2=0$, pullback commutes with both product and exterior derivative, and
\[
 d(\eta\wedge\zeta)=d\eta\wedge\zeta+(-1)^k\eta\wedge d\zeta
 \qquad(\eta\in\Omega^k_{\mathrm{amb}}(M)).
\]
Thus $\Omega^\bullet_{\mathrm{amb}}$ is a contravariant differential graded algebra construction in every characteristic.
\end{theorem}

\begin{proof}
Locality of analytic coefficient maps gives the sheaf property. Addition and multiplication by an analytic function are bounded multilinear operations on the coefficient spaces. For an analytic map $h$, the formula
\[
 h^*\eta=(\eta\circ h)\circ(Dh,\ldots,Dh)
\]
is analytic in the ambient space and has alternating values. The chain rule gives identity and composition laws for pullback.

To define the product, for $S\subseteq\{1,\ldots,k+l\}$ let $v_S$ list its entries in increasing order, and set
\[
 \epsilon(S)=(-1)^{\#\{(s,t):s\in S,\ t\notin S,\ s>t\}}.
\]
Use the shuffle formula
\[
 (\mu\wedge\nu)(v_1,\ldots,v_{k+l})
 =\sum_{\abs{S}=k}\epsilon(S)\mu(v_S)\nu(v_{S^c}).
\]
It defines a bounded bilinear operation on multilinear coefficient spaces, of norm at most $\binom{k+l}{k}$. For alternating inputs it is alternating: terms with two repeated arguments in one block vanish, and the remaining terms cancel in pairs by exchanging those arguments between blocks. This argument includes characteristic two. The formula commutes with simultaneous precomposition and hence glues in charts. Complementing $S$ proves graded commutativity. Expanding either bracketing of a triple product yields the same signed sum over ordered three-block partitions, proving associativity.

On an open subset of a model space $P$, write $\eta$ for an ambient coefficient map. Define
\begin{equation}\label{eq:exterior-derivative}
 (d\eta)_y(v_0,\ldots,v_k)
 =\sum_{i=0}^k(-1)^iD\eta(y)(v_i)
       (v_0,\ldots,\widehat v_i,\ldots,v_k).
\end{equation}
The operation applied to $D\eta$ is bounded linear with norm at most $k+1$, so the result is ambient analytic. Each $D\eta(y)(v_i)$ is alternating: difference quotients have alternating values, and that subspace is closed. If $v_a=v_b$, the summands other than $i=a,b$ vanish; alternation moves the remaining repeated argument across $b-a-1$ positions, making those two terms cancel. Hence $d\eta$ is a form.

The second derivative of an analytic map is symmetric in every characteristic. At the center of a power series, its quadratic contribution is $p_2(v,w)+p_2(w,v)$. In the expansion of $d(d\eta)$, pair the terms differentiated in directions $(v_a,v_b)$ and $(v_b,v_a)$ for $a<b$. Their values agree and their signs are $(-1)^{a+b-1}$ and $(-1)^{a+b}$, so $d^2\eta=0$.

Differentiate the pullback formula. Terms differentiating $\eta$ give $h^*(d\eta)$ after the signed sum~\eqref{eq:exterior-derivative}. The other terms insert $D^2h(v_a,v_b)$ into one argument of $\eta$. Symmetry of $D^2h$ and alternation of $\eta$ cancel the terms in pairs. Therefore $d(h^*\eta)=h^*(d\eta)$; applying this identity to chart changes proves that the local derivatives glue. Finally, differentiating the shuffle formula gives the Leibniz rule: placing the derivative index in the second block moves it past the first $k$ arguments and contributes $(-1)^k$. Every identity uses finite sums and cancellation, with no division by a factorial.
\end{proof}

The same proof gives forms with values in a commutative Banach algebra; for a general bounded bilinear product $\beta$ all properties except associativity and graded commutativity carry over. For such a product the shuffle bound is $\binom{k+l}{k}\norm{\beta}$. The scalar case already supplies the differential graded algebra relevant to the cotangent construction.

\begin{proposition}[comparison with alternating-bundle sections]\label{prop:ambient-agreement}
Every intrinsically analytic alternating-valued coefficient map is ambient analytic. If, for the model space $P$, the inclusion $j$ has a bounded linear retraction, then the natural alternating cotangent atlas is analytic and its analytic sections are precisely the ambient analytic forms. The same conclusion holds when $P$ has finite continuous coordinates.
\end{proposition}

\begin{proof}
The first statement follows by composition with $j$. For a bounded retraction $R$, the identity $\eta=R(j\eta)$ gives the converse. Moreover,
\[
 (u_1,\ldots,u_k)\longmapsto
 \bigl[m\longmapsto R\bigl(m\circ(u_1,\ldots,u_k)\bigr)\bigr]
\]
is a bounded multilinear lift of the precomposition action. Thus the induced alternating cotangent transitions are analytic. For finite continuous coordinates, analytic descent through closed subspaces applies to the coefficient map itself, and Corollary~\ref{cor:bundle-cases} supplies the analytic bundle atlas.
\end{proof}

The retraction hypothesis holds for invertible $k!$, by normalized antisymmetrization, and for spherically complete nonarchimedean $K$, by Theorem~\ref{thm:spherical}. These are sufficient conditions, not a necessity claim. An arbitrarily chosen analytic structure on a pointwise alternating bundle is weaker: the product bundles in Proposition~\ref{prop:tangent-obstruction}(2) have such structures, yet pullback fails to preserve their intrinsically analytic sections. The ambient definition avoids that dependence on an unjustified analytic descent, while preserving the geometric operations that motivate differential forms.

\section{A proposed Mathlib interface}\label{sec:mathlib-interface}

The preceding results suggest a single interface for constructing alternating bundles at every
regularity. Gou\"ezel's proposed construction~\cite{G2} uses characteristic zero
to prove regularity of the induced transition operators. Its assembly through
\texttt{VectorPrebundle.\allowbreak IsContMDiff} can instead use the family-preservation premise below.
The general constructor and its finite-coordinate application are proved in~\cite{LeanP}. The proof class and its instance are a proposal: they are not Mathlib
declarations and are not in~\cite{LeanP}.

\subsection{The family-preservation premise}

Fix a base $M$, its model with corners $I$, and a regularity $n\in\N\cup\{\infty,\omega\}$.
Let $E,E',F,F'$ be normed model fibers and put
\[
 \mathcal O=\cL(E',E)\times\cL(F,F'),\qquad
 \mathcal T=\cL\bigl(\Alt^k(E;F),\Alt^k(E';F')\bigr).
\]
Write $B\colon\mathcal O\to\mathcal T$ for the joint alternating action
\[
 B(u,v)(a)=v\circ a\circ(u,\ldots,u).
\]
The proposed premise is
\begin{equation}\label{eq:mathlib-family-premise}
 \begin{gathered}
 U\subseteq M\text{ open},\quad \gamma\colon M\to\mathcal O,\quad
 \gamma|_U\text{ is }C^n\text{ relative to }I\quad\Longrightarrow\quad
 (B\circ\gamma)|_U\text{ is }C^n.
 \end{gathered}
\end{equation}
The conclusion is regularity in the operator space $\mathcal T$.
Both the base $M$ and its model $I$ are parameters of the premise: $I$ alone does not determine $M$.
The four model fibers allow the same interface to handle morphisms between different alternating
bundles. For the bundle with source model $F_1$ and value model $F_2$, the required specialization is
$E=E'=F_1$ and $F=F'=F_2$.

The following complete definition uses the joint action \texttt{alternatingMapAction} from our
development. All parameters and typeclass hypotheses are explicit; the local names \texttt{Op}
and \texttt{Out} abbreviate $\mathcal O$ and $\mathcal T$. The definition typechecks against~\cite{LeanP}. Stating this predicate requires a charted base, but no completeness or
manifold-regularity assumption.

\Needspace{30\baselineskip}
\begin{leancode}
import AlternatingAnalytic.Geometry.FiniteCoordinateManifoldFamilies

open scoped Manifold ContDiff
open AlternatingAnalytic

def PreservesAlternatingFamilies
    {K P H : Type*} [NontriviallyNormedField K]
    [NormedAddCommGroup P] [NormedSpace K P]
    [TopologicalSpace H] (I : ModelWithCorners K P H)
    (n : ℕ∞ω) (M : Type*)
    [TopologicalSpace M] [ChartedSpace H M]
    (k : ℕ) (E E' F F' : Type*)
    [NormedAddCommGroup E] [NormedSpace K E]
    [NormedAddCommGroup E'] [NormedSpace K E']
    [NormedAddCommGroup F] [NormedSpace K F]
    [NormedAddCommGroup F'] [NormedSpace K F'] : Prop :=
  let Op := (E' →L[K] E) × (F →L[K] F')
  let Out := (E [⋀^Fin k]→L[K] F) →L[K]
    (E' [⋀^Fin k]→L[K] F')
  ∀ {U : Set M} {γ : M → Op}, IsOpen U →
    ContMDiffOn I 𝓘(K, Op) n γ U →
    ContMDiffOn I 𝓘(K, Out) n
      (fun x ↦ alternatingMapAction k (γ x)) U
\end{leancode}

With the same ambient hypotheses, the premise can be packaged as a proof class. The class and
bundle-interface sketches below suppress the repeated ambient binders;
\texttt{Alt k X Y} abbreviates the continuous alternating-map type indexed by \texttt{Fin k}.
These sketches specify the proposed organization rather than standalone Lean programs.

\begin{leancode}
class AlternatingFamilyRegularity
    (I) (n) (M) (k) (E E' F F') : Prop where
  preserves : PreservesAlternatingFamilies
    I n M k E E' F F'
\end{leancode}

This is a proof class. It contains no choice of norm, retraction, coordinates, bundle topology, or
multilinear lift. Such data may be used to prove its field, but the interface records only
\eqref{eq:mathlib-family-premise}. In particular, it does not require $B$ to be $C^n$ on its whole
operator domain: that stronger condition would fail to capture the finite-base analytic theorem
in its full generality.

\Needspace{31\baselineskip}
\subsection{The constructor and the single instance}

The theorem
\texttt{contMDiffVectorBundle\_\allowbreak alternating\_\allowbreak of\_\allowbreak family}
takes the family-preservation proof as an ordinary argument. It does not perform typeclass search
for a reason that the premise holds. Given input bundle transitions, it pairs the reversed source
transition with the forward target transition, applies the premise on their open overlap, and
identifies the resulting operator with the induced alternating transition. This proves
\texttt{VectorPrebundle.\allowbreak IsContMDiff}; Mathlib's theorem
\texttt{VectorPrebundle.\allowbreak contMDiffVectorBundle} then supplies the desired regularity on the existing
bundle.

With the same abbreviations, the proposed public assembly has the following form:
\begin{leancode}
theorem contMDiffVectorBundle_alternating_of_family
    (k : ℕ)
    (hfamily : PreservesAlternatingFamilies
      I n M k F₁ F₁ F₂ F₂)
    [ContMDiffVectorBundle n F₁ E₁ I]
    [ContMDiffVectorBundle n F₂ E₂ I] :
    ContMDiffVectorBundle n (Alt k F₁ F₂)
      (fun x ↦ Alt k (E₁ x) (E₂ x)) I := ...

instance ContMDiffVectorBundle.continuousAlternatingMap
    [ContMDiffVectorBundle n F₁ E₁ I]
    [ContMDiffVectorBundle n F₂ E₂ I]
    [h : AlternatingFamilyRegularity
      I n M k F₁ F₁ F₂ F₂] :
    ContMDiffVectorBundle n (Alt k F₁ F₂)
      (fun x ↦ Alt k (E₁ x) (E₂ x)) I :=
  contMDiffVectorBundle_alternating_of_family k h.preserves
\end{leancode}

Here $E_1,E_2$ are the actual fiber families over $M$, whereas $F_1,F_2$ are their normed model
fibers. The regularity premise mentions only the latter and the base. The constructor reuses the
canonical alternating \texttt{FiberBundle} and \texttt{VectorBundle} structures; it does not introduce
another topology or another atlas construction. Its preliminary prebundle proof can remain a named
theorem or a local instance inside the proof. Only the final assembly above needs to be registered
as the alternation-specific bundle instance.

An application with a specialized family argument may call the constructor directly. Alternatively,
it may install a local proof of \texttt{AlternatingFamilyRegularity} and let the single instance
perform the assembly. Neither usage requires a new global bundle instance.

\subsection{Supplying the regularity proofs}

The sufficient conditions proved in this paper belong at the family layer. They supply the
following proof constructors, with commonly used ones eligible for automatic instances:
\begin{enumerate}
\item For $n\le\infty$, Lemma~\ref{lem:closed-smooth} and the ambient polynomial action prove
preservation over every nontrivially normed field, for arbitrary normed model fibers and every
model with corners.
\item If $k!\ne0$ in $K$, Proposition~\ref{prop:split-lift} proves preservation at every regularity,
including $\omega$. Characteristic zero supplies this condition in every degree. Part~(2) of the same proposition proves the same when a source model fiber $E$ or $E'$ has finite continuous coordinates; this covers $\Alt^k(TM;F)$ over a finite-coordinate base.
\item For ordinary open-chart models $I=\mathcal I(K,P)$, continuous coordinates
$P\simeq_L K^d$ give analytic family preservation with arbitrary normed fibers
(Theorem~\ref{thm:finite-parameters}). The manifold application assumes the corresponding analytic
manifold structure. A convenience theorem can obtain the coordinates from completeness of $K$
and finite dimensionality of $P$; the primitive coordinate theorem needs neither completeness of
the field nor completeness of the fibers.
\item Over a nonarchimedean field, a nonarchimedean spherically complete final target $F'$ gives
preservation at every regularity by Theorem~\ref{thm:spherical}. The summable-parameter results of
Theorem~\ref{thm:summable-parameters} give further model-specific proof constructors under their
stated target hypotheses.
\end{enumerate}
No blanket Banach assumption belongs in the interface. Nor should the finite-coordinate theorem
be silently extended to arbitrary models with corners or arbitrary-set
\texttt{ContMDiffWithinAt} statements: the argument supplied here uses open parameter domains.

The distinction between a property of an action and a property of its input families also matters
for regularity reduction. Preservation of analytic families does not, merely by lowering the
regularity of the output, prove preservation of arbitrary $C^n$ families: the latter need not be
analytic. Thus one should not copy the regularity-lowering instances for
\texttt{ContMDiffVectorBundle} into the new family class without a separate theorem.
The smooth family proof is available independently from closed-range reflection. By contrast,
once a particular bundle has been proved analytic, its lower-order bundle regularity follows
through Mathlib's existing instances.

\subsection{Coherence, morphisms, and verification scope}\label{sec:coherence}

Mathlib's \texttt{ContMDiffVectorBundle} is itself a \texttt{Prop} class: it asserts regularity of
the transition operators and carries no new bundle data. Proof irrelevance identifies different
proofs of the same proposition~\cite{LeanRef}. Consequently overlapping proofs of the proposed
family class cannot create incompatible bundle topologies. They can still introduce alternative
typeclass-search paths. The design therefore gives one alternation-specific assembly instance,
not a claim that instance search has a unique path: ordinary regularity-lowering instances and
the general order-zero instance already provide other paths.

The instance graph should remain directed from analytic hypotheses to family regularity and then
to bundle regularity. Its family proofs should not ask for the output alternating-bundle instance.
Converse results should remain theorems, rather than reverse search rules. Specialized cases can
remain named proof constructors when automatic search would add unnecessary branching; priorities
and elaboration behavior of the common cases must be checked when implementing the proposal.

The four-endpoint premise also supports the morphism construction of
Section~\ref{sec:bundles}. Given regular input morphisms and regular alternating bundles at both
endpoints, apply it to their local operator coordinates to obtain a regular morphism between the
alternating bundles. The self-endpoint premise used for a bundle instance alone does not establish
this cross-endpoint condition. The development~\cite{LeanP} implements this distinction in
\texttt{alternatingBundleHom\_\allowbreak of\_\allowbreak family}.

For an individual chosen atlas the family premise is sufficient and may be stronger than
necessary; only the transition families occurring in that atlas must be regular. Universally
over all normed fibers and all compatible input frames, its necessity is detected by the shear
construction of Proposition~\ref{prop:shear}. An arbitrary family $u(x)\colon E\to D$ factors as
\[
 u(x)=\pi_D S_x i_E,\qquad
 S_x(d,e)=(d+u(x)e,e),\qquad S_x^{-1}(d,e)=(d-u(x)e,e).
\]
The two frames of the trivial $(D\oplus E)$-bundle defined by the identity and $S_x$ test the induced
action on $S_x$; fixed bounded compression then recovers the action on $u(x)$. Together with bounded
bilinear postcomposition this recovers the joint alternating family. This converse requires the
universal quantifiers and the enlarged source fiber, and is not a converse instance for one fixed
bundle or one fixed pair of models.

For the ordinary model $I=\mathcal I(K,P)$ and a nonempty $C^n$ manifold $M$, whether
\eqref{eq:mathlib-family-premise} holds depends only on $P$, $n$, $k$ and the model fibers, not on $M$.
Regularity on $M$ is checked in charts. Conversely, a family on an open subset of $P$ near $x_0$ can be
composed with $y\mapsto x_0+\lambda(y-y_0)$, for $\lambda\in K^\times$ small, and with a chart of $M$
whose image contains a ball around $y_0$; the premise on $M$ then gives regularity near $x_0$. With the
shear converse, if the alternating construction preserves $C^n$ bundles with all normed fibers over one
nonempty manifold modeled on $P$, it does so over every such manifold.

The theorem kernel in~\cite{LeanP} consists of the first three declarations below; the fourth is
the supplied-coordinate application.
\begin{leancode}
contMDiffOn_continuousAlternatingMapCoordChange_of_family
alternatingVectorPrebundle_isContMDiff_of_family
contMDiffVectorBundle_alternating_of_family
contMDiffVectorBundle_alternating_of_finiteCoordinates
\end{leancode}
These declarations use \texttt{Fin k}. An upstream presentation should retain
Mathlib's arbitrary finite index type, transporting the coordinate-based proofs along a bijection
with \texttt{Fin k}, where $k$ is its cardinality. The new proof class, its automatic instances, this
index-type adaptation, and the proposed search organization have not been implemented or checked
as an upstream patch. They are separate from the verification claims in
Section~\ref{sec:verification}.

\section{Further questions and verification}\label{sec:verification}

The descent formulation leaves three closely related questions in the obstructed range
$\operatorname{char}K=p>0$, $k\ge p$.

\paragraph{\textit{Which fibers remove the obstruction?}}
Call $F$ a universal target if $\Alt^k(-;F)$ is analytic on every hom space. A nonarchimedean,
spherically complete norm is sufficient, as is any equivalent norm of that kind. Over a complete
field that is not spherically complete, however, every universal target must have zero continuous dual.
Indeed, if $\lambda(v)=1$ for $\lambda\in F^*$ and $v\in F$, then
\[
 A^k_{E,D;K}(u)(a)=\lambda\circ A^k_{E,D;F}(u)(v a).
\]
The scalar counterexample excludes analyticity of the right-hand action. An intrinsic characterization
of universal targets remains open; zero dual alone has not been shown sufficient.
Remark~\ref{rem:nonspherical-analytic-target} shows separately that spherical completeness of the
given norm is not necessary: an equivalent spherically complete norm can suffice. Over a densely valued,
spherically complete field, it is also open whether every nonarchimedean Banach target is universal.
The general-target obstruction has no equivalent nonarchimedean norm and does not decide this question.

\paragraph{\textit{Which bases permit descent?}}
The finite-coordinate, $c_0$ and $\ell^1$ arguments impose different conditions on the final target.
The remaining problem is to characterize model spaces $P$ for which every analytic family of operators
parametrized by an open subset of $P$ is preserved by $\Alt^k$, separately for arbitrary normed,
complete, and complete nonarchimedean targets. Over $\F_q((u))$ with a $u$\nobreakdash-adic absolute
value, Corollary~\ref{lau:cor:dichotomy} answers this for arbitrary normed targets and Banach models
with an equivalent nonarchimedean norm: exactly the finite-dimensional models qualify.
Over a complete field, universal Banach-valued analytic reflection through closed linear inclusions is
sufficient for complete targets. Over a complete nonarchimedean field it is stronger than necessary for
this fixed-degree action, by Proposition~\ref{fam:prop:l1-optimal} and
Theorem~\ref{thm:summable-parameters}(2). Over $\mathbb R$ and $\mathbb C$ every Banach space has it, by symmetrization (Appendix~\ref{app:reflection}). For Banach parameter spaces,
Theorem~\ref{fam:thm:tensor-analytic} characterizes this reflection property by exponentially bounded
projections onto the closed spans of pure powers in projective tensor powers. The proof and the
sharp factorial projection bounds of Proposition~\ref{fam:prop:l1-optimal} appear in
Appendix~\ref{app:reflection}; they explain why the fixed-degree $\ell^1$ argument goes beyond
universal analytic reflection.

\paragraph{\textit{When do the two notions of analytic form coincide?}}
For a fixed model $P$, a bounded retraction onto alternating forms gives both an analytic cotangent
atlas and equality of ambient and intrinsic analytic forms. It is not known whether compatibility with
every analytic coordinate change on $P$ already forces that equality. A chosen analytic product
structure is insufficient, by Proposition~\ref{prop:tangent-obstruction}. The question concerns the
natural atlas and its full class of coordinate changes.

These questions concern fibers, bases and sections within the same descent problem. Restricting the
objects to a largest full analytic domain is a different extremal problem.
Theorem~\ref{dom:no-largest}, proved in Appendix~\ref{app:no-largest}, shows that such a largest domain
need not exist over the incomplete field $\F_p(t)$: two objects may each have an analytic endomorphism
action while their cross-action is nonanalytic. The construction uses finite algebraic dimension
without continuous coordinates. Whether an analogous failure occurs over complete fields or among
Banach pairs remains open.

\subsection*{Formal verification}
The Lean proofs are in~\cite{LeanP}. Its file
\texttt{challenge.md} is a ledger of 76 claims of this paper. Each claim has a challenge file, stating it in Lean with
\texttt{sorry} in place of the proof, and a solution file proving the identical statement. The command
\texttt{python3 scripts/check\_claim.py ID}, with \texttt{ID} such as \texttt{Thm4\_4}, builds both
files, checks that the proof uses only the axioms \texttt{propext}, \texttt{Classical.choice} and
\texttt{Quot.sound}, and checks that the two statements agree when printed with all implicit arguments shown (\texttt{pp.all}). All 76 claims
pass. The check uses Lean~\texttt{4.34.0-rc2} and the Mathlib fork \texttt{Deicyde/mathlib4} at
\texttt{2b73d98}, which is upstream \texttt{0383a80e64} plus the commit of~\cite{PR}. The six
statements of the earlier comparator challenge in the same repository pass the comparator
\url{https://github.com/leanprover/comparator} at revision \texttt{19e111e}.

\begin{center}
\small
\begin{tabular}{lp{0.62\textwidth}r}
\hline
Part & Ledger claims & Number\\
\hline
Sections 2--3 & \ref{thm:bundle-lifting}--\ref{prop:homogeneous-lift} & 5\\
Section 4 & \ref{prop:split-lift}--\ref{cor:bundle-cases}, with both parts of \ref{thm:summable-parameters} & 7\\
Section 5 & \ref{prop:exterior}, \ref{thm:duality} & 2\\
Section 6 & both parts of \ref{thm:operator-obstructions}; \ref{cor:classification}--\ref{prop:tangent-obstruction}; the bundle statement after \ref{prop:shear} & 7\\
Section 7 & \ref{thm:ambient-forms}, \ref{prop:ambient-agreement} & 2\\
Section 9 & zero dual of universal targets & 1\\
Appendices A--B & \ref{lem:unique}--\ref{lem:findim}, \ref{lem:detinj}--\ref{lem:D}, \ref{rem:FFunit} & 16\\
Appendix C & \ref{thm:laurent}--\ref{lem:Psi}, \ref{lau:prop:c0-multipliers}--\ref{lau:cor:dichotomy} & 7\\
Appendices D--E & \ref{lem:discdist}--\ref{prop:ascent}, \ref{lem:descent}, \ref{prop:nonultra}, \ref{rem:spaces} & 12\\
Appendix F & \ref{thm:scalar-counter}--\ref{scalar:certificate} & 5\\
Appendix G & \ref{fam:prop:tensor-homogeneous}, \ref{fam:thm:tensor-analytic}--\ref{fam:cor:no-common-radius} & 4\\
Appendix H & \ref{dom:split-incoming}--\ref{dom:degree-gap} & 5\\
Appendix I & \ref{prop:scalar-countable}--\ref{rem:nonspherical-analytic-target} & 3\\
\hline
\end{tabular}
\end{center}

Apart from Proposition~\ref{prop:tangent-values}, the ledger covers every theorem, proposition, lemma and corollary. It also covers three remarks,
the bundle statement after Proposition~\ref{prop:shear}, and one claim of this section. It does not cover the definitions or
Remarks~\ref{rem:FFhyp}, \ref{rem:infkappa} and~\ref{rem:ascent}. Proposition~\ref{prop:tangent-values}, the coefficient bound on retracts in Theorem~\ref{thm:summable-parameters}(1), and the remaining unnumbered arguments of the text are outside the ledger. Apart from the declarations named in Section~\ref{sec:mathlib-interface}, they have written proofs only.

Some ledger statements differ from the text, and \texttt{challenge.md} and the docstrings of the challenge files list the differences. The main
ones are these. Theorem~\ref{thm:ambient-forms} and Proposition~\ref{prop:ambient-agreement} are stated
on open subsets of a normed space, since Mathlib has no differential forms on manifolds. The Lean
statement of Theorem~\ref{thm:bundle-lifting} groups the variables by variance, uses the canonical trivializations, and
does not state independence of the trivializing cover. Corollary~\ref{lau:cor:diagonal-bundle} is
stated for the explicit transition and section maps, and Corollary~\ref{cor:bundle-cases} for objects
and morphisms without the functor laws.

\appendix

\section{Preliminaries for the obstruction proofs}\label{app:obstruction-preliminaries}
The constructions below use three elementary facts about valued fields and power series. We record their proofs to make explicit the scope of the coefficient argument, including incomplete base fields.

\begin{lemma}[one-variable uniqueness]\label{lem:unique}
Let $K$ be nontrivially normed and $Z$ a normed $K$\nobreakdash-space, not necessarily complete. Let $z_0,z_1,\ldots\in Z$, $\rho>0$
and $M\ge0$ with $\norm{z_n}\rho^n\le M$ for all $n$. If $\sum_nt^nz_n$ converges to $0$ for every $t\in K$ with
$\abs t<\rho$, then $z_n=0$ for all $n$.
\end{lemma}

\begin{proof}
Otherwise let $N$ be least with $z_N\neq0$. For $0<\abs t<\rho/2$ the series $\sum_{n>N}t^{n-N}z_n$ converges, being
$t^{-N}$ times the convergent tail $\sum_{n>N}t^nz_n$, and $t^N\bigl(z_N+\sum_{n>N}t^{n-N}z_n\bigr)=0$. Hence
\[
\norm{z_N}=\Bigl\lVert\sum_{n>N}t^{n-N}z_n\Bigr\rVert\le\sum_{n>N}M\rho^{-n}\abs t^{n-N}
=M\rho^{-N}\frac{\abs t/\rho}{1-\abs t/\rho},
\]
where the norm of a convergent series is bounded by the sum of the norms by continuity of the norm on partial sums.
Since $K$ is nontrivially normed, there are $t\ne0$ with $\abs t$ arbitrarily small, so $\norm{z_N}=0$, a
contradiction.
\end{proof}

\begin{lemma}\label{lem:ultra}
A normed field of characteristic $p>0$ is nonarchimedean.
\end{lemma}

\begin{proof}
Frobenius and the triangle inequality give, for every $m\ge0$,
\[
 \abs{x+y}^{p^m}=\abs{x^{p^m}+y^{p^m}}
 \le2\max(\abs x,\abs y)^{p^m}.
\]
Taking $p^m$\nobreakdash-th roots and letting $m\to\infty$ proves the assertion.
\end{proof}

\begin{lemma}\label{lem:findim}
Let $\mathbb K$ be a complete nontrivially normed field and $U$ a finite-dimensional normed $\mathbb K$\nobreakdash-space. Then every
linear functional on $U$ is bounded, and $U$ is complete.
\end{lemma}

\begin{proof}
We prove both assertions by induction on $n=\dim U$, the case $n=0$ being immediate. Suppose $n\ge1$ and let $g\ne0$
be a linear functional on $U$. Its kernel $H$ has dimension $n-1$, so the induction hypothesis makes $H$ complete
and therefore closed in $U$. Choose $v$ with $g(v)=1$. Since $v\notin H$,
$\delta:=\inf_{h\in H}\norm{v-h}>0$. Write $x=g(x)v+h$ with $h\in H$. If $g(x)\ne0$, then
\[
\norm x=\abs{g(x)}\,\norm{v+g(x)^{-1}h}\ge\abs{g(x)}\delta.
\]
Thus $\abs{g(x)}\le\delta^{-1}\norm x$ for every $x$, proving boundedness.

Choose a basis of $U$. Its coordinate functionals are now bounded, so the coordinate map
$U\to\mathbb K^n$, with the max norm on $\mathbb K^n$, is bounded. The inverse map
$(c_i)\mapsto\sum_ic_ie_i$ is bounded by $\sum_i\norm{e_i}$. Hence $U$ is bi-Lipschitz equivalent to the complete
space $\mathbb K^n$, and is itself complete.
\end{proof}

When the scalar field needs to be displayed, $A^{k,L}$ denotes the degree-$k$ precomposition action formed over $L$. For the precomposition action $A$, a \emph{bounded lift} means a bounded $k$\nobreakdash-linear map on its operator domain whose diagonal is $A$, as in Proposition~\ref{prop:homogeneous-lift}. That proposition applies at every base point; Lemma~\ref{lem:unique} provides the one-variable coefficient comparison used in its proof.

\section{The multiplier theorem over a finite field}\label{app:FF}\label{sec:FF}

This section proves the algebraic obstruction used in the Laurent-series construction: over a finite field in which
$k!=0$, coordinatewise multiplication admits no multilinear lift satisfying a polarized diagonal identity and a
uniform bound on exterior support. We give a roadmap of the complete construction, introduce the exterior
support estimates, and then distinguish the diagonal identities: over a finite field equality at all scalar
points need not be a polynomial identity.

\subsection*{Roadmap of the Banach counterexample}
The proof separates into an algebraic obstruction, its realization by a normed space, and passage to the prescribed
field. The analytic input is a lift criterion valid in every characteristic: if $A$ is analytic at any point, the
degree-$k$ term of its power series has diagonal $A$. One-variable uniqueness identifies this diagonal without
dividing by $k!$ (Proposition~\ref{prop:homogeneous-lift}).

\emph{The finite-field obstruction.}
Let $\kappa$ be finite and put $E_0=\kappa^\N$. Theorem~\ref{thm:FF} rules out a multilinear lift of coordinatewise
multiplication satisfying the polarized diagonal identity~\eqref{eq:Pol1}, with values in $\Lambda^k_\kappa E_0$ when $k!=0$ and every value of the lift lies in an exterior power of a
subspace of uniformly bounded dimension. Ramsey's theorem first makes the coefficients depend only on the order
and equality pattern of their indices. Signed sums over clusters of four indices cancel the terms not tied to the
output. A staircase of clusters converts the support bound into invariance under adjacent transpositions; the
polarized diagonal identity then forces $k!\,\chi=1$ for a common cluster value $\chi$, a contradiction.

\emph{The Laurent series construction.}
Over $K_1=\kappa((X))$, take $E_1=\linf(\N,K_1)$ and let $B$ be the completion of
$\Lambda^k_{K_1}E_1$ under its projective norm, defined using ordinary sums rather than maxima. The essential
estimate is that the constant-coefficient map
\[
\eta\colon B\longrightarrow\Lambda^k_\kappa E_0
\]
is well defined and has uniformly bounded support dimension on bounded sets
(Proposition~\ref{prop:supp}). Evaluating a hypothetical lift on multiplication operators and on the universal
wedge map, then applying $\eta$, produces exactly the finite-field lift excluded above. The coefficient map
is $\kappa$\nobreakdash-linear; no $K_1$\nobreakdash-linearity is asserted or needed.

\emph{Passage to the prescribed field.}
Choose $t\in K$ with $0<\abs t<1$. The closure of $\F_p(t)$ in the completion $\Kh$ is a spherically complete
Laurent series field $K_1\cong\F_p((X))$, with $\abs X=\abs t$. Ingleton's Hahn--Banach theorem~\cite{In} gives a
contracting $K_1$\nobreakdash-linear projection $\Kh\to K_1$. It makes the completed projective tensor products
\[
E=E_1\widehat\otimes_{\pi,K_1}\Kh,\qquad
F=B\widehat\otimes_{\pi,K_1}\Kh
\]
contain $E_1$ and $B$ isometrically, and allows a lift over $\Kh$ to descend to one over $K_1$. Density handles an
incomplete $K$. These are the explicit witnesses in Theorem~\ref{thm:operator-obstructions}(1). Finally, sums of unit wedges on disjoint
blocks have norms growing linearly with the number of blocks. This growth persists in $F$ and rules out every
equivalent nonarchimedean norm (Proposition~\ref{prop:nonultra}).

\subsection{Exterior powers, supports and determinant arrays}\label{ssec:ext}
Let $L$ be a field, $V$ an $L$\nobreakdash-vector space and $k\ge1$. We write $\Lambda^kV=\Lambda^k_LV$ for the algebraic exterior
power, in which $x_1\wedge\dots\wedge x_k=0$ whenever two $x_i$ coincide. We use the standard facts: for a subspace
$W\subseteq V$ the map $\Lambda^kW\to\Lambda^kV$ is injective, and we regard $\Lambda^kW$ as a subspace of $\Lambda^kV$;
if $u_1,\dots,u_d$ is a basis of a finite-dimensional space $U$, the wedges $u_I:=u_{i_1}\wedge\dots\wedge u_{i_k}$,
$I=\{i_1<\dots<i_k\}$, form a basis of $\Lambda^kU$; and every $\omega\in\Lambda^kV$ lies in $\Lambda^kU$ for some
finite-dimensional $U\subseteq V$.

\begin{definition}\label{def:supp}
The \emph{support dimension} of $\omega\in\Lambda^kV$ is
\[
\sdim(\omega):=\min\{\dim W:\ W\subseteq V\text{ finite-dimensional},\ \omega\in\Lambda^kW\}.
\]
\end{definition}

The definition gives $\sdim(0)=0$, $\sdim(c\omega)=\sdim(\omega)$ for $c\neq0$, $\sdim(\omega+\omega')\le\sdim(\omega)+\sdim(\omega')$
(use $W+W'$), and $\sdim(\omega)\le\dim\spn S$ if $\omega$ is a sum of wedges of vectors from a finite set $S$. If
$V\subseteq V'$, then $\sdim$ computed in $V'$ is at most $\sdim$ computed in $V$.

For spaces of functions, exterior vectors can be detected by determinants of coordinate evaluations. These arrays
also turn bounds on support dimension into ordinary matrix-rank bounds.
Let $S$ be a set and $V\subseteq L^S$ a space of functions on $S$. The \emph{determinant array} is the linear map
\[
\Omega\colon\Lambda^kV\to L^{S^k},\qquad \Omega_{y_1\wedge\dots\wedge y_k}(c_1,\dots,c_k):=\det\bigl(y_b(c_a)\bigr)_{a,b=1}^k .
\]
It is well defined because the right side is multilinear and alternating in $(y_1,\dots,y_k)$ (these are the
columns). For each $\omega$, the function $\Omega_\omega$ is alternating in $(c_1,\dots,c_k)$ (the rows); in particular
$\Omega_\omega(c_{\sigma(1)},\dots,c_{\sigma(k)})=\sgn(\sigma)\,\Omega_\omega(c_1,\dots,c_k)$.

\begin{lemma}\label{lem:detinj}
The determinant array $\Omega\colon\Lambda^kV\to L^{S^k}$ is injective.
\end{lemma}

\begin{proof}
Let $\Omega_\omega=0$ and choose a finite-dimensional $U\subseteq V$ with $\omega\in\Lambda^kU$. The evaluations
$\mathrm{ev}_s|_U$, $s\in S$, separate the points of $U$ (a function in $U$ vanishing at every $s$ is $0$), so they span
the dual space $U^*$. Choose $s_1,\dots,s_d\in S$ such that the $\mathrm{ev}_{s_i}|_U$ form a basis of $U^*$, and let
$u_1,\dots,u_d$ be the dual basis of $U$, so that $u_j(s_i)=\delta_{ij}$. For increasing $k$\nobreakdash-subsets $I,J$ of
$\{1,\dots,d\}$ we get $\Omega_{u_I}(s_{j_1},\dots,s_{j_k})=\det(u_{i_b}(s_{j_a}))_{a,b}=\delta_{IJ}$. Writing
$\omega=\sum_Ic_Iu_I$ therefore gives $0=\Omega_\omega(s_{j_1},\dots,s_{j_k})=c_J$ for every $J$, so $\omega=0$.
\end{proof}

\begin{lemma}[support bounds flattening rank]\label{lem:flat}
Let $V\subseteq L^S$, $\omega\in\Lambda^kV$ and $\alpha\in\{1,\dots,k\}$. Let $R\subseteq S$ and
$R'\subseteq S^{\{1,\dots,k\}\setminus\{\alpha\}}$ be finite. The $R\times R'$ matrix $M$ with entries
$M(c_\alpha,c'):=\Omega_\omega(c)$, where $c$ has $c_\alpha$ in slot $\alpha$ and the entries of $c'$ in the other
slots, has rank at most $\sdim(\omega)$.
\end{lemma}

\begin{proof}
Write $\omega=\sum_jy_{j,1}\wedge\dots\wedge y_{j,k}$ with all $y_{j,b}$ in a subspace $W$ of dimension $\sdim(\omega)$.
Expanding each determinant along row $\alpha$,
$\det(y_{j,b}(c_a))_{a,b}=\sum_b\pm y_{j,b}(c_\alpha)\cdot\mu_{j,b}(c')$, where the minor $\mu_{j,b}(c')$ does not
involve $c_\alpha$. So for fixed $c'$ the column $c_\alpha\mapsto M(c_\alpha,c')$ is a linear combination of the
restrictions $y_{j,b}|_R$. All columns lie in the restriction of $W$ to $R$, a space of dimension at most
$\sdim(\omega)$.
\end{proof}

A second way to obtain lower bounds on support is to contract all but one exterior factor. The resulting vectors
must lie in every subspace supporting the original exterior vector.

\begin{samepage}
\begin{lemma}[contraction]\label{lem:contr}
Let $k\ge2$ and $\varphi=(\varphi_1,\dots,\varphi_{k-1})\in(V^*)^{k-1}$. There is a linear map $c_\varphi\colon\Lambda^kV\to V$
with
\[
c_\varphi(x_1\wedge\dots\wedge x_k)=\sum_{i=1}^k(-1)^{k+i}\det\bigl(\varphi_a(x_b)\bigr)_{1\le a\le k-1,\ b\ne i}\;x_i ,
\]
and $c_\varphi(\Lambda^kW)\subseteq W$ for every subspace $W$. Consequently
$\sdim(\omega)\ge\dim\spn\{c_\varphi(\omega):\varphi\in(V^*)^{k-1}\}$.
\end{lemma}
\end{samepage}

\begin{proof}
The right side is multilinear in $(x_b)$. For $\psi\in V^*$, applying $\psi$ to it gives, by Laplace expansion along the
last row, the determinant of the $k\times k$ matrix with entries $\varphi_a(x_b)$ in the rows $a<k$ and $\psi(x_b)$ in the last
row; it vanishes when two columns coincide. Since $V^*$ separates the points of $V$, the right side
vanishes whenever $x_a=x_b$ for some $a\ne b$, so it factors through $\Lambda^kV$. Each term is a multiple of some $x_i$, which gives
$c_\varphi(\Lambda^kW)\subseteq W$. If $\omega\in\Lambda^kW$ with $\dim W=\sdim(\omega)$, all $c_\varphi(\omega)$ lie in $W$.
\end{proof}

\begin{proposition}[canonical exterior support]\label{prop:canonical-support}
Let $L$ be any field, $V$ an $L$\nobreakdash-vector space, $k\ge2$, and
$\omega\in\Lambda^kV$. Then
\[
S(\omega):=\spn\{c_\varphi(\omega):\varphi\in(V^*)^{k-1}\}
\]
is the unique smallest subspace supporting $\omega$. In particular,
$\omega\in\Lambda^k S(\omega)$ and
$\sdim(\omega)=\dim S(\omega)$.
\end{proposition}

\begin{proof}
Choose a finite-dimensional supporting subspace $U\subseteq V$.
Lemma~\ref{lem:contr} gives $S(\omega)\subseteq U$, so $S(\omega)$ is
finite-dimensional. Choose a basis $e_1,\dots,e_d$ of $U$ whose first
$r$ vectors form a basis of $S(\omega)$, and extend its coordinate
functionals algebraically to $V$. Write
$\omega=\sum_{|I|=k}a_Ie_I$ in the exterior basis of $U$.
If $j\in I$ and $j>r$, contract against the $k-1$ coordinate functionals
indexed by $I\setminus\{j\}$. The $e_j$\nobreakdash-coordinate of this
contraction is $\pm a_I$: no other basis wedge contributes to that
coordinate. The contraction belongs to $S(\omega)$, so this coordinate
is zero. Hence $a_I=0$ whenever $I$ meets $\{r+1,\dots,d\}$, proving
$\omega\in\Lambda^k S(\omega)$.
Every supporting subspace contains all the contractions by
Lemma~\ref{lem:contr}, and therefore contains $S(\omega)$. This proves
minimality and the dimension formula. Only unit signs occur, so the
argument is valid in every characteristic.
\end{proof}

The same conclusion holds in degree one: under $\Lambda^1V\simeq V$,
the empty-tuple contraction is the identity and $S(\omega)=L\omega$.
In degree zero, every element of $\Lambda^0V=L$ is supported by the
zero subspace, so its support dimension is zero; no contraction to
$V$ is needed.

For matrix calculations, choose any finite-dimensional supporting
subspace $U$ and a basis of it. Form the finite matrix whose columns
are the coordinates of the contractions against $(k-1)$\nobreakdash-tuples
of coordinate functionals, extended to $V$. Arbitrary functionals
restrict to linear combinations of these coordinates on $U$;
multilinearity in the covectors therefore shows that the column span
is exactly $S(\omega)$. The matrix consequently has rank
$\sdim(\omega)$. This uses ordinary matrix rank, and entails no
choice of a rank convention for an infinite matrix.

\subsection{Lifts and three diagonal identities}\label{ssec:diag}
Fix a field $L$, an integer $k\ge1$, $L$\nobreakdash-vector spaces $A$ and $V$, and a linear map $D\colon A\to\End_L(V)$,
$a\mapsto D_a$. A \emph{lift of $D$ in degree $k$} is a $2k$\nobreakdash-linear map $\Psi\colon A^k\times V^k\to\Lambda^kV$ over $L$ that is
alternating in the $V$\nobreakdash-slots; its diagonal identity will be specified separately. For every $n\ge1$ write
$[n]:=\{1,\dots,n\}$, and for $\alpha\in\N^m$ put $\abs\alpha:=\sum_i\alpha_i$. Given $m\ge1$, a map
$f\colon[k]\to[m]$, a family $b=(b_1,\dots,b_m)\in A^m$ and $\lambda\in L^m$, put
\[
b_f:=(b_{f(1)},\dots,b_{f(k)}),\quad \type f:=\bigl(\abs{f^{-1}(1)},\dots,\abs{f^{-1}(m)}\bigr)\in\N^m,\quad
\lambda_f:=\prod_{j=1}^k\lambda_{f(j)}=\lambda^{\type f}.
\]
The \emph{pointwise identity} is
\begin{equation*}\tag{Pw}\label{eq:Pw}
\Psi(a,\dots,a;x_1,\dots,x_k)=D_ax_1\wedge\dots\wedge D_ax_k\qquad\text{for all }a\in A,\ x\in V^k.
\end{equation*}
The \emph{polarized identity} is: for all $m\ge1$, $b\in A^m$, $\alpha\in\N^m$ with $\abs\alpha=k$, and $x\in V^k$,
\begin{equation*}\tag{Pol}\label{eq:Pol}
\sum_{\type f=\alpha}\Psi(b_f;x)=\sum_{\type f=\alpha}D_{b_{f(1)}}x_1\wedge\dots\wedge D_{b_{f(k)}}x_k .
\end{equation*}
Its \emph{multilinear instance} is: for all $b_1,\dots,b_k\in A$ and $x\in V^k$,
\begin{equation*}\tag{Pol1}\label{eq:Pol1}
\sum_{\sigma\in S_k}\Psi(b_{\sigma(1)},\dots,b_{\sigma(k)};x)
=\sum_{\sigma\in S_k}D_{b_{\sigma(1)}}x_1\wedge\dots\wedge D_{b_{\sigma(k)}}x_k .
\end{equation*}
Here $\sum_{\type f=\alpha}$ runs over the maps $f\colon[k]\to[m]$ of type $\alpha$. \eqref{eq:Pol1} is the instance $m=k$,
$\alpha=(1,\dots,1)$ of \eqref{eq:Pol}, since the maps of type $(1,\dots,1)$ are the permutations. For $b\in A^m$ and
$x\in V^k$ put $a_\lambda:=\sum_i\lambda_ib_i$ and
\begin{align*}
\Delta_{b,x}(\lambda)&:=\Psi(a_\lambda,\dots,a_\lambda;x)-D_{a_\lambda}x_1\wedge\dots\wedge D_{a_\lambda}x_k
=\sum_{\abs\alpha=k}\lambda^\alpha\,\delta_\alpha(b,x),\\
\delta_\alpha(b,x)&:=\sum_{\type f=\alpha}\Bigl[\Psi(b_f;x)-D_{b_{f(1)}}x_1\wedge\dots\wedge D_{b_{f(k)}}x_k\Bigr].
\end{align*}
Multilinearity of $\Psi$ in the $A$\nobreakdash-slots, linearity of $D$ and multilinearity of the wedge give the displayed
expansion. Thus $\Delta_{b,x}$ is a homogeneous polynomial of degree $k$ with coefficients in $\Lambda^kV$.
Identity~\eqref{eq:Pol} says that every coefficient $\delta_\alpha(b,x)$ is zero. Identity~\eqref{eq:Pw} says instead that
$\Delta_{b,x}$ vanishes at every point of $L^m$; conversely, this pointwise vanishing gives~\eqref{eq:Pw} by taking
$m=1$ and $\lambda=1$. The following two identity principles identify the cases in which pointwise vanishing forces
all coefficients to vanish.

\begin{lemma}[identity principle]\label{lem:idprinc}
Let $\mathbb K$ be an infinite field, $Y$ a $\mathbb K$\nobreakdash-vector space and $g(\lambda)=\sum_\alpha\lambda^\alpha y_\alpha$ a
polynomial in $\lambda\in\mathbb K^m$ with finitely many nonzero coefficients $y_\alpha\in Y$. If $g(\lambda)=0$ for all
$\lambda\in\mathbb K^m$, then every $y_\alpha=0$.
\end{lemma}

\begin{proof}
The finitely many coefficients $y_\alpha$ span a finite-dimensional subspace of $Y$. Applying its coordinate
functionals reduces the claim to scalar polynomials. For these, induction on $m$ reduces to the fact that a nonzero
polynomial in one variable has only finitely many roots.
\end{proof}

\begin{lemma}[homogeneous forms over a finite field]\label{lem:forms}
Let $q$ be a prime power, $L'\supseteq\F_q$ a field, $1\le e\le q$, and $g\in L'[\lambda_1,\dots,\lambda_m]$ homogeneous of
degree $e$ with $g(\lambda)=0$ for every $\lambda\in\F_q^m$. Then $g=0$.
\end{lemma}

\begin{proof}
Call a monomial $\lambda^\beta$ \emph{reduced} if every $\beta_i\le q-1$. The indicator function of a point
$a\in\F_q^m$ is $\prod_i\bigl(1-(\lambda_i-a_i)^{q-1}\bigr)$, which expands into reduced monomials. So the $q^m$ reduced
monomials span the $q^m$\nobreakdash-dimensional $L'$\nobreakdash-space of functions $\F_q^m\to L'$, and they are linearly independent as
functions. A monomial of degree $e\le q$ that is not reduced is $\lambda_i^q$ (and then $e=q\ge2$); as a function on
$\F_q^m$ it agrees with the reduced monomial $\lambda_i$, of degree $1<e$. If $e=q$, replace each $\lambda_i^q$ in $g$ by
$\lambda_i$. These degree-one images are distinct, and homogeneity separates them from all the remaining monomials,
which still have degree $q$. The replacement therefore sends the monomials of $g$ injectively to reduced monomials
without changing the associated function. Linear independence of the reduced monomials forces every coefficient to
vanish. If $e<q$, no replacement is needed.
\end{proof}

\begin{proposition}[pointwise versus polarized]\label{prop:diag}
Let $\Psi$ be a lift of $D$ in degree $k$ over a field $L$.
\begin{enumerate}
\item \eqref{eq:Pol}$\Rightarrow$\eqref{eq:Pw}$\Rightarrow$\eqref{eq:Pol1}, over every field.
\item If $L$ is infinite, or if $L=\F_q$ and $k\le q$, then \eqref{eq:Pw}$\Rightarrow$\eqref{eq:Pol}.
\item If $L=\F_q$ and $k\ge q+1$, then \eqref{eq:Pw} does not imply \eqref{eq:Pol} in general.
\item If $k!=0$ in $L$, then \eqref{eq:Pol1} does not imply \eqref{eq:Pw} in general.
\end{enumerate}
\end{proposition}

\begin{proof}
\emph{(1) Inclusion--exclusion.} The implication \eqref{eq:Pol}$\Rightarrow$\eqref{eq:Pw} is the case $m=1$ and
$\alpha=(k)$. To prove \eqref{eq:Pw}$\Rightarrow$\eqref{eq:Pol1}, fix $b_1,\dots,b_k\in A$ and $x\in V^k$, and put
$b_S:=\sum_{i\in S}b_i$ for $S\subseteq[k]$. Multilinearity gives
\begin{align*}
\sum_{S\subseteq[k]}(-1)^{k-\abs S}\Psi(b_S,\dots,b_S;x)&=\sum_S(-1)^{k-\abs S}\sum_{f\colon[k]\to S}\Psi(b_f;x)\\
&=\sum_{f\colon[k]\to[k]}\Psi(b_f;x)\sum_{S\supseteq f([k])}(-1)^{k-\abs S}=\sum_{f\text{ bijective}}\Psi(b_f;x),
\end{align*}
because $\sum_{S\supseteq T}(-1)^{k-\abs S}=(1-1)^{k-\abs T}$ is $1$ if $T=[k]$ and $0$ otherwise. The same computation
applies to the multilinear map $(a_1,\dots,a_k)\mapsto D_{a_1}x_1\wedge\dots\wedge D_{a_k}x_k$, whose diagonal is the right
side of \eqref{eq:Pw}. Evaluating \eqref{eq:Pw} at the $2^k$ subset sums $b_S$ therefore gives \eqref{eq:Pol1}.
The computation is valid in every characteristic and involves no division by $k!$.

\emph{(2) The identity principles.} The coefficients of $\Delta_{b,x}$ span a finite-dimensional subspace of $\Lambda^kV$; in a basis of it each
coordinate of $\Delta_{b,x}$ is a scalar form of degree $k$ vanishing on $L^m$. It is zero by Lemma~\ref{lem:idprinc}
if $L$ is infinite, and by Lemma~\ref{lem:forms} if $L=\F_q$ and $k\le q$.

\emph{(3) Pointwise vanishing without full polarization.} Let $\dim A\ge3$, with $b_1,b_2,b_3\in A$ and functionals $g_1,g_2,g_3$ on $A$ such that
$g_i(b_j)=\delta_{ij}$, and let $y_1,\dots,y_k\in V$ be linearly independent. Put
\[
\theta(a_1,\dots,a_k):=\bigl[g_1(a_1)\cdots g_1(a_q)\,g_2(a_{q+1})-g_1(a_1)\,g_2(a_2)\cdots g_2(a_{q+1})\bigr]
\cdot g_3(a_{q+2})\cdots g_3(a_k)
\]
(for $k=q+1$ the $g_3$\nobreakdash-factor is empty and $\dim A\ge2$ suffices), and $\Theta(a;x):=\theta(a)\,x_1\wedge\dots\wedge x_k$.
Then $\Theta$ is a lift of $D=0$. It satisfies \eqref{eq:Pw}, because
$\theta(a,\dots,a)=(g_1(a)^qg_2(a)-g_1(a)g_2(a)^q)\,g_3(a)^{k-q-1}=0$ since $c^q=c$ for $c\in\F_q$. It
violates \eqref{eq:Pol} for $\alpha=(q,1,k-q-1)$ and $b=(b_1,b_2,b_3)$: the first product contributes to
$\sum_{\type f=\alpha}\theta(b_f)$ exactly for $f=(1,\dots,1,2,3,\dots,3)$, with value $1$; the second product would need
$f(2)=\dots=f(q+1)=2$, that is, at least $q\ge2$ values equal to $2$, which is impossible for type $\alpha$. So the
coefficient $\delta_\alpha(b,y)$ of $\Theta$ is $y_1\wedge\dots\wedge y_k\neq0$. When $k=q+1$, the third label and its
functional can be omitted throughout, with $\alpha=(q,1)$.

\emph{(4) The multilinear identity without the pointwise identity.} Take $\bar a\in A$ and a functional $g_1$ with $g_1(\bar a)=1$, and $y_1,\dots,y_k$ linearly independent. The
lift $\Theta'(a;x):=g_1(a_1)\cdots g_1(a_k)\,x_1\wedge\dots\wedge x_k$ of $D=0$ satisfies \eqref{eq:Pol1}, since
$\sum_\sigma\Theta'(b_\sigma;x)=k!\,g_1(b_1)\cdots g_1(b_k)\,x_1\wedge\dots\wedge x_k=0$ for all
$b_1,\dots,b_k\in A$ and $x\in V^k$, but $\Theta'(\bar a,\dots,\bar a;y)=y_1\wedge\dots\wedge y_k\neq0$.
\end{proof}

The finite-field obstruction below requires only~\eqref{eq:Pol1}, which follows from either~\eqref{eq:Pw}
or~\eqref{eq:Pol}. The distinction between the latter two remains important: over $\F_q$ they need not agree when
$k>q$. The lift constructed in Appendix~\ref{sec:laurent} nevertheless satisfies the full identity~\eqref{eq:Pol} in every
degree: the identity is polarized over an infinite Laurent-series field before constant coefficients are taken
(Lemma~\ref{lem:Psi}).

\subsection{Statement and proof strategy}
For the proof of the theorem, assume that $L$ is a \emph{finite} field with $k!=0$ in $L$, that is, $\operatorname{char}L=p\le k$; in
particular $k\ge2$. Put $V:=L^\N$ and let $V_{\mathrm{fin}}:=L^{(\N)}\subseteq V$ be the subspace of finitely supported
sequences. For $u,x\in V$ we write $ux$ for the coordinatewise product; so $D_ux=ux$ is the multiplier family.

\Needspace{16\baselineskip}
\begin{samepage}
\begin{theorem}[multipliers over a finite field]\label{thm:FF}
Let $L$ be a finite field and $k\ge1$ with $k!=0$ in $L$.
\begin{enumerate}
\item There is no $2k$\nobreakdash-linear map $\Psi\colon V_{\mathrm{fin}}^k\times V_{\mathrm{fin}}^k\to\Lambda^k_LV$ over $L$ such that
\begin{enumerate}
\item $\Psi$ is antisymmetric in the last $k$ slots;
\item $\Psi$ satisfies \eqref{eq:Pol1} for $D_ux=ux$: for all $u_1,\dots,u_k,x_1,\dots,x_k\in V_{\mathrm{fin}}$,
\[
\sum_{\sigma\in S_k}\Psi(u_{\sigma(1)},\dots,u_{\sigma(k)};x_1,\dots,x_k)=\sum_{\sigma\in S_k}(u_{\sigma(1)}x_1)\wedge\dots\wedge(u_{\sigma(k)}x_k);
\]
\item for some $d\in\N$, $\sdim(\Psi(u;x))\le d$ for all $u,x\in V_{\mathrm{fin}}^k$ with entries in $\{0,\pm1\}$.
\end{enumerate}
\item Let $E_0=V$ or $E_0=V_{\mathrm{fin}}$. There is no $2k$\nobreakdash-linear map $\Psi\colon E_0^k\times E_0^k\to\Lambda^k_LE_0$ over $L$
that is antisymmetric in the last $k$ slots, satisfies \eqref{eq:Pol1} for all finitely supported $u,x$, and has
$\sdim(\Psi(u;x))\le d$ for all $u,x\in E_0^k$, for some $d$.
\end{enumerate}
\end{theorem}
\end{samepage}

\emph{Reduction to (1).} Given a map as in (2), restrict its inputs to
$V_{\mathrm{fin}}^k\times V_{\mathrm{fin}}^k$ and compose its output with the injective map
$\Lambda^kE_0\to\Lambda^kV$. This preserves antisymmetry and~\eqref{eq:Pol1}. Support dimension can only decrease under
the inclusion of the ambient space (\S\ref{ssec:ext}), so the resulting map satisfies (a)--(c).

We therefore prove (1). Suppose that $\Psi$ satisfies its hypotheses, and fix the bound $d$ in (c). Write
$e_i\in V_{\mathrm{fin}}$ for the unit vectors, $e_a:=(e_{a_1},\dots,e_{a_k})$ for $a\in\N^k$, and $\Omega$ for the
determinant array of \S\ref{ssec:ext} with $S=\N$.

The proof has three steps. Ramsey's theorem first makes the coefficients of $\Omega\Psi$ depend only on the relative
order and equalities among their indices. Special weights on four-point clusters then cancel every contribution
from an occupied cluster containing no output index. This cancellation produces a two-valued staircase matrix from
one exterior vector. Its bounded support forces the two values to agree, and hence makes all cluster order patterns
have a common value. Finally, antisymmetry and~\eqref{eq:Pol1} say that the sum of these $k!$ equal values is $1$,
contradicting $k!=0$ in $L$.

\subsection{Order-homogeneity from Ramsey's theorem}
We use Ramsey's theorem~\cite{Ra} in its infinite form: \emph{for every $n\ge1$ and every coloring of the
$n$\nobreakdash-element subsets of $\N$ with finitely many colors there is an infinite $H\subseteq\N$ all of whose $n$\nobreakdash-element
subsets have the same color.}

For a tuple $z\in\N^N$ with distinct values $y_1<\dots<y_r$, its \emph{pattern} is the map
$i\mapsto j$ determined by $z_i=y_j$. In particular, repeated entries are retained in the pattern. Two tuples $z$
and $z'$ have the same pattern if and only if $z_i<z_j\Leftrightarrow z'_i<z'_j$ and $z_i=z_j\Leftrightarrow z'_i=z'_j$ for
all $i,j$.

\begin{lemma}[pattern homogeneity]\label{lem:H}
For every $N\ge1$ and every function $T\colon\N^N\to L$ there is an infinite $H\subseteq\N$ such that $T(z)=T(z')$
whenever $z,z'\in H^N$ have the same pattern.
\end{lemma}

\begin{proof}
Color an $N$\nobreakdash-element set $s=\{s_1<\dots<s_N\}\subseteq\N$ by the vector
\[
\bigl(T(s_{\beta(1)},\dots,s_{\beta(N)})\bigr)_{\beta\colon[N]\to[N]}\in L^{N^N}.
\]
The coloring records all maps $\beta$, including those that repeat indices. It has finitely many values because
$L$ is finite; this is the only use of the finiteness of $L$ in the obstruction proof. Ramsey's theorem gives an
infinite $H$ on which the color is constant, say equal to $(t_\beta)_\beta$. Let $z\in H^N$ have value set
$\{y_1<\dots<y_r\}$. Since $H$ is unbounded, add $N-r$ elements of $H$ above $y_r$ to obtain an
$N$\nobreakdash-element set $s\subseteq H$. Padding above the largest value leaves the ranks unchanged: $s_j=y_j$ for
$j\le r$, and
$z=(s_{\beta(1)},\dots,s_{\beta(N)})$, where $\beta(i)$ is the rank of $z_i$ in $\{y_1,\dots,y_r\}$. The map $\beta$
depends only on the pattern of $z$, and $T(z)=t_\beta$.
\end{proof}

\emph{Coefficients.} For $(a;y;c)\in\N^k\times\N^k\times\N^k$ put
\[
T(a;y;c):=\Omega_{\Psi(e_a;e_y)}(c_1,\dots,c_k)\in L .
\]
Lemma~\ref{lem:H} with $N=3k$ gives an infinite set $H=\{h_0<h_1<h_2<\cdots\}$ on which $T$ depends only on the pattern.
For $u_j,v_j\in V_{\mathrm{fin}}$ supported in $H$, multilinearity of $\Psi$ and linearity of $\Omega$ give
\begin{equation}\label{eq:expand}
\Omega_{\Psi(u;v)}(c)=\sum_{a,y\in H^k}\ \prod_{j=1}^ku_j(a_j)\prod_{j=1}^kv_j(y_j)\cdot T(a;y;c).
\end{equation}
Although the sum is indexed by $H^k\times H^k$, only finitely many summands are nonzero, because all the input
vectors have finite support. Every output tuple $c$ used below also lies in $H^k$, so pattern homogeneity applies to
all $3k$ indices in~\eqref{eq:expand}.

\subsection{Clusters and the cluster value}
A \emph{cluster} is a set $C=\{h_n,h_{n+1},h_{n+2},h_{n+3}\}$ of four
\emph{consecutive} elements of $H$; we write $C(i):=h_{n+i-1}$ for $i=1,\dots,4$. Since a cluster is an interval of
$H$, two disjoint clusters are \emph{comparable}: every element of one is smaller than every element of the other; we
write $C<C'$. Distinct clusters need not be adjacent in $H$.

We choose two weight vectors, one for each kind of input slot:
\begin{align*}
s&=(s_1,s_2,s_3,s_4):=(1,0,-1,0)&&\text{on the operator slots},\\
w&=(w_1,w_2,w_3,w_4):=(1,-1,1,-1)&&\text{on the vector slots}.
\end{align*}
Over $\Z$, hence over every field,
\begin{equation*}\tag{W}\label{eq:W}
\sum_is_i=\sum_iw_i=\sum_is_iw_i=\sum_{i<i'}s_iw_{i'}=\sum_{i>i'}s_iw_{i'}=0,\qquad s_1w_1=1.
\end{equation*}
Indeed $\sum_is_i=1-1$, $\sum_iw_i=1-1+1-1$, $s_1w_1=1$, and
\begin{align*}
\sum_is_iw_i&=s_1w_1+s_3w_3=1-1=0,\\
\sum_{i<i'}s_iw_{i'}&=s_1(w_2+w_3+w_4)+s_3w_4=-1+1=0,\\
\sum_{i>i'}s_iw_{i'}&=s_3(w_1+w_2)=-(1-1)=0 .
\end{align*}
The first two zero sums will cancel a cluster occupied by only one input index. The other three will cancel a
cluster occupied by an operator index and a vector index, according to their relative order. For a cluster $C$ put
\[
s_C:=\sum_{i=1}^4s_ie_{C(i)},\qquad w_C:=\sum_{i=1}^4w_ie_{C(i)} .
\]
These are vectors of $V_{\mathrm{fin}}$ with entries in $\{0,\pm1\}$, and $s_Cw_C=e_{C(1)}-e_{C(3)}$.

\emph{The cluster value.} For pairwise disjoint clusters $C_1,\dots,C_k$, one for each slot, define
\begin{align*}
\chi(C_1,\dots,C_k)&:=\Omega_{\Psi(s_{C_1},\dots,s_{C_k};\,w_{C_1},\dots,w_{C_k})}\bigl(C_1(1),\dots,C_k(1)\bigr)\\
&=\sum_{i,i'\in[4]^k}\ \prod_{j=1}^ks_{i_j}w_{i'_j}\cdot
T\bigl(C_1(i_1),\dots,C_k(i_k);\,C_1(i'_1),\dots,C_k(i'_k);\,C_1(1),\dots,C_k(1)\bigr).
\end{align*}
Let $\tau\in S_k$ be the \emph{order pattern} of $(C_1,\dots,C_k)$: $\tau(j)$ is the position of $C_j$ among
$C_1,\dots,C_k$. The pattern of every $3k$\nobreakdash-tuple in the sum is determined by $(i,i')$ and $\tau$: entries in different
clusters compare as the clusters do, and the entries $C_j(i_j)$, $C_j(i'_j)$, $C_j(1)$ in the same cluster compare as
their indices do, including equalities. Pattern homogeneity therefore shows that the displayed sum depends only
on $\tau$. Every permutation occurs as the order pattern of some disjoint clusters in $H$; we denote the resulting
scalar by $\chi(\tau)$.

\begin{lemma}[cancellation of output-free clusters]\label{lem:C}
Let $\cC_1,\dots,\cC_k$ be finite families of clusters that are pairwise disjoint as sets of clusters, such that all
clusters in $\bigcup_j\cC_j$ are pairwise disjoint. Put $u_j:=\sum_{C\in\cC_j}s_C$ and $v_j:=\sum_{C\in\cC_j}w_C$. Let
$C_l\in\cC_l$ for $l=1,\dots,k$ and $c:=(C_1(1),\dots,C_k(1))$. Then $\Omega_{\Psi(u;v)}(c)=\chi(C_1,\dots,C_k)$.
\end{lemma}

\begin{proof}
We show that every cluster assignment except the one matching the output clusters has total contribution zero.
Expand by~\eqref{eq:expand}. Since the clusters are disjoint, a nonzero term has, for each slot $j$, an operator index
$a_j=\Cop_j(i_j)$ and a vector index $y_j=\Cvec_j(i'_j)$ with $\Cop_j,\Cvec_j\in\cC_j$, and weight
$\prod_js_{i_j}w_{i'_j}$. Group the terms according to the \emph{cluster assignment} $(\Cop_j,\Cvec_j)_{j}$.

Call an occupied cluster \emph{output-free} if it is none of $C_1,\dots,C_k$. Suppose a group has such a cluster
$\Gamma$. Disjointness of the families assigns $\Gamma$ to exactly one slot $j$, so the only tuple entries it can
contain are $a_j$ and $y_j$. Disjointness of the clusters excludes every output index from $\Gamma$.

Fix the positions of all indices outside $\Gamma$ and sum over the positions inside it. Since $\Gamma$ consists of
consecutive elements of $H$, every outside entry lies either below the whole cluster or above it. Moving an index
within $\Gamma$ cannot change any comparison with an outside entry. By pattern homogeneity, the coefficient $T$
therefore depends on the internal positions only as follows:
\begin{enumerate}
\item[(i)] if only $a_j\in\Gamma$, not at all, and the partial sum is $(\sum_is_i)\cdot T=0$;
\item[(ii)] if only $y_j\in\Gamma$, not at all, and the partial sum is $(\sum_iw_i)\cdot T=0$;
\item[(iii)] if both lie in $\Gamma$, only through whether $i_j<i'_j$, $i_j=i'_j$ or $i_j>i'_j$; with the corresponding
values $T_<,T_=,T_>$ of $T$ the partial sum is
\[
T_<\sum_{i<i'}s_iw_{i'}+T_=\sum_is_iw_i+T_>\sum_{i>i'}s_iw_{i'}=0 .
\]
\end{enumerate}
In each case the weights of the other slots are a common factor, and the displayed partial sum vanishes
by~\eqref{eq:W}. Thus every assignment with an output-free cluster contributes zero. In any surviving assignment,
every occupied cluster must be an output cluster. The only output cluster belonging to $\cC_j$ is $C_j$, which forces
$\Cop_j=\Cvec_j=C_j$ for every $j$. The contribution of this one assignment is exactly
$\chi(C_1,\dots,C_k)$.
\end{proof}

\subsection{The staircase and the diagonal}

We now use the support bound to compare two order patterns that differ by an adjacent exchange. The key is to
realize arbitrarily large staircase matrices as flattenings of a single value of $\Psi$.

\begin{lemma}[staircase]\label{lem:S}
Let $\tau\in S_k$, let $\alpha\neq\beta$ be slots with $\tau(\alpha)=\ell$ and $\tau(\beta)=\ell+1$, and let $\tau'$ be $\tau$
with these two positions exchanged, so $\tau'(\alpha)=\ell+1$, $\tau'(\beta)=\ell$ and $\tau'(\gamma)=\tau(\gamma)$ otherwise.
Then $\chi(\tau)=\chi(\tau')$.
\end{lemma}

\begin{proof}
Let $m\ge1$. Choose pairwise disjoint clusters, in increasing order in $H$: first a cluster $C_\gamma$ for each slot
$\gamma$ with $\tau(\gamma)<\ell$, ordered as $\tau$ prescribes; then $2m$ clusters $A_1<B_1<A_2<B_2<\dots<A_m<B_m$; then a
cluster $C_\gamma$ for each slot $\gamma$ with $\tau(\gamma)>\ell+1$, ordered as $\tau$ prescribes. Put
$\cC_\alpha:=\{A_1,\dots,A_m\}$, $\cC_\beta:=\{B_1,\dots,B_m\}$ and $\cC_\gamma:=\{C_\gamma\}$ for the other slots, and let
$u,v$ be as in Lemma~\ref{lem:C} and $\omega:=\Psi(u;v)$. Disjointness of the clusters ensures that the entries of
$u$ and $v$ still lie in $\{0,\pm1\}$, so $\sdim(\omega)\le d$ by hypothesis~(c) of Theorem~\ref{thm:FF}(1).
For this fixed $m$, keep $u$, $v$ and hence $\omega$ fixed while varying only the output indices.

Consider the $m\times m$ matrix $M$ with rows indexed by $c_\alpha=A_n(1)$ and columns by $c_\beta=B_{n'}(1)$
($n,n'=1,\dots,m$), all other output indices being fixed at $c_\gamma=C_\gamma(1)$, and entries $M_{nn'}:=\Omega_\omega(c)$.
By Lemma~\ref{lem:C}, applied with the output clusters $A_n\in\cC_\alpha$, $B_{n'}\in\cC_\beta$ and $C_\gamma\in\cC_\gamma$,
$M_{nn'}$ is the cluster value of the tuple with $A_n$ in slot $\alpha$, $B_{n'}$ in slot $\beta$, and $C_\gamma$ in
every other slot. The interlacing order gives $A_n<B_{n'}$ exactly when $n\le n'$. The tuple therefore has pattern
$\tau$ on and above the matrix diagonal, and pattern $\tau'$ below it.

Put $\delta:=\chi(\tau)-\chi(\tau')$, let $\mathbf 1\in L^m$ be the all-ones column, and let $U$ be the upper
unitriangular matrix with $U_{nn'}=1$ for $n\le n'$ and $U_{nn'}=0$ otherwise. Then
\[
M=\begin{pmatrix}
\chi(\tau)&\chi(\tau)&\cdots&\chi(\tau)\\
\chi(\tau')&\chi(\tau)&\cdots&\chi(\tau)\\
\vdots&\ddots&\ddots&\vdots\\
\chi(\tau')&\cdots&\chi(\tau')&\chi(\tau)
\end{pmatrix}
=\chi(\tau')\,\mathbf 1\mathbf 1^{\mathsf T}+\delta\,U,
\]
The rows vary only slot $\alpha$; each column fixes all the other slots, including the chosen value in slot $\beta$.
Thus $M$ is a submatrix of the flattening of this one array $\Omega_\omega$ along slot $\alpha$. Lemma~\ref{lem:flat}
gives $\rank M\le\sdim(\omega)\le d$.

If $\delta\ne0$, the matrix $\delta U$ is invertible, whereas the constant matrix
$\chi(\tau')\mathbf 1\mathbf 1^{\mathsf T}$ has rank at most one. Applying rank subadditivity to
$\delta U=M-\chi(\tau')\mathbf 1\mathbf 1^{\mathsf T}$ gives
\[
m=\rank(\delta U)\le\rank M+1\le d+1.
\]
The construction is available for every $m\ge1$; choosing $m=d+2$ is a contradiction. Hence $\delta=0$, as required.
\end{proof}

The support bound has now forced equality between adjacent order patterns. The remaining identity comes from
\eqref{eq:Pol1}; it does not use the support bound.

\begin{lemma}[the diagonal]\label{lem:D}
$\sum_{\tau\in S_k}\chi(\tau)=1$.
\end{lemma}

\begin{proof}
Take clusters $C_1<\dots<C_k$, and put $u_j:=s_{C_j}$, $v_j:=w_{C_j}$ and
$c:=(C_1(1),\dots,C_k(1))$. Apply $\omega\mapsto\Omega_\omega(c)$ to both sides of~\eqref{eq:Pol1}, with the operator
vectors $u_1,\dots,u_k$ and the vector arguments $v_1,\dots,v_k$.

\emph{Right side.} If $\sigma\ne\id$, some slot $j$ has $\sigma(j)\ne j$. Disjoint supports then give
$u_{\sigma(j)}v_j=0$, so the corresponding wedge vanishes. Only $\sigma=\id$ survives. It gives $\Omega_{(u_1v_1)\wedge\dots\wedge(u_kv_k)}(c)=\det\bigl((u_bv_b)(C_a(1))\bigr)_{a,b}$. Since
$u_bv_b=e_{C_b(1)}-e_{C_b(3)}$ and $C_a(1)\ne C_b(3)$ for all $a,b$, this is the determinant of the identity matrix, so
the right side is $1$.

\emph{Left side.} Fix $\sigma\in S_k$ and write $u_\sigma:=(u_{\sigma(1)},\dots,u_{\sigma(k)})$ and
$v_\sigma:=(v_{\sigma(1)},\dots,v_{\sigma(k)})$. To identify the summand with a cluster value, we must permute both the
vector arguments and the output coordinates to match the operator arguments. Each permutation contributes the same
sign. More precisely, antisymmetry of $\Psi$ in the vector slots gives
$\Psi(u_\sigma;v)=\sgn(\sigma)\Psi(u_\sigma;v_\sigma)$. For every $\omega\in\Lambda^kV$ the function $\Omega_\omega$ is
alternating in its arguments, so $\Omega_\omega(c)=\sgn(\sigma)\,\Omega_\omega(c_{\sigma(1)},\dots,c_{\sigma(k)})$; we
apply this with $\omega=\Psi(u_\sigma;v_\sigma)$. Since $\sgn(\sigma)^2=1$,
\begin{align*}
\Omega_{\Psi(u_\sigma;\,v)}(c)&=\sgn(\sigma)\,\Omega_{\Psi(u_\sigma;\,v_\sigma)}(c)
=\sgn(\sigma)^2\,\Omega_{\Psi(u_\sigma;\,v_\sigma)}(c_{\sigma(1)},\dots,c_{\sigma(k)})\\
&=\Omega_{\Psi(u_\sigma;\,v_\sigma)}(c_{\sigma(1)},\dots,c_{\sigma(k)}).
\end{align*}
In the last expression slot $j$ has operator $s_{C_{\sigma(j)}}$, vector $w_{C_{\sigma(j)}}$ and output
$C_{\sigma(j)}(1)$, so it equals $\chi(C_{\sigma(1)},\dots,C_{\sigma(k)})$ by definition. Its order pattern is $\sigma$,
since $C_{\sigma(j)}$ is the $\sigma(j)$\nobreakdash-th cluster in increasing order. Thus the convention for patterns
gives $\sigma$, not $\sigma^{-1}$. No cluster cancellation is needed in this calculation.

As $\sigma$ runs over $S_k$, each order pattern occurs exactly once. Comparing the two sides gives
\[
\sum_{\sigma\in S_k}\Omega_{\Psi(u_\sigma;\,v)}(c)=\sum_{\tau\in S_k}\chi(\tau)=1 .\qedhere
\]
\end{proof}

\begin{proof}[Proof of Theorem~\ref{thm:FF}]
Exchanging adjacent ranks acts on an order pattern by
$\tau\mapsto(\ell\ \ell{+}1)\circ\tau$, $1\le\ell<k$. These transpositions generate $S_k$, and Lemma~\ref{lem:S}
makes $\chi$ invariant under each exchange. Hence $\chi$ is constant; write its value as $\chi_0$.
Lemma~\ref{lem:D} now gives $k!\,\chi_0=1$ in $L$, contradicting $k!=0$. This proves (1), and the preceding reduction
gives (2).
\end{proof}

\begin{remark}[where the hypotheses enter]\label{rem:FFhyp}
Ramsey homogeneity is the only step in the obstruction proof that uses the finiteness of $L$.
Multilinearity then gives the cluster expansion and cancellation. The support bound is used in Lemma~\ref{lem:S}
only at the finitely many test inputs obtained by taking $m=d+2$; all have entries in $\{0,\pm1\}$.
Antisymmetry and~\eqref{eq:Pol1} enter in Lemma~\ref{lem:D}, which does not use that bound. Finally, $k!=0$ is used only
to turn the diagonal identity into a contradiction.

The inclusion--exclusion argument in Proposition~\ref{prop:diag}(1) uses only multilinearity in the
operator slots. The theorem therefore remains true with the pointwise identity \eqref{eq:Pw} at $L$\nobreakdash-points in place of
\eqref{eq:Pol1}, for this family of multipliers. In characteristic $2$ antisymmetry is the same as symmetry, and the two signs in
Lemma~\ref{lem:D} are trivial.
\end{remark}

\begin{remark}[when $k!$ is invertible]\label{rem:FFunit}
The vanishing of $k!$ is essential. If instead $k!\ne0$ in $L$, the \emph{normalized alternatization}
\[
\Psi_{\mathrm{alt}}(u;x):=\frac1{k!}\sum_{\sigma\in S_k}(u_{\sigma(1)}x_1)\wedge\dots\wedge(u_{\sigma(k)}x_k)
\]
is alternating in $x$ and satisfies~\eqref{eq:Pol1}. Every summand is supported in the span of the $k^2$ vectors
$u_ix_j$, so $\sdim(\Psi_{\mathrm{alt}}(u;x))\le k^2$. The preceding homogeneity, cancellation, staircase and diagonal
arguments apply, but now yield $\chi\equiv1/k!$: for pairwise disjoint clusters only $\sigma=\id$
survives in $\Psi_{\mathrm{alt}}(s_{C_1},\dots,s_{C_k};w_{C_1},\dots,w_{C_k})$, and the determinant at
$(C_1(1),\dots,C_k(1))$ is that of the identity. The final equation $k!\cdot(1/k!)=1$ is then consistent. So the
obstruction is exactly the vanishing of $k!$, not positive characteristic as such. For example, at $(p,k)=(5,3)$ every
step is consistent, since $3!=6=1$ in $\F_5$ and $\chi_0=1$.
\end{remark}

\section{The Banach exterior target over a Laurent series field}\label{app:laurent}\label{sec:laurent}

We now realize the finite-field obstruction in a Banach-space setting. The target is the completion of an exterior
power with its ordinary projective norm. Its determinant arrays admit a constant-coefficient map into the
\emph{algebraic} exterior power over the coefficient field, with support dimension controlled by the Banach norm.
This estimate turns a hypothetical bounded lift into the uniformly bounded-support map ruled out by
Theorem~\ref{thm:FF}. The passage to the completion is the key point: constant coefficients stabilize exactly once
the approximation error is less than one.

\begin{theorem}\label{thm:laurent}
Let $\kappa$ be a finite field of characteristic $p$, let $k\ge p$ and $r\in(0,1)$, and let $K_1=\kappa((X))$ with
$\abs X=r$. Let $E_1=\linf(\N,K_1)$, and let $B$ be the completion of $\Lambda^k_{K_1}E_1$ under the projective exterior
norm. Then $A^{k,K_1}_{E_1,E_1;B}$ admits no bounded $k$\nobreakdash-linear lift over $K_1$. Consequently it is analytic at no point of
$\cL(E_1,E_1)$.
\end{theorem}

\subsection{Setting}\label{ssec:laurentsetting}
Throughout this section $k\ge1$, $\kappa$ is a field, $r\in(0,1)$, and $K_1:=\kappa((X))$ is the Laurent series field with
absolute value $\abs a:=r^{\ord a}$ ($\abs0=0$). This field is complete, nontrivially normed and nonarchimedean, and
$\abs c=1$ for every $c\in\kappa^\times$. Write $\coeff_n(a)\in\kappa$ for the $X^n$\nobreakdash-coefficient of $a\in K_1$.
For a set $S$ and $y\in\linf(S,K_1)$ we take the constant coefficient coordinatewise:
\[
\coeff_0(y):=\bigl(\coeff_0(y(s))\bigr)_{s\in S}\in\kappa^S.
\]
The resulting \emph{constant-coefficient map} $\coeff_0$ is $\kappa$\nobreakdash-linear but not $K_1$\nobreakdash-linear.
A nonzero Laurent series of absolute value less than one has order at least one. Consequently,
\begin{equation}\label{eq:small}
\norm y_\infty<1\ \Longrightarrow\ \coeff_0(y)=0,\qquad\text{hence}\qquad \norm{y-y'}_\infty<1\ \Longrightarrow\ \coeff_0(y)=\coeff_0(y').
\end{equation}
Put $E_1:=\linf(\N,K_1)$ with the sup norm, a nonarchimedean $K_1$\nobreakdash-Banach space, and let $E_0:=\kappa^\N\subseteq E_1$ be
the $\kappa$\nobreakdash-subspace of sequences with entries in $\kappa$. Every nonzero $z\in E_0$ has $\norm z_\infty=1$. For
$x\in E_1$ and $n\in\Z$ let $x_{[n]}:=(\coeff_n(x(i)))_{i\in\N}\in E_0$ be the $n$\nobreakdash-th Laurent component. For $x\ne0$
the set $\{\ord x(i):x(i)\ne0\}$ is bounded below, because $r^{\ord x(i)}\le\norm x_\infty$, so it has a minimum
$\nu(x)$; then $\norm x_\infty=r^{\nu(x)}$ and $x_{[n]}=0$ for $n<\nu(x)$. Finally put
\[
M_r:=\max_{l\in\N}\,(l+1)r^l .
\]
The maximum exists because $(l+1)r^l\to0$. This constant will control the number of Laurent components that can
contribute to a constant coefficient. For $r\le\frac12$ one has $(l+1)r^l\le(l+1)2^{-l}\le1$, so $M_r=1$.

\subsection{The projective exterior norm}\label{ssec:projnorm}
On $\Lambda:=\Lambda^k_{K_1}E_1$ put
\[
\norm\omega_\pi:=\inf\Bigl\{\sum_j\prod_{i=1}^k\norm{x_{j,i}}_\infty\ :\ \omega=\sum_jx_{j,1}\wedge\dots\wedge x_{j,k}\Bigr\},
\]
where the infimum is over finite decompositions and the cost is an ordinary sum. Let $\Omega^{K_1}\colon\Lambda\to\linf(\N^k,K_1)$ be the determinant
array of \S\ref{ssec:ext} with $L=K_1$, $S=\N$ and $V=E_1$, so that
$\Omega^{K_1}(x_1\wedge\dots\wedge x_k)(i_1,\dots,i_k)=\det(x_b(i_a))_{a,b}$; its values are bounded by the next lemma.

\needspace{9\baselineskip}
\begin{lemma}\label{lem:proj}
$\norm\cdot_\pi$ is a $K_1$\nobreakdash-norm on $\Lambda$ with $\norm{x_1\wedge\dots\wedge x_k}_\pi\le\prod_i\norm{x_i}_\infty$, and
$\Omega^{K_1}$ is $K_1$\nobreakdash-linear and injective with $\norm{\Omega^{K_1}(\omega)}_\infty\le\norm\omega_\pi$. Let $B$ be the completion of
$(\Lambda,\norm\cdot_\pi)$, a $K_1$\nobreakdash-Banach space, $J\colon B\to\linf(\N^k,K_1)$ the continuous extension of $\Omega^{K_1}$ (so
$\norm J\le1$), and $W_B\colon E_1^k\to B$, $W_B(x):=x_1\wedge\dots\wedge x_k$. Then $W_B\in\Alt^k_{K_1}(E_1;B)$,
$\norm{W_B}\le1$, and $J(W_B(x))=\Omega^{K_1}(x_1\wedge\dots\wedge x_k)$.
\end{lemma}

\begin{proof}
For a pure wedge, each entry of the determinant array is a signed sum of products
$x_1(i_{\sigma(1)})\cdots x_k(i_{\sigma(k)})$. Each product has absolute value at most
$\prod_b\norm{x_b}_\infty$, so the nonarchimedean inequality in $K_1$ gives the same bound for the determinant.
For any finite decomposition of $\omega$, it follows that
\[
\norm{\Omega^{K_1}(\omega)}_\infty
\le\max_j\prod_i\norm{x_{j,i}}_\infty
\le\sum_j\prod_i\norm{x_{j,i}}_\infty.
\]
Thus the array is bounded, and taking the infimum gives
$\norm{\Omega^{K_1}(\omega)}_\infty\le\norm\omega_\pi$.
The map $\Omega^{K_1}$ is linear and injective by Lemma~\ref{lem:detinj}, with $L=K_1$, $S=\N$ and $V=E_1$.
In particular, $\norm\omega_\pi>0$ whenever $\omega\ne0$.

The remaining norm laws follow directly from decompositions. Concatenation gives the triangle inequality.
Since $k\ge1$, a scalar can be absorbed into the first factor of each wedge; for $\lambda\ne0$ this gives
$\norm{\lambda\omega}_\pi\le\abs\lambda\norm\omega_\pi$, and applying the same inequality to $\lambda^{-1}$ gives
equality. The pure-wedge bound follows by using a single summand.

The array space $\linf(\N^k,K_1)$ is complete, so $\Omega^{K_1}$ extends to the contraction $J$ on $B$.
The wedge map into $B$ is $K_1$\nobreakdash-multilinear and vanishes whenever two arguments coincide. The pure-wedge
bound makes it continuous with norm at most one, and the asserted identity with $J$ holds on its values by
construction.
\end{proof}

For $k\ge2$, $B$ is not nonarchimedean (Proposition~\ref{prop:nonultra}).

\subsection{The support estimate}\label{ssec:suppest}
We next show that the constant coefficient of every completed determinant array is represented by a finite
algebraic exterior vector. The quantitative bound comes from counting the Laurent components needed as wedge
\emph{factors}; the number of wedge summands may be much larger.

Let $\Omega^\kappa\colon\Lambda^k_\kappa E_0\to\kappa^{\N^k}$ be the determinant array over $\kappa$; it is injective by
Lemma~\ref{lem:detinj} with $L=\kappa$ and $V=E_0$. For $y_1,\dots,y_k\in E_0$ we have
$\Omega^{K_1}(y_1\wedge\dots\wedge y_k)=\Omega^\kappa(y_1\wedge\dots\wedge y_k)$: the determinants are the same, with entries in
$\kappa\subseteq K_1$. Support dimensions of elements of $\Lambda^k_\kappa E_0$ are computed in $E_0$ over $\kappa$.

\begin{proposition}[support estimate]\label{prop:supp}
For every $b\in B$ there is a unique $\eta(b)\in\Lambda^k_\kappa E_0$ with $\Omega^\kappa(\eta(b))=\coeff_0(Jb)$. The map $\eta$ is
$\kappa$\nobreakdash-linear, and
\begin{equation}\label{eq:suppest}
\sdim\bigl(\eta(b)\bigr)\le kM_r\,\norm b_B\qquad\text{for all }b\in B .
\end{equation}
\end{proposition}

\begin{proof}
\emph{Step 1: one decomposable.} Let $x_1,\dots,x_k\in E_1$. If some $x_b=0$, the wedge is zero and we take $\eta=0$.
Otherwise put $\nu_b:=\nu(x_b)$ and $l:=-\sum_b\nu_b$, so that
$\prod_b\norm{x_b}_\infty=r^{-l}$. The possible exponent tuples contributing to a constant coefficient form the set
\[
\mathcal N:=\Bigl\{n\in\Z^k:\ n_b\ge\nu_b\text{ for all }b,\ \sum_bn_b=0\Bigr\}.
\]
For such a tuple, the nonnegative integers $n_b-\nu_b$ sum to $l$. Hence $\mathcal N$ is empty if $l<0$;
if $l\ge0$, every coordinate satisfies $\nu_b\le n_b\le\nu_b+l$, so $\mathcal N$ is finite.

Fix $(i_1,\dots,i_k)\in\N^k$ and expand each entry as
$x_b(i_a)=\sum_{n\ge\nu_b}\coeff_n(x_b(i_a))X^n$. In the determinant expansion
\[
\det(x_b(i_a))_{a,b}=\sum_\sigma\sgn(\sigma)\prod_bx_b(i_{\sigma(b)}),
\]
the constant coefficient of each product is its finite coefficient convolution over $\mathcal N$.
Exchanging this finite sum with the permutation sum gives
\[
\coeff_0\det\bigl(x_b(i_a)\bigr)_{a,b}
=\sum_{n\in\mathcal N}\det\bigl(x_{b,[n_b]}(i_a)\bigr)_{a,b}.
\]
The same finite set $\mathcal N$ works for every array coordinate. Thus
$\coeff_0\bigl(\Omega^{K_1}(x_1\wedge\dots\wedge x_k)\bigr)=\Omega^\kappa(\eta)$, where
\[
\eta:=\sum_{n\in\mathcal N}x_{1,[n_1]}\wedge\dots\wedge x_{k,[n_k]}\in\Lambda^k_\kappa E_0.
\]
This is a finite sum in the algebraic exterior power. When $l<0$ it is zero. When $l\ge0$, all its wedge factors
belong to the set
$\{x_{b,[n]}:1\le b\le k,\ \nu_b\le n\le\nu_b+l\}$, which has at most $k(l+1)$ elements. Taking their span gives
\[
\sdim(\eta)\le k(l+1)=k\,(l+1)r^l\cdot r^{-l}\le kM_r\prod_b\norm{x_b}_\infty.
\]
The final bound therefore holds in both cases.

\emph{Step 2: an element of $\Lambda$.} Write $\omega=\sum_jw_j$, where
$w_j=x_{j,1}\wedge\dots\wedge x_{j,k}$, and take $\eta_j$ from Step~1. Additivity gives
\[
\coeff_0(\Omega^{K_1}(\omega))
=\sum_j\coeff_0(\Omega^{K_1}(w_j))
=\Omega^\kappa\Bigl(\sum_j\eta_j\Bigr).
\]
Injectivity of $\Omega^\kappa$ shows that $\eta(\omega):=\sum_j\eta_j$ is independent of the chosen decomposition.
Subadditivity of support dimension now yields
\[
\sdim(\eta(\omega))\le\sum_j\sdim(\eta_j)
\le kM_r\sum_j\prod_i\norm{x_{j,i}}_\infty.
\]
The left side is fixed while the decomposition varies. Taking the infimum therefore gives
$\sdim(\eta(\omega))\le kM_r\norm\omega_\pi$, without requiring a minimizing decomposition.

\emph{Step 3: completion.} Choose $\omega_n\in\Lambda$ with $\norm{b-\omega_n}_B\to0$. Since $J$ is a contraction,
\[
\norm{Jb-\Omega^{K_1}(\omega_n)}_\infty\le\norm{b-\omega_n}_B<1
\qquad(n\ge n_0).
\]
Equation~\eqref{eq:small} therefore gives the \emph{exact} identities
\[
\coeff_0(Jb)=\coeff_0(\Omega^{K_1}(\omega_n))=\Omega^\kappa(\eta(\omega_n))
\qquad(n\ge n_0).
\]
By injectivity of $\Omega^\kappa$, the algebraic vectors $\eta(\omega_n)$ are eventually equal. Define $\eta(b)$
to be this common value. Its array is $\coeff_0(Jb)$, so the same injectivity also proves independence of the
approximating sequence and uniqueness. In particular, no limit in $\Lambda^k_\kappa E_0$ is involved.

For every $n\ge n_0$ the fixed vector $\eta(b)$ satisfies
$\sdim(\eta(b))\le kM_r\norm{\omega_n}_\pi$. Since $\norm{\omega_n}_\pi\to\norm b_B$, this gives~\eqref{eq:suppest}.
Finally, the array characterization proves $\kappa$\nobreakdash-linearity: $J$ is $K_1$\nobreakdash-linear, $\coeff_0$ is
$\kappa$\nobreakdash-linear, and $\Omega^\kappa$ is linear and injective.
\end{proof}

\subsection{The coefficient lift}\label{ssec:Psi}
The support estimate converts norm control into an algebraic obstruction. We apply it to a hypothetical lift,
evaluated on coordinatewise multiplication operators and on the universal wedge map $W_B$.

Suppose that $P$ is a bounded $k$\nobreakdash-linear lift of $A^{k,K_1}_{E_1,E_1;B}$ over $K_1$. For $a\in E_1$ let
$D_a\in\cL(E_1,E_1)$ be coordinatewise multiplication, $D_ax=ax$. The map $a\mapsto D_a$ is $K_1$\nobreakdash-linear and
$\norm{D_a}\le\norm a_\infty$. Define
\begin{align*}
\Phi(a_1,\dots,a_k;x_1,\dots,x_k)&:=P(D_{a_1},\dots,D_{a_k})(W_B)(x_1,\dots,x_k)\in B,\\
\Psi&:=\eta\circ\Phi|_{E_0^{2k}}\colon E_0^k\times E_0^k\to\Lambda^k_\kappa E_0 .
\end{align*}
Since $\Omega^\kappa\circ\eta=\coeff_0\circ J$, the map $\Psi$ is $\coeff_0\circ J\circ\Phi|_{E_0^{2k}}$, read in $\Lambda^k_\kappa E_0$
through the injective map $\Omega^\kappa$.

\begin{lemma}\label{lem:Psi}
The map $\Psi$ has the following properties.
\begin{enumerate}
\item[($\Psi$1)] $\Psi$ is $2k$\nobreakdash-linear over $\kappa$.
\item[($\Psi$2)] $\Psi$ is alternating in the last $k$ slots.
\item[($\Psi$3)] $\sdim(\Psi(z))\le d:=\lfloor kM_r\norm P\rfloor$ for \emph{all} $z\in E_0^{2k}$.
\item[($\Psi$4)] $\Psi$ satisfies \eqref{eq:Pol} for the multiplier family $D^\kappa\colon E_0\to\End_\kappa(E_0)$,
$D^\kappa_ax=ax$: for all $m\ge1$, $b_1,\dots,b_m\in E_0$, $\alpha\in\N^m$ with $\abs\alpha=k$, and $x\in E_0^k$,
\[
\sum_{\type f=\alpha}\Psi(b_f;x)=\sum_{\type f=\alpha}(b_{f(1)}x_1)\wedge\dots\wedge(b_{f(k)}x_k)\qquad\text{in }\Lambda^k_\kappa E_0 .
\]
In particular $\Psi$ satisfies \eqref{eq:Pol1}.
\end{enumerate}
\end{lemma}

\begin{proof}
The map $\Phi$ is $K_1$\nobreakdash-multilinear, so its restriction to $E_0^{2k}$ is $\kappa$\nobreakdash-multilinear.
Composing with the $\kappa$\nobreakdash-linear map $\eta$ proves ($\Psi$1). Since
$P(D_{a_1},\dots,D_{a_k})(W_B)$ belongs to $\Alt^k_{K_1}(E_1;B)$, it vanishes when two of the last $k$ arguments
coincide; the same holds after applying $\eta$, proving ($\Psi$2).

For ($\Psi$3), all inputs from $E_0$ have norm at most one. Thus
\[
\norm{\Phi(a;x)}\le\norm P\prod_i\norm{D_{a_i}}\,\norm{W_B}\prod_i\norm{x_i}_\infty\le\norm P.
\]
Apply~\eqref{eq:suppest} and use that support dimension is an integer to obtain the stated bound by $d$.

To prove ($\Psi$4), we first compare coefficients over the infinite field $K_1$, while the values still lie in
$B$. Fix $b_1,\dots,b_m\in E_0$ and $x\in E_0^k$. For $\lambda\in K_1^m$ put
$a_\lambda:=\sum_i\lambda_ib_i\in E_1$. The diagonal identity for the lift gives
\[
\Phi(a_\lambda,\dots,a_\lambda;x)
=P(D_{a_\lambda},\dots,D_{a_\lambda})(W_B)(x)
=A(D_{a_\lambda})(W_B)(x)
=W_B(a_\lambda x_1,\dots,a_\lambda x_k).
\]
Expanding both sides by multilinearity gives
$\sum_\alpha\lambda^\alpha(L_\alpha-R_\alpha)=0$ for all $\lambda\in K_1^m$, where
\[
L_\alpha:=\sum_{\type f=\alpha}\Phi(b_f;x)\in B,\qquad
R_\alpha:=\sum_{\type f=\alpha}W_B(b_{f(1)}x_1,\dots,b_{f(k)}x_k)\in B.
\]
The field $K_1$ is infinite even when $\kappa$ is finite. Lemma~\ref{lem:idprinc}, applied in the
$K_1$\nobreakdash-vector space $B$, therefore yields $L_\alpha=R_\alpha$ for every $\alpha$.

We can now apply $\eta$. Additivity gives
$\eta(L_\alpha)=\sum_{\type f=\alpha}\Psi(b_f;x)$. To identify the other side, take $y_1,\dots,y_k\in E_0$.
By Lemma~\ref{lem:proj} and the determinant-array agreement preceding Proposition~\ref{prop:supp},
\[
J(W_B(y))=\Omega^{K_1}(y_1\wedge\dots\wedge y_k)
=\Omega^\kappa(y_1\wedge\dots\wedge y_k).
\]
This array has entries in $\kappa$, so taking its constant coefficient leaves it unchanged. Uniqueness in
Proposition~\ref{prop:supp} then gives
$\eta(W_B(y))=y_1\wedge\dots\wedge y_k\in\Lambda^k_\kappa E_0$.
Using $y_j=b_{f(j)}x_j$ in each summand of $R_\alpha$ proves
\[
\eta(R_\alpha)=\sum_{\type f=\alpha}(b_{f(1)}x_1)\wedge\dots\wedge(b_{f(k)}x_k),
\]
which is the required identity.
\end{proof}

The order of these operations is essential for the full identity ($\Psi$4). Evaluating the diagonal only at
$\kappa$\nobreakdash-points would give \eqref{eq:Pw} there, which for $\kappa=\F_q$ and $k>q$ is strictly weaker than
\eqref{eq:Pol} (Proposition~\ref{prop:diag}(3)). The proof below uses only the consequence~\eqref{eq:Pol1}.

\subsection{No bounded lift over a Laurent series field}

We now combine the support estimate (Proposition~\ref{prop:supp}) and the coefficient lift (Lemma~\ref{lem:Psi}) with the finite-field theorem of Appendix~\ref{app:FF} (Theorem~\ref{thm:FF}) to prove Theorem~\ref{thm:laurent}.

\begin{proof}[Proof of Theorem~\ref{thm:laurent}]
Suppose $P$ exists, and form $\Psi$ as in \S\ref{ssec:Psi}. By Lemma~\ref{lem:Psi},
$\Psi\colon E_0^k\times E_0^k\to\Lambda^k_\kappa E_0$ with $E_0=\kappa^\N$ is $2k$\nobreakdash-linear over $\kappa$, alternating and hence
antisymmetric in the last $k$ slots, satisfies \eqref{eq:Pol1}, and has $\sdim(\Psi(z))\le d$ for all inputs $z$. Moreover
$k!=0$ in $\kappa$, since $k\ge p=\operatorname{char}\kappa$. This contradicts Theorem~\ref{thm:FF}(2) with $L=\kappa$
and $E_0=\kappa^\N$. The last sentence follows from Proposition~\ref{prop:homogeneous-lift}.
\end{proof}

\begin{remark}\label{rem:infkappa}
Lemmas~\ref{lem:proj} and~\ref{lem:Psi} and Proposition~\ref{prop:supp} hold for every field $\kappa$, and nothing in them
uses a normalization of $r$: the constant $M_r$ makes the argument uniform in $r\in(0,1)$. Only the final appeal to
Theorem~\ref{thm:FF} needs $\kappa$ finite. We do not need, and do not discuss, the analogue of
Theorem~\ref{thm:laurent} for infinite $\kappa$ with this particular pair $(E_1,B)$. For an infinite $\kappa$ of
characteristic $p$ the field $\kappa((X))$ is covered by Theorem~\ref{thm:operator-obstructions}(1), with the pair obtained by ascent from its
subfield $\F_p((X))$.
\end{remark}

\subsection{Multiplier families and \texorpdfstring{$c_0$}{c0} bases}\label{lau:ssec:multipliers}
The proof of Theorem~\ref{thm:laurent} uses the hypothetical lift only on coordinatewise multiplication operators and
only at the single form $W_B$. Moreover, Theorem~\ref{thm:FF}(1) needs only finitely supported test vectors with
entries in $\{0,\pm1\}$. Both observations give a much smaller bad parameter space. Throughout this subsection
$\kappa$, $k$, $r$, $K_1$, $E_1$, $B$ and $W_B$ are as in Theorem~\ref{thm:laurent}, $c_0(\N,K_1)\subseteq E_1$ carries the
supremum norm, and $D_a\in\cL(E_1,E_1)$ denotes multiplication by $a$.

\begin{proposition}[multipliers on $c_0$]\label{lau:prop:c0-multipliers}
The homogeneous map
\[
 \sigma\colon c_0(\N,K_1)\to\Alt^k_{K_1}(E_1;B),\qquad
 \sigma(a)(x_1,\dots,x_k)=W_B(ax_1,\dots,ax_k)=(ax_1)\wedge\dots\wedge(ax_k),
\]
has no bounded $k$\nobreakdash-linear lift with values in $\Alt^k_{K_1}(E_1;B)$. Consequently, for every
$u_0\in\cL(E_1,E_1)$, the map $a\mapsto A^{k,K_1}_{E_1,E_1;B}(u_0+D_a)(W_B)$ is analytic at no point of $c_0(\N,K_1)$. In
particular $a\mapsto A^{k,K_1}_{E_1,E_1;B}(u_0+D_a)$ is analytic at no point.
\end{proposition}



\begin{proof}[Proof of Proposition~\ref{lau:prop:c0-multipliers}]
Suppose $p\colon c_0(\N,K_1)^k\to\Alt^k_{K_1}(E_1;B)$ is a bounded $k$\nobreakdash-linear lift of $\sigma$. Let
$V_{\mathrm{fin}}=\kappa^{(\N)}$ be the finitely supported sequences with entries in $\kappa$; it lies in both $c_0(\N,K_1)$
and $E_0$. Define $\Psi(a;x)=\eta\bigl(p(a_1,\dots,a_k)(x_1,\dots,x_k)\bigr)$ for $a,x\in V_{\mathrm{fin}}^k$. The proof of
Lemma~\ref{lem:Psi} applies verbatim. The map $\Psi$ is $2k$\nobreakdash-linear over $\kappa$ and alternating in the last
$k$ slots. Its inputs have supremum norm at most one, so~\eqref{eq:suppest} gives
$\sdim\Psi(a;x)\le\lfloor kM_r\norm p\rfloor$. The polarization points $a_\lambda=\sum_i\lambda_ib_i$, with
$b_i\in V_{\mathrm{fin}}$, are finitely supported and hence lie in $c_0(\N,K_1)$, so the identity~\eqref{eq:Pol}, and in
particular~\eqref{eq:Pol1}, follows from the diagonal identity as before. This contradicts Theorem~\ref{thm:FF}(1), since
$k!=0$ in $\kappa$.

For the second assertion, suppose that the map is analytic at $a_0$, with expansion $\sum_n p_n$ there. For
$h\in c_0(\N,K_1)$ and $t\in K_1$, multilinearity gives, in $\Mult^k_{K_1}(E_1;B)$,
\[
 W_B\circ(u_0+D_{a_0}+tD_h,\dots,u_0+D_{a_0}+tD_h)=\sum_{r=0}^kt^rC_r(h),\qquad C_k(h)=\sigma(h).
\]
Comparing this with the given expansion on scalar lines, Lemma~\ref{lem:unique} shows, exactly as in
Proposition~\ref{prop:homogeneous-lift}, that $p_k$ is a bounded lift of $\sigma$. This is impossible. The operator-valued
statement follows by evaluation at $W_B$.
\end{proof}

\begin{corollary}[diagonal transitions over a $c_0$ base]\label{lau:cor:diagonal-bundle}
Let $U$ be the open unit ball of $c_0(\N,K_1)$, an analytic manifold with one chart. Let $\mathcal V\to U$ be the
trivial bundle with fiber $E_1$ and the two global trivializations $\tau_0(a,z)=(a,z)$ and $\tau_1(a,z)=(a,D_{1+a}z)$,
and let $\mathcal B=U\times B$. Then $\mathcal V$ is an analytic normed vector bundle whose transitions are isometric
diagonal automorphisms of $E_1$. The naturally induced atlas on $\Alt^k(\mathcal V;\mathcal B)$ is not analytic: the section
that is constantly $W_B$ in the chart induced by $\tau_1$ is, in the chart induced by $\tau_0$, the map
$a\mapsto W_B\circ(D_{1+a},\dots,D_{1+a})$, which is analytic at no point of $U$.
\end{corollary}

\begin{proof}
For $a\in U$ every coordinate satisfies $\abs{1+a_i}=1$, so $D_{1+a}$ is an isometry of $E_1$. Its inverse
$D_{(1+a)^{-1}}=\sum_{n\ge0}(-D_a)^n$ is a power series of radius one in $a$, since $D_a^n=D_{a^n}$ has norm at most
$\norm a^n$. Hence both transitions are analytic. If $\omega_a\in\Alt^k(\mathcal V_a;B)$ has coordinates $m_0(a)$ and
$m_1(a)$ in the two induced charts, then $m_0(a)=m_1(a)\circ(D_{1+a},\dots,D_{1+a})$. With $m_1\equiv W_B$, this is the
map of Proposition~\ref{lau:prop:c0-multipliers} with $u_0=\id$, which is analytic at no point.
\end{proof}


The obstruction concerns the atlas obtained from the two analytic frames. The underlying pointwise bundle is trivial and has an analytic product structure; the failure is compatibility of that structure with the natural construction from all analytic frames. The complete target $B$ admits no equivalent nonarchimedean norm by Proposition~\ref{prop:nonultra}, so this example does not contradict the $c_0$ parameter theorem with nonarchimedean targets.

\begin{corollary}[finite dimension is sharp for nonarchimedean bases]\label{lau:cor:dichotomy}
Let $K=\F_q((u))$, where $q$ is a power of $p$, with a $u$\nobreakdash-adic absolute value, and let $k\ge p$. Let $P$ be a $K$\nobreakdash-Banach space admitting an equivalent nonarchimedean norm. The alternating construction preserves analytic bundles and their operator-valued analytic morphisms over every analytic manifold modeled on $P$, with arbitrary normed fibers, if and only if $P$ is finite-dimensional.
\end{corollary}

\begin{proof}
If $P$ is finite-dimensional, Lemma~\ref{lem:findim} supplies continuous coordinates and Corollary~\ref{cor:bundle-cases} applies. Conversely, suppose that $P$ is infinite-dimensional. Round an equivalent nonarchimedean norm upwards to the discrete value group $\abs{K^\times}$; the result is still an equivalent nonarchimedean $K$\nobreakdash-norm, and its nonzero values lie in $\abs{K^\times}$. Serre's orthonormal-basis theorem~\cite[\S1, Proposition~1]{Se} then identifies $P$ topologically with $c_0(I,K)$ for an infinite set $I$. Choosing a countable subset of $I$ gives bounded linear maps
\[
 \iota\colon c_0(\N,K)\longrightarrow P,\qquad
 \pi\colon P\longrightarrow c_0(\N,K),\qquad \pi\iota=\id.
\]
Apply the Laurent construction with $\kappa=\F_q$ and $K_1=K$. On the open set $V=\{x\in P:\norm{\pi x}<1\}$, take the trivial bundle with fiber $E_1$ and two analytic frames whose transition is $D_{1+\pi(x)}$, as in Corollary~\ref{lau:cor:diagonal-bundle}, and the trivial target bundle with fiber $B$. If the induced alternating transition were analytic at $x_0\in V$, restriction along $b\mapsto x_0+\iota b$ would make
\[
 b\longmapsto A^k_{E_1,E_1;B}\bigl(D_{1+\pi(x_0)}+D_b\bigr)
\]
analytic at zero. This contradicts Proposition~\ref{lau:prop:c0-multipliers} with $u_0=D_{1+\pi(x_0)}$.
\end{proof}

The dichotomy permits arbitrary normed target fibers. In the nonarchimedean Banach target category, Corollary~\ref{cor:discrete-targets} gives analyticity over every base for this discretely valued field.


\section{Ascent to an arbitrary base field}\label{app:ascent}\label{sec:ascent}

To pass from a Laurent series field to a prescribed field, we work first inside its completion. A suitable Laurent
subfield is spherically complete, so Ingleton's theorem supplies a contractive scalar retraction. The ordinary
projective tensor norm then gives isometric scalar extensions of the source and target, together with a retraction
onto the original Banach target. Operators and alternating forms extend compatibly with these embeddings. As a
result, any lift after scalar extension descends to the Laurent subfield, and nonexistence of a lift ascends.
The final density argument removes the completeness assumption on the prescribed field.

\subsection{The Laurent subfield}
We use spherical completeness in the sense used in Theorem~\ref{thm:spherical}: every nonempty family of pairwise-meeting
closed balls has a common point. For the field case of the following lemma, see
Schneider~\cite[Lemmas~1.3 and~1.6]{Sch}.

\begin{lemma}[discrete distances]\label{lem:discdist}
Let $Y$ be a complete nonarchimedean space whose distances lie in $r^\Z\cup\{0\}$ for some $r\in(0,1)$. Then $Y$ is
spherically complete.
\end{lemma}

\begin{proof}
Write $\cD(a,\rho)=\{y\in Y:\mathrm d(y,a)\le\rho\}$. In a nonarchimedean space, every point of a ball is a center.
Consequently, whenever two closed balls meet, the one with the smaller radius lies in the other.
The discrete distance assumption also gives
\[
\cD(a,\rho)=\cD(a,\rho^\flat),\qquad
\rho^\flat:=\max\{s\in r^\Z\cup\{0\}:s\le\rho\}.
\]
We may therefore replace every radius by a value in $r^\Z\cup\{0\}$.

Let $(\cD(a_i,\rho_i))_{i\in I}$ be a nonempty family of pairwise-meeting closed balls. If its radii have a least
element $\rho_{i_0}$, then $\cD(a_{i_0},\rho_{i_0})$ lies in every ball of the family, so $a_{i_0}$ is a common point.
Otherwise, discreteness implies that all the radii are positive and that their infimum is zero. Choose
$i_1,i_2,\dots$ so that $\rho_{i_n}$ decreases strictly to zero. The corresponding balls are nested, and
\[
\mathrm d(a_{i_{n'}},a_{i_n})\le\rho_{i_n}\qquad(n'\ge n).
\]
Their centers form a Cauchy sequence. By completeness it has a limit $a$, and closedness puts $a$ in every ball
of the sequence. For any $i\in I$, choose $n$ with $\rho_{i_n}\le\rho_i$. The $n$\nobreakdash-th ball then lies in
$\cD(a_i,\rho_i)$, so $a$ belongs to every ball of the original family.
\end{proof}



\begin{lemma}[Laurent subfield]\label{lem:laurentsub}
Let $K'$ be a complete nontrivially normed field of characteristic $p$, and let $t\in K'$ with $0<\abs t=r<1$.
\begin{enumerate}
\item $K'$ is nonarchimedean, and $\abs c=1$ for every $c\in\F_p^\times$.
\item For $a_m,\dots,a_M\in\F_p$, not all zero, $\bigl\lvert\sum_{i=m}^Ma_it^i\bigr\rvert=r^{\min\{i:a_i\ne0\}}$.
\item Give $\F_p((X))$ the absolute value with $\abs X=r$. There is a continuous isometric ring homomorphism
$\mathrm{ev}_t\colon\F_p((X))\to K'$ such that
$\mathrm{ev}_t\bigl(\sum_{i\ge m}a_iX^i\bigr)=\sum_{i\ge m}a_it^i$. Its image
$K_1$ is a closed subfield of $K'$; it is the closure of $\F_p(t)$, and it is isometrically isomorphic to $\F_p((X))$.
\item $K_1$ is spherically complete.
\end{enumerate}
\end{lemma}

\begin{proof}
(1) Lemma~\ref{lem:ultra} makes $K'$ nonarchimedean. For $c\in\F_p^\times$, the identity $c^{p-1}=1$ gives $\abs c=1$.

(2) The nonzero terms $a_it^i$ have pairwise distinct absolute values $r^i$. In a nonarchimedean field,
$\abs{x+y}=\abs x$ whenever $\abs y<\abs x$. Repeatedly applying this identity shows that the sum has the absolute
value of its largest term, namely $r^{\min\{i:a_i\ne0\}}$.

(3) Evaluation on Laurent polynomials defines a ring homomorphism
$\mathrm{ev}_t\colon\F_p[X,X^{-1}]\to K'$, and (2) shows that
$\abs{\mathrm{ev}_t(g)}=r^{\ord g}=\abs g$. Laurent polynomials are dense in $\F_p((X))$, since truncations converge.
Completeness of $K'$ therefore gives a unique continuous extension to $\F_p((X))$. Passing to limits preserves the
isometry and the ring operations. It also gives the stated evaluation formula, because the partial sums of
$\sum_{i\ge m}a_iX^i$ map to the partial sums of $\sum_{i\ge m}a_it^i$.

The image is complete as the isometric image of a complete space, hence closed in $K'$. It is a subfield: for
$g\ne0$, the identity $\mathrm{ev}_t(g)\mathrm{ev}_t(g^{-1})=1$ provides the inverse. This subfield contains
$\F_p(t)$, while the Laurent polynomials $\F_p[t,t^{-1}]\subseteq\F_p(t)$ are dense in it. The image is consequently
exactly the closure of $\F_p(t)$.

(4) By (3), $K_1$ is complete and has value group $\abs{K_1^\times}=r^\Z$. It is nonarchimedean by (1), so its distances
lie in $r^\Z\cup\{0\}$. Lemma~\ref{lem:discdist} gives spherical completeness.
\end{proof}

\subsection{Ingleton's theorem and its consequences}
Throughout the rest of this section we fix fields $K_1\subseteq K'$ with the following properties.
\begin{enumerate}
\item[(H1)] $K_1$, with the restricted absolute value, is nontrivially normed, complete and spherically complete.
\item[(H2)] $K'$ is nonarchimedean.
\end{enumerate}
By multiplication, $K'$ is a nonarchimedean normed $K_1$\nobreakdash-space. In our application $K'=\Kh$ and $K_1$ is the field of
Lemma~\ref{lem:laurentsub}.

\begin{theorem}[Ingleton~\cite{In}; see also~\cite{vR}]\label{thm:ingleton}
Let $K_1$ be a spherically complete nonarchimedean normed field, $V$ a $K_1$\nobreakdash-vector space with a nonarchimedean norm,
$V_0\subseteq V$ a subspace, $C\ge0$, and $T_0\colon V_0\to K_1$ linear with $\abs{T_0v}\le C\norm v$ for $v\in V_0$. Then
$T_0$ extends to a linear map $T\colon V\to K_1$ with $\abs{Tx}\le C\norm x$ for all $x\in V$.
\end{theorem}

Schneider~\cite[Prop.~9.2]{Sch} gives the more general formulation with a dominating nonarchimedean seminorm.

We will use Theorem~\ref{thm:ingleton} on the nonarchimedean $K_1$\nobreakdash-space $K'$ in two ways: to obtain the scalar
retraction below, and to extend finite-dimensional coordinate functionals in the proof of positivity of the
projective tensor norm. Both applications have codomain $K_1$. The spaces $E_1$ and $F_1$ introduced later need not
be nonarchimedean.

\begin{corollary}[Ingleton projection]\label{cor:proj}
There is a $K_1$\nobreakdash-linear map $\varpi\colon K'\to K_1$ with $\varpi|_{K_1}=\id$ and $\abs{\varpi(\lambda)}\le\abs\lambda$ for all
$\lambda\in K'$.
\end{corollary}

\begin{proof}
Theorem~\ref{thm:ingleton} with $V=K'$, $V_0=K_1$, $T_0=\id$ and $C=1$; its hypotheses hold by (H1) and (H2).
\end{proof}

\begin{samepage}
\begin{lemma}[extension]\label{lem:ext}
Let $\mathbb K\in\{K_1,K'\}$ and let $C\ge0$.
\begin{enumerate}
\item Let $V_0$ be a dense $\mathbb K$\nobreakdash-subspace of a normed $\mathbb K$\nobreakdash-space $V$, let $Y$ be a
$\mathbb K$\nobreakdash-Banach space and $T_0\colon V_0\to Y$ a $\mathbb K$\nobreakdash-linear map with $\norm{T_0x}\le C\norm x$. Then $T_0$ extends
uniquely to a continuous map $T\colon V\to Y$, which is $\mathbb K$\nobreakdash-linear with $\norm{Tx}\le C\norm x$.
\item Let $V_1,\dots,V_k$ be normed $K'$\nobreakdash-spaces with dense subspaces $V_i^0$, $Y$ a $K'$\nobreakdash-Banach space, and
$M\colon\prod_iV_i^0\to Y$ a $K'$\nobreakdash-multilinear map with $\norm{M(d)}\le C\prod_i\norm{d_i}$. Then $M$ extends uniquely to a
continuous $K'$\nobreakdash-multilinear map $\overline M\colon\prod_iV_i\to Y$ with $\norm{\overline M}\le C$.
\end{enumerate}
\end{lemma}
\end{samepage}

\begin{proof}
(1) The bound makes $T_0$ Lipschitz, hence uniformly continuous. Density of $V_0$ and completeness of $Y$ give a
unique continuous extension. Linearity and the norm bound pass to limits by continuity of addition, scalar
multiplication and the norm.

(2) The case $k=0$ is immediate. For $k\ge1$, multilinearity gives the telescoping estimate
\[
\norm{M(d)-M(d')}\le C\rho^{k-1}\sum_i\norm{d_i-d'_i}
\]
whenever $\norm{d_i},\norm{d'_i}\le\rho$ for every $i$, with $\rho>0$. For any $v\in\prod_iV_i$, choose tuples
$d^{(n)}\in\prod_iV_i^0$ converging to $v$. Their coordinates are uniformly bounded, so the estimate makes
$M(d^{(n)})$ Cauchy in $Y$. Define $\overline M(v)$ to be its limit. The same estimate, applied to two approximating
sequences, proves that this definition is independent of the choices and is continuous on bounded sets.
Multilinearity and the product norm bound pass to limits. Finally, density of $\prod_iV_i^0$ proves uniqueness.
\end{proof}

\subsection{Projective base change}
The next construction extends scalars while retaining the original norm and a contractive retraction whenever
the original space is complete. The scalar retraction from Corollary~\ref{cor:proj} is the essential ingredient;
the construction applies to arbitrary normed spaces.

For a normed $K_1$\nobreakdash-space $V$ let $V_{K'}:=V\otimes_{K_1}K'$, a $K'$\nobreakdash-vector space through the second factor, and put
\[
\norm u_\pi:=\inf\Bigl\{\sum_j\norm{x_j}\,\abs{\lambda_j}\ :\ u=\sum_jx_j\otimes\lambda_j\Bigr\},
\]
the infimum over finite decompositions, with an ordinary sum. For a bounded $K_1$\nobreakdash-linear $\psi\colon K'\to K_1$ let
$T_\psi\colon V_{K'}\to V$ be the $K_1$\nobreakdash-linear map with $T_\psi(x\otimes\lambda)=\psi(\lambda)x$; it exists because
$(x,\lambda)\mapsto\psi(\lambda)x$ is $K_1$\nobreakdash-bilinear. Let $\iota_V\colon V\to V_{K'}$, $x\mapsto x\otimes1$; it is
$K_1$\nobreakdash-linear, since $cx\otimes1=x\otimes c=c\cdot(x\otimes1)$ for $c\in K_1$.

\begin{lemma}[projective base change]\label{lem:bc}
Let $V$ be a normed $K_1$\nobreakdash-space; no nonarchimedean hypothesis is made on $V$.
\begin{enumerate}
\item $\norm{T_\psi u}\le\norm\psi\,\norm u_\pi$ for every bounded $K_1$\nobreakdash-linear $\psi\colon K'\to K_1$.
\item $\norm\cdot_\pi$ is a $K'$\nobreakdash-norm on $V_{K'}$.
\item $\iota_V$ is isometric.
\item Let $V\widehat\otimes_\pi K'$ be the completion of $(V_{K'},\norm\cdot_\pi)$. If $V$ is complete, then $T_\varpi$, with
$\varpi$ as in Corollary~\ref{cor:proj}, extends to a $K_1$\nobreakdash-linear map $\Pi_V\colon V\widehat\otimes_\pi K'\to V$ with
$\norm{\Pi_V}\le1$ and $\Pi_V\circ\iota_V=\id_V$.
\end{enumerate}
\end{lemma}

\begin{proof}
(1) For every decomposition $u=\sum_jx_j\otimes\lambda_j$,
\[
\norm{T_\psi u}=\Bigl\lVert\sum_j\psi(\lambda_j)x_j\Bigr\rVert
\le\sum_j\norm\psi\abs{\lambda_j}\norm{x_j}.
\]
Taking the infimum over decompositions proves the contraction estimate.

(2) Concatenating decompositions gives the triangle inequality. Multiplying their second factors by
$\mu\in K'$ gives $\norm{\mu u}_\pi\le\abs\mu\norm u_\pi$; for $\mu\ne0$, applying the same inequality with
$\mu^{-1}$ gives equality. The zero-scalar case is immediate.

It remains to prove positivity. Let $u\ne0$. Choose a finite tensor representation and rewrite its second
factors in a basis of their $K_1$\nobreakdash-span. This gives
\[
u=\sum_{j=1}^nx_j\otimes\lambda_j
\]
with $\lambda_1,\dots,\lambda_n$ linearly independent over $K_1$. Some $x_{j_0}$ is nonzero, since otherwise
$u=0$. On the finite-dimensional subspace $U:=\operatorname{span}_{K_1}\{\lambda_1,\dots,\lambda_n\}\subseteq K'$,
take the coordinate functional $\psi_0(\lambda_j)=\delta_{jj_0}$. It is bounded by some $C\ge0$ by
Lemma~\ref{lem:findim}, using completeness of $K_1$ in (H1). Apply Theorem~\ref{thm:ingleton} to the nonarchimedean
$K_1$\nobreakdash-space $K'$ to extend it to $\psi\colon K'\to K_1$ with $\norm\psi\le C$. Then
$T_\psi u=x_{j_0}$, so (1) gives
\[
0<\norm{x_{j_0}}\le C\norm u_\pi.
\]
Thus $\norm u_\pi>0$. This argument detects a nonzero vector coefficient by a functional on the \emph{scalar}
factor $K'$; it requires no separating continuous dual on $V$.

(3) The elementary tensor representation gives $\norm{x\otimes1}_\pi\le\norm x$. For the reverse inequality,
use $\varpi(1)=1$ and $\norm\varpi\le1$ in (1):
\[
\norm x=\norm{T_\varpi(x\otimes1)}\le\norm{x\otimes1}_\pi.
\]

(4) When $V$ is complete, the contraction $T_\varpi$ on the dense tensor subspace extends by
Lemma~\ref{lem:ext}(1), with $\mathbb K=K_1$, to a $K_1$\nobreakdash-linear contraction $\Pi_V$ with values in $V$.
On the embedded copy of $V$ it satisfies
$\Pi_V\iota_Vx=T_\varpi(x\otimes1)=\varpi(1)x=x$.
\end{proof}

We now fix the spaces to which the descent argument will be applied. Let $E_1$ be a normed $K_1$\nobreakdash-space
and $F_1$ a $K_1$\nobreakdash-Banach space, and put
\[
E:=E_1\widehat\otimes_\pi K',\qquad F:=F_1\widehat\otimes_\pi K',
\]
the completions of $((E_1)_{K'},\norm\cdot_\pi)$ and $((F_1)_{K'},\norm\cdot_\pi)$. They are $K'$\nobreakdash-Banach spaces;
$\iota_E\colon E_1\to E$ and $\iota_F\colon F_1\to F$ are $K_1$\nobreakdash-linear isometries, and $\Pi_F\colon F\to F_1$ is $K_1$\nobreakdash-linear
with $\norm{\Pi_F}\le1$ and $\Pi_F\circ\iota_F=\id_{F_1}$ (Lemma~\ref{lem:bc}).

\begin{lemma}[operators]\label{lem:ops}
For $f\in\cL_{K_1}(E_1,E_1)$ the map $f\otimes\id$ on $(E_1)_{K'}$ is $K'$\nobreakdash-linear with
$\norm{(f\otimes\id)u}_\pi\le\norm f\norm u_\pi$. Its extension $f_{K'}\in\cL_{K'}(E,E)$ satisfies $\norm{f_{K'}}\le\norm f$
and $f_{K'}\circ\iota_E=\iota_E\circ f$, and $f\mapsto f_{K'}$ is $K_1$\nobreakdash-linear.
\end{lemma}

\begin{proof}
For any decomposition $u=\sum_jx_j\otimes\lambda_j$, the image under $f\otimes\id$ has the decomposition
$\sum_jf(x_j)\otimes\lambda_j$, whose cost is at most
$\norm f\sum_j\norm{x_j}\abs{\lambda_j}$. Taking the infimum proves the norm bound on the algebraic tensor product.
Since $E$ is complete, Lemma~\ref{lem:ext}(1), with $\mathbb K=K'$, gives the continuous extension $f_{K'}$ with the
same bound. Its action on $x\otimes1$ proves the embedding identity. Additivity and $K_1$\nobreakdash-homogeneity in
$f$ hold on elementary tensors, hence on their span, and then on $E$ by continuity.
\end{proof}

\begin{lemma}[forms]\label{lem:forms-bc}
For $m\in\Alt^k_{K_1}(E_1;F_1)$ there is $m_{K'}\in\Alt^k_{K'}(E;F)$ with $\norm{m_{K'}}\le\norm m$ and
\[
m_{K'}(\iota_Ex_1,\dots,\iota_Ex_k)=\iota_F\bigl(m(x_1,\dots,x_k)\bigr)\qquad(x_i\in E_1),
\]
and $m\mapsto m_{K'}$ is $K_1$\nobreakdash-linear.
\end{lemma}

\begin{proof}
We first define the form on algebraic tensors, then establish a bound independent of their decompositions, and
finally pass to the completions. Fix a $K_1$\nobreakdash-basis $(\lambda_\beta)_{\beta\in\mathfrak B}$ of $K'$.
Every $u\in(E_1)_{K'}$ has a unique finite expansion
$u=\sum_\beta u^\beta\otimes\lambda_\beta$, with $u^\beta\in E_1$, and each coordinate map $u\mapsto u^\beta$ is
$K_1$\nobreakdash-linear. Define
\[
m^{\mathrm{alg}}(u_1,\dots,u_k):=
\sum_{\beta_1,\dots,\beta_k}
m\bigl(u_1^{\beta_1},\dots,u_k^{\beta_k}\bigr)\otimes
\lambda_{\beta_1}\cdots\lambda_{\beta_k}\in(F_1)_{K'}.
\]
The sum is finite, since each input has finitely many nonzero coordinates.

\textup{(i)} Linearity of the coordinate maps and multilinearity of $m$ make $m^{\mathrm{alg}}$
$K_1$\nobreakdash-multilinear.

\textup{(ii)} On elementary tensors, the definition gives
\[
m^{\mathrm{alg}}(x_1\otimes\mu_1,\dots,x_k\otimes\mu_k)
=m(x_1,\dots,x_k)\otimes\mu_1\cdots\mu_k
\qquad(x_i\in E_1,\ \mu_i\in K').
\]
Indeed, writing $\mu_i=\sum_\beta c_{i\beta}\lambda_\beta$ with $c_{i\beta}\in K_1$ gives
$x_i\otimes\mu_i=\sum_\beta(c_{i\beta}x_i)\otimes\lambda_\beta$. Substitution in the defining sum yields
\[
\sum_{\beta_1,\dots,\beta_k}
\Bigl(\prod_ic_{i\beta_i}\Bigr)m(x)\otimes\prod_i\lambda_{\beta_i}
=m(x)\otimes\prod_i\Bigl(\sum_\beta c_{i\beta}\lambda_\beta\Bigr),
\]
which is the asserted formula.

\textup{(iii)} The map $m^{\mathrm{alg}}$ is $K'$\nobreakdash-multilinear. By additivity in each slot, it suffices
to check $K'$\nobreakdash-homogeneity when every argument is an elementary tensor; there it follows from
\textup{(ii)}. The same formula determines the map on all tuples of algebraic tensors, so the initial basis
choice does not affect it.

\textup{(iv)} For arbitrary decompositions $u_i=\sum_jx_{ij}\otimes\mu_{ij}$, multilinearity and \textup{(ii)} give
\[
m^{\mathrm{alg}}(u)=
\sum_{j_1,\dots,j_k}m(x_{1j_1},\dots,x_{kj_k})\otimes\prod_i\mu_{ij_i}.
\]
The cost of this tensor decomposition is at most
$\norm m\prod_i\bigl(\sum_j\norm{x_{ij}}\abs{\mu_{ij}}\bigr)$.
Taking the infimum separately over the decompositions in each slot proves
\[
\norm{m^{\mathrm{alg}}(u_1,\dots,u_k)}_\pi\le\norm m\prod_i\norm{u_i}_\pi.
\]

\textup{(v)} To prove strong alternation, suppose two distinct slots $a,b$ both contain
$u=\sum_jx_j\otimes\mu_j$. Use this \emph{same} decomposition in both slots and expand all arguments into
elementary tensors. Terms using the same index $j$ in slots $a$ and $b$ vanish because $m$ has a repeated
argument. For $j\ne j'$, the term using indices $(j,j')$ in these slots pairs with the term using $(j',j)$.
Their scalar factors agree, while the values of $m$ are negatives of one another, so the pair sums to zero.
Thus $m^{\mathrm{alg}}$ vanishes whenever any two arguments coincide. Pair cancellation remains valid in
characteristic $2$; merely checking $m^{\mathrm{alg}}(u,\dots,u)=0$ would not establish strong alternation for
$k\ge3$.

\textup{(vi)} Regard $m^{\mathrm{alg}}$ as taking values in the completion $F$.
By \textup{(iii)}, \textup{(iv)} and Lemma~\ref{lem:ext}(2), it extends to a continuous
$K'$\nobreakdash-multilinear map $m_{K'}\colon E^k\to F$ with $\norm{m_{K'}}\le\norm m$.
Strong alternation also passes to the limit: if $v_a=v_b$, approximate this common vector by one sequence in
$(E_1)_{K'}$, and approximate the other entries separately. Every approximating tuple has equal entries in
slots $a,b$, so \textup{(v)} and continuity give $m_{K'}(v)=0$.

\textup{(vii)} Taking every scalar factor equal to one in \textup{(ii)} gives
\[
m_{K'}(\iota_Ex_1,\dots,\iota_Ex_k)
=m^{\mathrm{alg}}(x_1\otimes1,\dots,x_k\otimes1)
=m(x)\otimes1=\iota_F(m(x)).
\]
The defining formula is $K_1$\nobreakdash-linear in $m$. These linear identities pass to the continuous
extensions by uniqueness, proving $K_1$\nobreakdash-linearity of $m\mapsto m_{K'}$.
\end{proof}

\begin{proposition}[descent of a lift]\label{prop:ascent}
Assume \textup{(H1)} and \textup{(H2)}. Let $E_1$ be a normed $K_1$\nobreakdash-space, $F_1$ a $K_1$\nobreakdash-Banach space, and $E,F$ as
above. If $A^{k,K'}_{E,E;F}$ admits a bounded $k$\nobreakdash-linear lift $P$ over $K'$, then $A^{k,K_1}_{E_1,E_1;F_1}$ admits a
bounded $k$\nobreakdash-linear lift $P_1$ over $K_1$ with $\norm{P_1}\le\norm P$. Equivalently, by Proposition~\ref{prop:homogeneous-lift} over $K_1$ and over
$K'$: if $A^{k,K_1}_{E_1,E_1;F_1}$ is analytic at no point, then $A^{k,K'}_{E,E;F}$ is analytic at no point either.
\end{proposition}

\begin{proof}
Extend the operator and form inputs to $K'$, apply the given lift, and retract the value to $F_1$. Explicitly,
for $f_1,\dots,f_k\in\cL_{K_1}(E_1,E_1)$, $m\in\Alt^k_{K_1}(E_1;F_1)$ and $x\in E_1^k$, set
\[
P_1(f_1,\dots,f_k)(m)(x_1,\dots,x_k):=
\Pi_F\Bigl(P\bigl((f_1)_{K'},\dots,(f_k)_{K'}\bigr)(m_{K'})(\iota_Ex_1,\dots,\iota_Ex_k)\Bigr).
\]
We check the norm bound and then the diagonal identity.

For fixed $f$ and $m$, the map $P((f_i)_{K'})(m_{K'})$ is a continuous alternating $K'$\nobreakdash-multilinear form
on $E$. Precomposing each slot with the $K_1$\nobreakdash-linear isometry $\iota_E$ and postcomposing with
$\Pi_F$ gives a $K_1$\nobreakdash-multilinear map on $E_1$, still vanishing whenever two arguments coincide.
Lemmas~\ref{lem:bc}--\ref{lem:forms-bc} give
\[
\norm{P_1(f)(m)}
\le\norm{\Pi_F}\norm P\prod_i\norm{(f_i)_{K'}}\norm{m_{K'}}
\le\norm P\prod_i\norm{f_i}\norm m.
\]
In particular, $P_1(f)(m)\in\Alt^k_{K_1}(E_1;F_1)$.
The form extension is $K_1$\nobreakdash-linear in $m$, the operator extension is $K_1$\nobreakdash-linear in each
$f_i$, and $P$ is $K'$\nobreakdash-multilinear with values in $K'$\nobreakdash-linear operators.
Consequently $P_1$ is $K_1$\nobreakdash-linear in $m$ and $K_1$\nobreakdash-multilinear in $(f_1,\dots,f_k)$.
The displayed estimate proves boundedness at both levels and yields $\norm{P_1}\le\norm P$.

On the diagonal, use $P(g,\dots,g)=A^{k,K'}(g)$ with $g=f_{K'}$, followed by the embedding identities in
Lemmas~\ref{lem:ops} and~\ref{lem:forms-bc} and the retraction identity $\Pi_F\iota_F=\id$:
\begin{align*}
P_1(f,\dots,f)(m)(x)
&=\Pi_F\bigl(m_{K'}(f_{K'}\iota_Ex_1,\dots,f_{K'}\iota_Ex_k)\bigr)\\
&=\Pi_F\bigl(m_{K'}(\iota_Efx_1,\dots,\iota_Efx_k)\bigr)\\
&=\Pi_F\bigl(\iota_F(m(fx_1,\dots,fx_k))\bigr)\\
&=m(fx_1,\dots,fx_k)=A^{k,K_1}(f)(m)(x).\qedhere
\end{align*}
\end{proof}

\begin{remark}\label{rem:ascent}
Proposition~\ref{prop:ascent} is characteristic-free and uses no hypothesis on $k$. It transports negative results
upward only, from $K_1$ to $K'$. Ingleton's theorem is used only on $K'$, with no nonarchimedean norm hypothesis on $F_1$. The completeness of $F_1$ is used essentially: it is what makes the retraction $\Pi_F$ land in $F_1$.
\end{remark}

\subsection{Incomplete base fields}
The preceding construction takes place over the completion $\Kh$ of the prescribed field. To return to $K$,
we restrict scalars on the resulting Banach spaces. Density of $K$ in $\Kh$ forces every continuous
$K$\nobreakdash-multilinear map between these spaces to be $\Kh$\nobreakdash-multilinear, so a lift cannot appear under
this restriction.

\begin{lemma}\label{lem:descent}
Let $K$ be a nontrivially normed field with completion $\Kh$, and let $E,E',F$ be $\Kh$\nobreakdash-Banach spaces, regarded as
$K$\nobreakdash-Banach spaces by restriction of scalars (with the same norms). Then $\cL_K(E,E')=\cL_{\Kh}(E,E')$,
$\Alt^k_K(E;F)=\Alt^k_{\Kh}(E;F)$ and $\Alt^k_K(E';F)=\Alt^k_{\Kh}(E';F)$ with the same norms, $A^{k,K}_{E,E';F}=A^{k,\Kh}_{E,E';F}$,
and every bounded $k$\nobreakdash-linear lift of $A^{k,K}_{E,E';F}$ over $K$ is a bounded $k$\nobreakdash-linear lift of $A^{k,\Kh}_{E,E';F}$ over
$\Kh$.
\end{lemma}

\begin{proof}
Let $g$ be a continuous $K$\nobreakdash-linear map between normed spaces that already carry compatible
$\Kh$\nobreakdash-scalar actions. For $\lambda\in\Kh$, choose $\lambda_n\in K$ with $\lambda_n\to\lambda$.
Since $\norm{\lambda_nx-\lambda x}=\abs{\lambda_n-\lambda}\norm x$, continuity gives
\[
g(\lambda x)=\lim_n g(\lambda_nx)=\lim_n\lambda_ng(x)=\lambda g(x).
\]
Thus $g$ is $\Kh$\nobreakdash-linear. Applying the same argument with all but one slot fixed shows that every
continuous $K$\nobreakdash-multilinear map is $\Kh$\nobreakdash-multilinear. Alternation is the same equal-input
vanishing condition over either field.

The spaces of bounded maps therefore coincide as sets. Their operator norms agree as well, since the norms are
defined by the same pointwise bounds on the same vectors. They are $\Kh$\nobreakdash-Banach spaces, and
$A^{k,K}_{E,E';F}$ and $A^{k,\Kh}_{E,E';F}$ agree by their defining precomposition formula.
Finally, apply the slotwise density argument once more to a bounded $k$\nobreakdash-linear lift $P$ between these
map spaces. It becomes $\Kh$\nobreakdash-multilinear with the same norm and the same diagonal identity, and hence
is a lift over $\Kh$.
\end{proof}

\needspace{8\baselineskip}

\section{Proof of the counterexample theorem}\label{app:assembly}\label{sec:proofs}

Appendices~\ref{app:FF}--\ref{app:ascent} give the obstruction to a bounded lift over a Laurent series field and its transfer to a prescribed nontrivially normed field of characteristic $p$. We first establish the remaining property of the target, that it admits no equivalent nonarchimedean norm (Proposition~\ref{prop:nonultra}), and then assemble the proof of Theorem~\ref{thm:operator-obstructions}(1). The classification in Corollary~\ref{cor:classification} follows from this obstruction and the invertible-factorial construction.

\subsection{The target has no equivalent nonarchimedean norm}

\begin{proposition}\label{prop:nonultra}
Let $\kappa$ be any field, $k\ge2$, $r\in(0,1)$, and let $K_1=\kappa((X))$, $E_1$, $E_0$ and $B$ be as in
Appendix~\ref{sec:laurent}. Let $e_i\in E_0$ be the unit vectors and $\omega_j:=e_{kj+1}\wedge\dots\wedge e_{kj+k}\in B$
($j\ge0$), wedges on pairwise disjoint blocks of indices. Then $\norm{\omega_j}_B=1$ and
\[
\frac N{M_r}\ \le\ \Bigl\lVert\sum_{j<N}\omega_j\Bigr\rVert_B\ \le\ N\qquad(N\ge1).
\]
Hence the norm of $B$ is not nonarchimedean, and $B$ admits no equivalent nonarchimedean norm. The same holds for
$F=B\widehat\otimes_\pi K'$ as in Appendix~\ref{sec:ascent}, for any $K'\supseteq K_1$ satisfying \textup{(H1)} and \textup{(H2)},
regarded as a normed space over $K'$ or over any subfield of $K'$.
\end{proposition}

\begin{proof}
Each $\omega_j$ is a wedge of unit vectors, so $\norm{\omega_j}_B\le1$ by the defining projective norm. Evaluation
of its determinant array at the block indices gives $\det I=1$, hence
$\norm{\omega_j}_B\ge\norm{J\omega_j}_\infty\ge1$. The triangle inequality now gives the upper bound $N$.

For the lower bound, write $\xi_N$ for the sum of the same block wedges in $\Lambda^k_\kappa E_0$, and put
$\omega:=\sum_{j<N}\omega_j\in B$. Its array $J\omega$ is the constant Laurent array $\Omega^\kappa(\xi_N)$.
Thus $\coeff_0(J\omega)=\Omega^\kappa(\xi_N)$, and uniqueness in Proposition~\ref{prop:supp} gives
$\eta(\omega)=\xi_N$.

We apply Lemma~\ref{lem:contr} to $\xi_N$. Fix one of the first $N$ blocks and an index $i_0$ in it. Take for
$\varphi$ the $k-1$ coordinate functionals corresponding to the other indices of that block. On every other
block all these functionals vanish, so its contraction is zero; here $k\ge2$ ensures that the minors have
positive size. On the chosen block, the coefficient of $e_{i_0}$ is $\pm1$, while every other cofactor contains
the zero column obtained by evaluating the functionals at $e_{i_0}$. Consequently
$c_\varphi(\xi_N)=\pm e_{i_0}$. These contractions span all $kN$ unit vectors in the first $N$ blocks, so
\[
kN\le\sdim(\xi_N)=\sdim(\eta(\omega))\le kM_r\norm\omega_B
\]
by~\eqref{eq:suppest}. This proves the lower bound.

If $\norm\cdot_u$ were an equivalent nonarchimedean norm, choose positive constants $c,c'$ with
$\norm\cdot_B\le c\norm\cdot_u\le c'\norm\cdot_B$. Then
\[
\frac{N}{M_r}\le\Bigl\lVert\sum_{j<N}\omega_j\Bigr\rVert_B
\le c\max_{j<N}\norm{\omega_j}_u\le c'
\]
for every $N\ge1$, a contradiction. In particular, the original norm is not nonarchimedean. This argument uses
only the nonarchimedean inequality for $\norm\cdot_u$, not its scalar homogeneity.

Finally, $\iota_F\colon B\to F$ is a $K_1$\nobreakdash-linear isometry by Lemma~\ref{lem:bc}(3). The images
$\iota_F\omega_j$ satisfy the same unit-norm and sum-growth estimates, so the argument applies to $F$ as well.
Restriction to a subfield leaves these estimates unchanged.
\end{proof}

If $r\le\frac12$, then $M_r=1$ and $\norm{\sum_{j<N}\omega_j}_B=N$ exactly.

\subsection{Proof of Theorem~\ref{thm:operator-obstructions}(1)}
\begin{proof}
Let $K$ be a nontrivially normed field of characteristic $p>0$, and let $k\ge p$.
\begin{enumerate}
\item \emph{Choose a Laurent subfield.} By Lemma~\ref{lem:ultra}, $K$ is nonarchimedean. Its completion $\Kh$ is
a complete, nontrivially normed, nonarchimedean field of characteristic $p$. Choose $t\in K$ with
$0<\abs t=r<1$, and let $K_1$ be the closure of $\F_p(t)$ in $\Kh$. Lemma~\ref{lem:laurentsub} identifies $K_1$
isometrically with $\F_p((X))$, with $\abs X=r$, by $\mathrm{ev}_t$. In particular, $K_1$ is complete and
spherically complete.
\item \emph{Obtain the obstruction over $K_1$.} Put $E_1:=\linf(\N,K_1)$, and let $B$ be the completion of
$\Lambda^k_{K_1}E_1$ under the projective exterior norm. Transport of scalars along $\mathrm{ev}_t$ preserves
the norms, exterior powers and completions in this construction. It also identifies the operator and
alternating-map spaces, their precomposition maps, and bounded lifts. Theorem~\ref{thm:laurent}, with
$\kappa=\F_p$, therefore shows that $A^{k,K_1}_{E_1,E_1;B}$ has no bounded $k$\nobreakdash-linear lift.
\item \emph{Ascend to $\Kh$.} The inclusion $K_1\subseteq\Kh$ satisfies (H1) and (H2), and $B$ is complete.
Proposition~\ref{prop:ascent}, with $K'=\Kh$, gives Banach spaces
\[
E:=E_1\widehat\otimes_\pi\Kh,\qquad F:=B\widehat\otimes_\pi\Kh
\]
for which $A^{k,\Kh}_{E,E;F}$ has no bounded $k$\nobreakdash-linear lift.
\item \emph{Descend to $K$.} Regard the same $E$ and $F$ as $K$\nobreakdash-Banach spaces. By
Lemma~\ref{lem:descent}, a bounded lift of $A^{k,K}_{E,E;F}$ would extend to a bounded lift of
$A^{k,\Kh}_{E,E;F}$, which is impossible. Proposition~\ref{prop:homogeneous-lift} applies over the possibly incomplete
field $K$: $A^{k,K}_{E,E;F}$ is analytic at no point. Reindexing the arguments by a bijection gives the same conclusion, with the same spaces $E,F$, for every finite index set with $k$ elements.
\item \emph{Exclude an equivalent nonarchimedean target norm.} Since $k\ge p\ge2$, Proposition~\ref{prop:nonultra}
applies to $B$ and $F$, after the scalar identification in step~(2). Its conclusion persists when $F$ is
regarded over $K$.\qedhere
\end{enumerate}
\end{proof}

\begin{remark}\label{rem:spaces}
If $K_1=\Kh$, as when $K=\F_p((u))$ and $t=u$, the completed tensor products identify isometrically with
$E_1$ and $B$: use Lemma~\ref{lem:bc}(3) and $x\otimes\lambda=\lambda x\otimes1$. Thus $E$ is nonarchimedean in
this case. For $K=\F_q((u))$, $q=p^n$, with its $u$\nobreakdash-adic absolute value, one can instead apply
Theorem~\ref{thm:laurent} directly with $\kappa=\F_q$ and $X=u$, obtaining $E=\linf(\N,K)$ and $F=B$.
The proof over a general base field does not require, or assert, that $E$ is nonarchimedean.
\end{remark}

\section{A counterexample with scalar values}\label{app:scalar}

We now prove Theorem~\ref{thm:operator-obstructions}(2), in the more explicit form below. The obstruction persists when the value space is the ground field. The construction combines a
finite-dimensional incompatibility with a family of limit conditions on a countable tree. The finite
lemma applies directly in every degree whose factorial vanishes; no passage from degree $p$ to a higher
degree is needed.

\begin{theorem}[scalar counterexample]\label{thm:scalar-counter}
Let $K$ be a complete, nontrivially normed field that is not spherically complete, and let $k\ge1$ satisfy
$k! = 0$ in $K$. There are nonarchimedean Banach spaces $E,E'$ for which
\[
 A^k_{E,E';K}\colon\cL(E,E')\longrightarrow
 \cL\bigl(\Alt^k(E';K),\Alt^k(E;K)\bigr)
\]
is analytic at no point. The spaces may be taken to be closed subspaces of
$\linf(\Lambda;K^{k+1})$ and $\linf(\Lambda;K^k)$, respectively, containing the corresponding $c_0$ spaces,
where $\Lambda$ is countable. Both spaces are nonseparable.
\end{theorem}

The factorial hypothesis says precisely that $K$ has characteristic $p>0$ and $k\ge p$. In particular
$k\ge2$, and $K$ is nonarchimedean by Lemma~\ref{lem:ultra}. We first establish the three ingredients
of the construction.

\subsection{The countable multilinear tail property}

For a set $I$, a family $(a_i)_{i\in I}$ is \emph{null} if, for every $\varepsilon>0$, only finitely many
indices satisfy $\abs{a_i}\ge\varepsilon$. Over a complete nonarchimedean field, every null scalar family
is unconditionally summable: its finite partial sums form a Cauchy net, since the norm of a difference is
at most the largest omitted term. Its sum has norm at most $\sup_i\abs{a_i}$, and such sums may be regrouped.

We derive the multilinear tail property from van der Put's description of the dual of countable $\linf$. The induction is independent of the characteristic.

\begin{lemma}[multilinear tails]\label{scalar:tail}
Let $K$ be complete, nonarchimedean and not spherically complete, let $I$ be countably infinite, and let
$d\ge1$. Every bounded $d$\nobreakdash-linear map $\lambda\colon\linf(I,K)^d\to K$ has an expansion
\begin{equation}\label{scalar:array}
 \lambda(x^1,\dots,x^d)
 =\sum_{(i_1,\dots,i_d)\in I^d}a_{i_1,\dots,i_d}x^1_{i_1}\cdots x^d_{i_d},
 \qquad a_{i_1,\dots,i_d}=\lambda(e_{i_1},\dots,e_{i_d}),
\end{equation}
where $a\in c_0(I^d,K)$. Consequently, for every $\varepsilon>0$, there is a finite $S\subseteq I$ such that
$\abs{\lambda(x^1,\dots,x^d)}\le\varepsilon$ whenever all $\norm{x^r}\le1$ and at least one $x^r$ vanishes
on $S$. In particular $\lambda(e_i,\dots,e_i)\to0$ outside finite subsets of $I$.
\end{lemma}

\begin{proof}
After reindexing, take $I=\N$. Van der Put's theorem~\cite[Th\'eor\`eme~4.1]{VdP}, applied to countably
many copies of $K$ (a countable index set satisfies its nonmeasurability hypothesis), gives
\[
 \varphi(x)=\sum_i x_i\varphi(e_i),\qquad (\varphi(e_i))_i\in c_0(\N,K)
 \quad\bigl(\varphi\in\cL(\linf,K)\bigr).
\]
This is the case $d=1$. Suppose the assertion about the array and expansion holds in smaller degrees.
Fixing a coordinate vector in any slot of $\lambda$ gives a form of degree $d-1$. Hence, for every
$\varepsilon>0$, each slice of the set
\[
 W=\{(i_1,\dots,i_d):\abs{a_{i_1,\dots,i_d}}\ge\varepsilon\}
\]
obtained by fixing one coordinate is finite.

If $W$ were infinite, we could choose $w^{(0)},w^{(1)},\ldots\in W$ with the coordinates
$w^{(n)}_r$ distinct as $n$ varies, for each $r$: a finite list of forbidden coordinates excludes only
finitely many elements of $W$. Define an isometric embedding $T_r\colon\linf(\N,K)\to\linf(\N,K)$
by putting $(T_rx)_{w^{(n)}_r}=x_n$ and zero elsewhere. The pulled-back form
$\mu=\lambda\circ(T_1,\dots,T_d)$ has array $b$ with $\abs{b_{n,\dots,n}}\ge\varepsilon$ for all $n$.
Its slices are null by induction. Thus
\[
 F_n=\{v\in\N^d:\abs{b_v}\ge\varepsilon,\ v_r=n\text{ for some }r\}
\]
is finite. Choose $n_1<n_2<\cdots$ so that $n_s$ exceeds every coordinate occurring in
$F_{n_1}\cup\cdots\cup F_{n_{s-1}}$. Every nonconstant tuple in $\{n_s:s\ge1\}^d$ then has coefficient
of norm $<\varepsilon$: its least coordinate puts a large tuple in the corresponding $F_{n_r}$,
whereas its greatest coordinate violates the choice of the subsequence.

Let $y$ be the indicator of $\{n_s:s\ge1\}$ and set $\psi(x)=\mu(x,y,\dots,y)$. The expansion of each
fixed-coordinate slice gives
\[
 \psi(e_{n_s})
 =b_{n_s,\dots,n_s}
   +\sum_{v'\in\{n_r:r\ge1\}^{d-1},\ v'\ne(n_s,\dots,n_s)}b_{n_s,v'}.
\]
The summands in the last sum form a null family and all have norm $<\varepsilon$. Their supremum is
also $<\varepsilon$: outside a finite set they have norm $<\varepsilon/2$, and the remaining finite
maximum is $<\varepsilon$. The nonarchimedean inequality therefore gives
$\abs{\psi(e_{n_s})}\ge\varepsilon$ for every $s$, contradicting the displayed dual formula.
Thus $a$ is null. Applying that formula in the first slot and the inductive expansion in the remaining
slots gives~\eqref{scalar:array}; the null full array permits regrouping the iterated sums.

Finally choose a finite $S$ containing all coordinates occurring in tuples with $\abs{a_{i_1,\dots,i_d}}
>\varepsilon$. If one input vanishes on $S$, all these terms vanish. Every remaining term has norm at
most $\varepsilon$ on the product of unit balls, so the same holds for the sum. Apply this conclusion to
coordinate vectors outside $S$ to obtain the last assertion.
\end{proof}

\subsection{Chains and a finite-dimensional incompatibility}

Let $L$ be a nonempty finite set and $\Lambda=L^{<\omega}$ the set of finite words in $L$, including the
empty word. Write $s\preceq t$ when $t$ extends $s$, and write $sj$ for the word obtained by appending
$j\in L$. A \emph{$j$-chain} is a set $T=\{s_0,s_1,\dots\}$ with $s_rj\preceq s_{r+1}$ for every $r$.
Let $\Gamma_j$ denote all such chains.

\begin{lemma}[chain gap]\label{scalar:gap}
Chains with different labels have at most one common word. If sets $a_j\subseteq\Lambda$, $j\in L$,
contain all but finitely many words of every chain in $\Gamma_j$, then $\bigcap_{j\in L}a_j\ne\varnothing$.
\end{lemma}

\begin{proof}
If distinct words $s,t$ belong to both a $j$-chain and an $h$-chain, take $s$ shorter than $t$.
Then $t$ extends both $sj$ and $sh$, forcing $j=h$.

For each word $r$ and label $j$, there is an extension $s\succeq r$ such that every extension of $sj$
belongs to $a_j$. Otherwise start with $s_0=r$ and choose successively an extension
$s_{m+1}\succeq s_mj$ outside $a_j$. This is a $j$-chain with infinitely many words outside $a_j$,
a contradiction. Enumerate $L=\{j_1,\dots,j_q\}$ and start at the empty word $r_0$. Choose $s_m\succeq
r_{m-1}$ with the asserted property for $j_m$, and put $r_m=s_mj_m$. The word $r_q$ extends every
$s_mj_m$ and hence belongs to every $a_j$.
\end{proof}

The next lemma is purely algebraic; the alternating maps in its statement need no topology.
Over a field $\kappa$, put $V=\kappa^{k+1}$ and $V'=\kappa^k$. Write $e_1,\dots,e_{k+1}$ and
$e'_1,\dots,e'_k$ for their bases, $\varepsilon^a$ for the coordinate forms on $V$, and $\delta'$ for the
determinant on $V'$. Set
\begin{align*}
 N&=\{\varepsilon^a:1\le a\le k+1\}
       \cup\{\varepsilon^a+\varepsilon^b:1\le a<b\le k+1\},\\
 \Sigma_\nu&=\ker\nu,\qquad C'_s=\{z:z_s=0\},\qquad H'_{cd}=\{z:z_c=z_d\},
\end{align*}
and consider the finite family of pairs
\begin{equation}\label{scalar:family}
 \mathcal F=\{(V,H'_{cd}):c<d\}
 \cup\{(\Sigma_\nu,C'_s):\nu\in N,\ 1\le s\le k\}.
\end{equation}
Let $g_0\colon V\to V'$ be the coordinate projection onto the first $k$ coordinates.

\begin{lemma}[finite fiber obstruction]\label{scalar:fibre}
Suppose $k! = 0$ in $\kappa$. There is no $k$\nobreakdash-linear map
\[
 \tau\colon\operatorname{Hom}_\kappa(V,V')^k\longrightarrow\Alt^k_\kappa(V;\kappa)
\]
satisfying both
\begin{enumerate}
\item $\tau(g_0,\dots,g_0)=\delta'\circ(g_0,\dots,g_0)$;
\item for every $(\Sigma,\Sigma')\in\mathcal F$, the restriction
$\tau(g^1,\dots,g^k)|_{\Sigma^k}$ is zero whenever $g^r(\Sigma)\subseteq\Sigma'$ for every $r$.
\end{enumerate}
\end{lemma}

\begin{proof}
Let $D$ be the determinant on $V$. Every alternating $k$\nobreakdash-form $\alpha$ has a unique expression
$\alpha=D(w_\alpha,\cdot,\dots,\cdot)$: its evaluations on the basis tuples omitting one vector give
\[
 (w_\alpha)_a=(-1)^{a-1}\alpha(e_1,\dots,\widehat e_a,\dots,e_{k+1}).
\]
For a nonzero $\nu$, this description also gives
\begin{equation}\label{scalar:restriction}
 \alpha|_{(\ker\nu)^k}=0\quad\Longleftrightarrow\quad\nu(w_\alpha)=0.
\end{equation}
Indeed choose $\zeta$ with $\nu(\zeta)=1$ and a basis $y_1,\dots,y_k$ of $\ker\nu$. Writing
$w_\alpha=\nu(w_\alpha)\zeta+z$ with $z\in\ker\nu$ gives
$D(w_\alpha,y_1,\dots,y_k)=\nu(w_\alpha)D(\zeta,y_1,\dots,y_k)$, whose last factor is nonzero.

Suppose $\tau$ exists. For a function $a\colon\{1,\dots,k\}\to\{1,\dots,k\}$ define
\[
 B_a(\varphi_1,\dots,\varphi_k)
 =\tau(\varphi_1\otimes e'_{a(1)},\dots,\varphi_k\otimes e'_{a(k)}),
 \qquad (\varphi\otimes e'_c)(x)=\varphi(x)e'_c.
\]
If $a$ omits $s$, every input map takes $V$ into $C'_s$. Applying condition (2) to
$(\ker\varepsilon^b,C'_s)$ for each $b$ and using~\eqref{scalar:restriction} makes every coordinate of
$w_{B_a(\varphi)}$ zero. Hence $B_a=0$ unless $a$ is a permutation.

For a permutation $a$ and distinct $c,d$, replace the target vector in slots $a^{-1}(c)$ and $a^{-1}(d)$
by $e'_c+e'_d$. All input maps then take $V$ into $H'_{cd}$, so their image under $\tau$ is zero.
Expansion in these two slots gives $B_a+B_{(c\ d)\circ a}$ and two terms with a missing target coordinate.
The latter vanish by the preceding paragraph. Thus, with $B=B_{\id}$,
\begin{equation}\label{scalar:target-sign}
 B_a=\sgn(a)B\qquad(a\in S_k).
\end{equation}

Define the $(k+1)$\nobreakdash-linear form
\[
 \beta(\varphi_1,\dots,\varphi_k;\chi)
   =\chi\bigl(w_{B(\varphi_1,\dots,\varphi_k)}\bigr).
\]
If the $s$th argument is $\nu\in N$, its rank-one input vanishes on $\Sigma_\nu$, and every other
rank-one input takes values in $C'_s$. Conditions (2) and~\eqref{scalar:restriction} give
$\beta(\varphi_1,\dots,\nu,\dots,\varphi_k;\nu)=0$. Use this identity successively for
$\varepsilon^a$, $\varepsilon^b$ and $\varepsilon^a+\varepsilon^b$ to obtain
\[
 \beta(\dots,\varepsilon^a,\dots;\varepsilon^b)
 =-\beta(\dots,\varepsilon^b,\dots;\varepsilon^a).
\]
Consequently
$t(\sigma)=\beta(\varepsilon^{\sigma(1)},\dots,\varepsilon^{\sigma(k)};
\varepsilon^{\sigma(k+1)})$ changes sign under every transposition of a slot with the last slot.
These transpositions generate $S_{k+1}$, so $t(\sigma)=\sgn(\sigma)t(\id)$.

Expand $g_0=\sum_{c=1}^k\varepsilon^c\otimes e'_c$ in condition (1). The missing-coordinate terms vanish,
and~\eqref{scalar:target-sign} gives
\[
 \varepsilon^1\wedge\cdots\wedge\varepsilon^k
 =\sum_{a\in S_k}\sgn(a)B(\varepsilon^{a(1)},\dots,\varepsilon^{a(k)}).
\]
Apply $\alpha\mapsto\varepsilon^{k+1}(w_\alpha)$. Its value on the left is $(-1)^k$; on the right it is
$\sum_{a\in S_k}\sgn(a)^2t(\id)=k!t(\id)=0$. This is a contradiction. All identities remain valid in
characteristic $2$: repeated arguments vanish by alternation, and no step divides by $2$ or by a factorial.
\end{proof}

We will need a quantitative consequence of this finite incompatibility. Give $K^{k+1}$ and $K^k$ their
maximum norms and use the same family~\eqref{scalar:family} over $K$.

\begin{lemma}[finite test certificate]\label{scalar:certificate}
Assume $k!=0$ in $K$, as in Theorem~\ref{thm:scalar-counter}. There are finite sets of tests $\mathcal T_j$, one for each $j=(\Sigma^j,\Sigma'^j)\in\mathcal F$, consisting of
tuples $(g^1,\dots,g^k;\xi_1,\dots,\xi_k)$ with $g^r\Sigma^j\subseteq\Sigma'^j$ and $\xi_r\in\Sigma^j$,
such that every $K$\nobreakdash-multilinear $\tau$ satisfying condition (1) of Lemma~\ref{scalar:fibre} obeys
\begin{equation}\label{scalar:threshold}
 \max_{j\in\mathcal F}\max_{(g;\xi)\in\mathcal T_j}
 \abs{\tau(g^1,\dots,g^k)(\xi_1,\dots,\xi_k)}\ge1.
\end{equation}
\end{lemma}

\begin{proof}
Write $\operatorname{char}K=p$. All spaces in~\eqref{scalar:family} are defined over $\F_p$.
For each $j$, choose an $\F_p$\nobreakdash-basis of the subspace $\Sigma^j$ defined over $\F_p$, and one of
\[
 \{g\in\operatorname{Hom}_{\F_p}(\F_p^{k+1},\F_p^k):g\Sigma^j\subseteq\Sigma'^j\}.
\]
Their scalar extensions are bases of the corresponding $K$\nobreakdash-spaces: row reduction of the
defining prime-field matrices has the same pivot positions over $K$. Take all $k$\nobreakdash-tuples from
each basis to form $\mathcal T_j$. Multilinearity in the maps and the vectors makes these tests equivalent
to condition (2).

Choose coordinates $\theta$ for the finite-dimensional space of possible $\tau$. Condition (1) is a
system $A\theta=b$, and the tests for condition (2) form a system $C\theta=0$, with all matrices and
constants over $\F_p$. Lemma~\ref{scalar:fibre} over $\F_p$ makes the combined system inconsistent.
Row reduction of its augmented matrix therefore gives row vectors $\lambda,\mu$ over $\F_p$ with
\[
 \lambda A+\mu C=0,\qquad \lambda b=1.
\]
For every $K$\nobreakdash-valued solution of $A\theta=b$ it follows that $\mu C\theta=-1$.
Every nonzero coefficient of $\mu$ has absolute value $1$, since $\F_p^\times$ is finite.
The nonarchimedean inequality now gives~\eqref{scalar:threshold}.
\end{proof}

\subsection{The Banach spaces defined by chain limits}

Set $L=\mathcal F$ and $\Lambda=L^{<\omega}$. Write $\Sigma^j\subseteq V=K^{k+1}$ and
$\Sigma'^j\subseteq V'=K^k$ for the pair indexed by $j\in L$. Define
\begin{equation}\label{scalar:spaces}
 \begin{split}
 E&=\{x\in\linf(\Lambda;V):
       \operatorname{dist}(x_i,\Sigma^j)\longrightarrow0
       \text{ along every }T\in\Gamma_j,\ j\in L\},\\
 E'&=\{z\in\linf(\Lambda;V'):
       \operatorname{dist}(z_i,\Sigma'^j)\longrightarrow0
       \text{ along every }T\in\Gamma_j,\ j\in L\}.
 \end{split}
\end{equation}
Limits along a chain are taken outside its finite subsets. We verify the properties of these full
solution spaces that will enter the proof.

Each subspace in the definition admits a projection $R$ with $\norm R\le1$ and $\norm{1-R}\le1$.
For $\ker\varepsilon^a$, delete coordinate $a$; for $\ker(\varepsilon^a+\varepsilon^b)$, replace coordinate
$a$ by the negative of coordinate $b$; for $V$, use the identity. On $V'$, delete coordinate $s$ for
$C'_s$, and replace coordinate $c$ by coordinate $d$ for $H'_{cd}$. These formulas are projections onto
the stated spaces. Their norm bounds follow directly from the maximum norm and the nonarchimedean
inequality. Moreover
\begin{equation}\label{scalar:distance}
 \operatorname{dist}(v,\Sigma)=\norm{(1-R)v}.
\end{equation}
Indeed $Rv\in\Sigma$ gives one inequality, and $(1-R)(v-y)=(1-R)v$ for every $y\in\Sigma$ gives the other.

For fixed $j$ and $T\in\Gamma_j$, the corresponding condition in~\eqref{scalar:spaces} is thus the inverse
image of $c_0(T;V)$ under the bounded linear map $x\mapsto((1-R_j)x_i)_{i\in T}$, and similarly for $V'$.
The space $c_0$ is closed in the sup norm: a uniform approximation within $\varepsilon/2$ by a null family
leaves only finitely many coordinates of norm at least $\varepsilon$. Each condition is therefore closed
and linear. Their intersection is a closed subspace of the Banach space $\linf(\Lambda;V)$, respectively
$\linf(\Lambda;V')$, so $E$ and $E'$ are nonarchimedean Banach spaces.

Both contain their corresponding $c_0$ spaces. They are also modules over $\linf(\Lambda,K)$ under
coordinatewise multiplication, with $\norm{tx}\le\norm t\norm x$: each residual family in
\eqref{scalar:distance} is multiplied by the same bounded scalar family and remains null.
Finally, if $T\in\Gamma_j$, every bounded $\Sigma^j$\nobreakdash-valued family supported on $T$ belongs
to $E$. For chains with label $j$ its residual is zero; for a chain with another label, its support meets
that chain in at most one word by Lemma~\ref{scalar:gap}. The identical assertion holds for $E'$ and
$\Sigma'^j$.

Since $k\ge2$, the label $j=(V,H'_{12})$ occurs and both its spaces contain a vector of norm $1$.
Fix a $j$-chain $T$. Multiplying either such vector by arbitrary bounded scalar families on $T$ gives
an isometric embedding of $\linf(\N,K)$ into $E$, respectively $E'$. The binary-valued sequences form
an uncountable set at pairwise distance $1$, so both spaces are nonseparable.

\subsection{The global contradiction}

\begin{proof}[Proof of Theorem~\ref{thm:scalar-counter}]
Use the spaces~\eqref{scalar:spaces}. Suppose that $A=A^k_{E,E';K}$ has a bounded $k$\nobreakdash-linear
lift $P$. For $i\in\Lambda$, let $\xi_{[i]}\in E$ denote the family equal to $\xi\in V$ at $i$ and zero
elsewhere. Single-coordinate families belong to $c_0\subseteq E$, and likewise in $E'$. For
$g\in\operatorname{Hom}_K(V,V')$ let $g_{[i]}\in\cL(E,E')$ send $x$ to the family equal to $g(x_i)$ at
$i$ and zero elsewhere. Its norm is at most $\norm g$. Set
\[
 \delta'_i(z_1,\dots,z_k)=\delta'(z_{1,i},\dots,z_{k,i})\in\Alt^k(E';K).
\]
The determinant estimate for the maximum norm gives $\norm{\delta'_i}\le1$. Define the fiber map
\begin{equation}\label{scalar:fibre-map}
 \pi_i(g^1,\dots,g^k)(\xi_1,\dots,\xi_k)
 =P(g^1_{[i]},\dots,g^k_{[i]})(\delta'_i)
       (\xi_{1[i]},\dots,\xi_{k[i]}).
\end{equation}
It is $k$\nobreakdash-linear in the maps and alternating in the vectors, and the lift identity gives
\begin{equation}\label{scalar:fibre-diagonal}
 \pi_i(g,\dots,g)=\delta'\circ(g,\dots,g)
 \qquad\text{for every }g.
\end{equation}
In particular every $\pi_i$ satisfies condition (1) of Lemma~\ref{scalar:fibre}.

Fix $j=(\Sigma,\Sigma')\in L$, a chain $T\in\Gamma_j$, maps $g^1,\dots,g^k$ with
$g^r\Sigma\subseteq\Sigma'$, and vectors $\xi_1,\dots,\xi_k\in\Sigma$. We claim that
\begin{equation}\label{scalar:chain-decay}
 \pi_i(g^1,\dots,g^k)(\xi_1,\dots,\xi_k)\longrightarrow0
 \qquad(i\in T).
\end{equation}
For $t\in\linf(\Lambda,K)$ define $u^r_t$ by
\[
 (u^r_t x)_i=\begin{cases}t_i g^r(x_i),&i\in T,\\0,&i\notin T.\end{cases}
\]
This is a bounded map from $E$ to $E'$. To see that its values lie in $E'$, let $R$ be the projection
onto $\Sigma$ above and decompose, on $T$,
\[
 t_i g^r(x_i)=t_i g^r(Rx_i)+t_i g^r((1-R)x_i).
\]
The first term is a bounded $\Sigma'$\nobreakdash-valued family supported on $T$, hence belongs to $E'$.
The second is null on $T$ by~\eqref{scalar:distance}, hence its zero extension belongs to $c_0\subseteq E'$.
The coordinate formula gives $\norm{u^r_t}\le\norm t\norm{g^r}$, so $t\mapsto u^r_t$ is bounded linear
with values in $\cL(E,E')$.

For $c\in\linf(\Lambda,K)$ put
\begin{equation}\label{scalar:chain-form}
 m_c(z_1,\dots,z_k)=\sum_{i\in T}c_i\delta'(z_{1,i},\dots,z_{k,i}).
\end{equation}
This sum converges. In fact, if $R'$ projects onto $\Sigma'$, write
$z_{r,i}=a_{r,i}+b_{r,i}$ with $a_{r,i}=R'z_{r,i}$ and $b_{r,i}=(1-R')z_{r,i}$. Both families are bounded
by $\norm{z_r}$, and $b_{r,i}\to0$ along $T$. Expanding the determinant, the term with only $a$'s is
zero because $\dim\Sigma'=k-1$. Every remaining term contains an error factor, giving
\[
 \abs{\delta'(z_{1,i},\dots,z_{k,i})}
 \le\max_{1\le r\le k}\left(\norm{b_{r,i}}\prod_{q\ne r}\norm{z_q}\right)
 \longrightarrow0.
\]
Multiplication by bounded $c$ preserves this limit. Completeness of $K$ now gives the unconditional
scalar sum in~\eqref{scalar:chain-form}. Linearity in each argument and in $c$ follows by taking limits
of finite partial sums; repeated arguments make every summand zero. Moreover the determinant estimate gives
\[
 \abs{m_c(z_1,\dots,z_k)}\le\norm c\prod_{r=1}^k\norm{z_r}.
\]
Thus $c\mapsto m_c$ is a bounded linear map into $\Alt^k(E';K)$, of norm at most $1$.

Let $x^r$ be the family equal to $\xi_r$ on $T$ and zero elsewhere; it belongs to $E$ by the
chain-generator inclusion already proved. Coordinatewise multiplication defines a bounded linear map
$\rho\mapsto\rho x^r$ from $\linf(\Lambda,K)$ to $E$. Consequently
\[
 \begin{split}
 \Phi(t^1,\dots,t^k,c,r^1,\dots,r^k)
   =P(u^1_{t^1},\dots,u^k_{t^k})(m_c)(r^1x^1,\dots,r^kx^k)
 \end{split}
\]
is a bounded $(2k+1)$\nobreakdash-linear scalar form on $\linf(\Lambda,K)$, with bound
\[
 \norm P\prod_{s=1}^k\bigl(\norm{g^s}\norm{\xi_s}\bigr)
       \prod_{s=1}^k\norm{t^s}\,\norm c\,\prod_{s=1}^k\norm{r^s}.
\]
The set $\Lambda$ is countably infinite, so Lemma~\ref{scalar:tail} gives
$\Phi(e_i,\dots,e_i)\to0$ outside finite sets. For $i\in T$ we have
$u^s_{e_i}=g^s_{[i]}$, $m_{e_i}=\delta'_i$ and $e_i x^s=\xi_{s[i]}$.
Equation~\eqref{scalar:fibre-map} therefore identifies this diagonal value with the expression in
\eqref{scalar:chain-decay}, proving the claim.

Choose the finite tests of Lemma~\ref{scalar:certificate}, and define
\[
 G_j(i)=\max_{(g;\xi)\in\mathcal T_j}
       \abs{\pi_i(g^1,\dots,g^k)(\xi_1,\dots,\xi_k)},
 \qquad a_j=\{i\in\Lambda:G_j(i)<1\}.
\]
For fixed $j$ and $T\in\Gamma_j$, equation~\eqref{scalar:chain-decay} applies to every test in the finite
set $\mathcal T_j$. Hence $G_j(i)\to0$ along $T$, and $T\setminus a_j$ is finite. Lemma~\ref{scalar:gap}
supplies an $i$ belonging to every $a_j$. But~\eqref{scalar:fibre-diagonal} and
Lemma~\ref{scalar:certificate} force $G_j(i)\ge1$ for some $j$, a contradiction.

There is therefore no bounded $k$\nobreakdash-linear lift of $A$. Proposition~\ref{prop:homogeneous-lift} applies in
this degree without dividing by $k!$: analyticity at any point would make the degree-$k$ Taylor coefficient
such a lift. Thus $A$ is analytic at no point, as claimed.
\end{proof}



\section{Universal analytic reflection and projection norms}\label{app:reflection}\label{fam:sec:tensors}
The parameter theorems for alternating bundles exploit the fixed degree of the outer operator
polynomial. Reflection of every Banach-valued analytic map through every closed linear inclusion is a
stronger property. We characterize that property and give an explicit example separating the two.

Throughout this appendix, $K$ is complete and nontrivially normed, and $P$ is Banach.
For $n\ge1$, let $T_n(P)$ be the separated completed projective tensor power representing bounded
$n$\nobreakdash-linear maps from $P^n$ to Banach spaces. We use the ordinary projective seminorm,
defined by infima of sums of products of norms, and take its separated completion before invoking
the universal property. Its canonical multilinear map has norm at most one.
Define the closed subspace
\[
 \Delta_n(P)=\overline{\spn\{x^{\otimes n}:x\in P\}}\subseteq T_n(P).
\]
A projection onto $\Delta_n(P)$ will mean a bounded linear map
$R_n\colon T_n(P)\to\Delta_n(P)$ whose restriction to $\Delta_n(P)$ is the identity.

\begin{proposition}[homogeneous reflection]\label{fam:prop:tensor-homogeneous}
For a fixed $n\ge1$, the following conditions are equivalent.
\begin{enumerate}
\item For every Banach space $Z$, closed subspace $W\subseteq Z$, and bounded $n$\nobreakdash-linear map
$B\colon P^n\to Z$ with diagonal values in $W$, there is a bounded $W$\nobreakdash-valued
$n$\nobreakdash-linear map with the same diagonal.
\item The subspace $\Delta_n(P)$ is the range of a bounded projection in $T_n(P)$.
\end{enumerate}
\end{proposition}

\begin{proof}
Suppose first that a projection $R_n$ exists. Linearize $B$ to a bounded map
$\widetilde B\colon T_n(P)\to Z$. The diagonal hypothesis and closedness of $W$ imply
$\widetilde B(\Delta_n(P))\subseteq W$.
Composing $\widetilde B|_{\Delta_n(P)}$ with $R_n$ and the canonical tensor map gives the required
multilinear representative, of norm at most $\norm B\norm{R_n}$.
Conversely, take $Z=T_n(P)$, $W=\Delta_n(P)$, and $B$ equal to the canonical tensor map.
Linearizing the representative supplied by the first condition gives a bounded map into $\Delta_n(P)$
fixing every pure power. It fixes their linear span and, by continuity, their closed span, so it is a projection.
\end{proof}

\begin{definition}\label{fam:def:universal-reflection}
The Banach parameter space $P$ has \emph{universal analytic reflection} if, for every Banach space $Z$, closed
subspace $W\subseteq Z$, every open $U\subseteq P$, and every map $f\colon U\to W$, analyticity of
the composite $U\to W\hookrightarrow Z$ implies analyticity of $f$.
\end{definition}

\begin{theorem}[analytic tensor criterion]\label{fam:thm:tensor-analytic}
A Banach space $P$ has universal analytic reflection if and only if there are projections
$R_n\colon T_n(P)\to\Delta_n(P)$ for all $n\ge1$ such that
\begin{equation}\label{fam:eq:exponential-projections}
 \limsup_{n\to\infty}\norm{R_n}^{1/n}<\infty.
\end{equation}
\end{theorem}

\begin{proof}
Assume first that the projections exist. Expand an ambient analytic map with values in $W$ at a point.
The scalar-line uniqueness argument in Theorem~\ref{thm:coefficient-descent} puts every homogeneous diagonal in $W$.
The construction in Proposition~\ref{fam:prop:tensor-homogeneous} replaces its coefficient $p_n$ by a
$W$\nobreakdash-valued coefficient of norm at most $\norm{p_n}\norm{R_n}$ and with the same diagonal.
The exponential bound preserves a positive radius of convergence. The degree-zero coefficient is
already in $W$, and the resulting expansion represents the given map.

Conversely, form the Banach spaces with supremum norms
\[
 Z=c_0\bigl(\{T_n(P)\}_{n\ge1}\bigr),\qquad
 W=c_0\bigl(\{\Delta_n(P)\}_{n\ge1}\bigr)\subseteq Z.
\]
Here $c_0$ denotes families tending to zero in their respective Banach coordinate spaces; $W$ is closed.
On the open unit ball of $P$, define
\[
 f(x)=(x^{\otimes n})_{n\ge1}.
\]
Since $\norm{x^{\otimes n}}\le\norm x^n$, this vector lies in $W$.
As a $Z$\nobreakdash-valued map it is analytic: its degree-$n$ coefficient is the canonical tensor map
placed in coordinate $n$, and has norm at most one. This gives an expansion at zero. To re-expand
at $a$ with $\norm a<1$, choose $r>0$ with $\norm a+r<1$. Expanding each tensor power in $a+h$
gives translated coefficients, and the sum of all term bounds is at most
$\sum_{n\ge1}(\norm a+r)^n<\infty$. Completeness of $Z$ supplies the coefficients and the
rearranged sum, proving analyticity throughout the unit ball. Universal reflection gives a $W$\nobreakdash-valued
expansion $\sum_mq_m(x,\ldots,x)$ near zero with $\norm{q_m}\le C\rho^{-m}$ for some $C,\rho>0$.
Compose this expansion with the projection $\pi_n\colon W\to\Delta_n(P)$ and compare on scalar lines
with $x\mapsto x^{\otimes n}$. Uniqueness of scalar power-series coefficients gives
\[
 \pi_nq_n(x,\ldots,x)=x^{\otimes n}\qquad(x\in P).
\]
Linearizing $\pi_nq_n$ therefore gives a projection $R_n$ onto $\Delta_n(P)$ with
$\norm{R_n}\le\norm{q_n}\le C\rho^{-n}$, proving~\eqref{fam:eq:exponential-projections}.
\end{proof}

\subsection{Projection bounds and the special role of fixed degree}
For $P=K^d$ with the maximum norm, the coordinate construction in
Theorem~\ref{thm:finite-parameters}, applied to the canonical tensor map, gives
$\norm{R_n}\le d^n$ for $d\ge1$. A bounded linear change of coordinates introduces only a fixed
distortion factor to the $n$\nobreakdash-th power, so it does not change the exponential-growth condition.

For the ordinary sum-norm space $P=\ell^1(I,K)$, the projective tensor power is isometrically
$\ell^1(I^n,K)$. Indeed, the coordinate tensor map has norm at most one, while expansion in coordinate
tensors gives the reverse projective-norm estimate; completion then gives the identification.
In these coordinates, $\Delta_n(P)$ consists exactly of arrays constant on each permutation orbit in $I^n$.
One inclusion follows because every pure power is constant on such orbits. For the other, finite-coordinate
polynomial interpolation puts each finite orbit sum in the span of pure powers, and finite sums of orbit
sums are dense in the space of orbit-constant $\ell^1$ arrays.

Choose one representative of each orbit, send its basis vector to the sum of all basis vectors in that
orbit, and send the other basis vectors in the orbit to zero. This gives a projection with norm at most $n!$.
Over a nonarchimedean field and for infinite $I$, this estimate is optimal even among projections that
do not preserve the individual orbit blocks.

\begin{proposition}\label{fam:prop:l1-optimal}
If $K$ is nonarchimedean and complete, and $I$ is infinite, then
\[
 \inf\{\norm R:R\colon T_n(\ell^1(I,K))\to\Delta_n(\ell^1(I,K))
                  \text{ is a bounded projection}\}=n!.
\]
Consequently $\ell^1(I,K)$ does not have universal analytic reflection.
\end{proposition}

\begin{proof}
The preceding construction gives the upper bound. For the lower bound, take an orbit $\mathcal O$
of a tuple with $n$ distinct indices; it has $n!$ elements. For each $w\in\mathcal O$, let $\lambda_w$
be the common coefficient on $\mathcal O$ of $R(e_w)$.
The orbit sum $\sum_{w\in\mathcal O}e_w$ belongs to $\Delta_n(P)$ and is fixed by $R$, hence
\[
 \sum_{w\in\mathcal O}\lambda_w=1.
\]
The nonarchimedean inequality implies that some $\abs{\lambda_w}\ge1$.
For this $w$, the ordinary sum norm satisfies
$\norm{R(e_w)}_1\ge n!\abs{\lambda_w}\ge n!$.
Thus $\norm R\ge n!$. Finally, $(n!)^{1/n}$ is unbounded, so
Theorem~\ref{fam:thm:tensor-analytic} excludes universal analytic reflection.
\end{proof}

\begin{corollary}[entire polynomial approximation of a nonanalytic smooth map]
\label{fam:cor:no-common-radius}
Let $K$ be a complete nonarchimedean nontrivially normed field and put
$P=\ell^1(\N,K)$, with the ordinary sum norm. There are Banach spaces
$W\subseteq Z$, with $W$ closed, and a map $f$ on the open unit ball of $P$
such that $f$ is analytic into $Z$ and $C^\infty$ into $W$, but is not analytic
into $W$ at zero. Moreover, $f$ is a uniform limit on every smaller ball of
entire $W$\nobreakdash-valued polynomials.
\end{corollary}

\begin{proof}
Take
\[
 Z=c_0\bigl(\{T_n(P)\}_{n\ge1}\bigr),\qquad
 W=c_0\bigl(\{\Delta_n(P)\}_{n\ge1}\bigr),\qquad
 f(x)=(x^{\otimes n})_{n\ge1}.
\]
The spaces are Banach and $W$ is closed in $Z$. As in the proof of
Theorem~\ref{fam:thm:tensor-analytic}, the ambient expansion has coefficient
norms at most one and defines an analytic map on the open unit ball.
Smooth reflection through the closed isometric inclusion makes $f$
$C^\infty$ as a map into $W$.

Suppose a $W$\nobreakdash-valued expansion at zero has coefficients $q_n$.
Uniqueness on scalar lines gives
\[
 \pi_nq_n(x,\ldots,x)=x^{\otimes n},
\]
where $\pi_n\colon W\to\Delta_n(P)$ is the coordinate projection.
Linearizing $\pi_nq_n$ gives a projection
$T_n(P)\to\Delta_n(P)$ of norm at most $\norm{q_n}$.
Proposition~\ref{fam:prop:l1-optimal} therefore gives
$\norm{q_n}\ge n!$. Since $(n!)^{1/n}$ is unbounded, these coefficients
give a formal multilinear series of radius zero, contradicting the assumed analytic expansion.

For each $n$, choose the sorted-orbit projection
$R_n\colon T_n(P)\to\Delta_n(P)$. Its composition with the canonical
tensor map is a bounded multilinear representative of $x\mapsto x^{\otimes n}$
with values in $\Delta_n(P)$. Consequently the truncation
\[
 f_N(x)=(x,x^{\otimes2},\ldots,x^{\otimes N},0,\ldots)
\]
is an entire $W$\nobreakdash-valued polynomial. For $0<r<1$ and $\norm x\le r$,
\[
 \norm{f(x)-f_N(x)}_W\le r^{N+1},
\]
which proves the asserted uniform convergence.
\end{proof}

This example is characteristic-free within the nonarchimedean setting.
Every individual homogeneous term has an analytic lift, but their best
multilinear bounds grow factorially. The fixed-degree argument for alternating
families avoids this obstruction by using one outer degree throughout.

There is no conflict with Theorem~\ref{thm:summable-parameters}(2): every alternating-map family on these
parameter spaces is still analytic as an operator family when its final output fiber is complete. The ambient alternating action
has a fixed polynomial degree. Factoring a general analytic parameter map through a free $\ell^1$ space
therefore incurs a single factorial bound, rather than a sequence of factorial bounds across Taylor degrees.
Over a complete nonarchimedean field, the alternating-family property is consequently strictly weaker
than universal analytic reflection. Over $\mathbb R$ and $\mathbb C$ every Banach space has universal analytic reflection, by Theorem~\ref{fam:thm:tensor-analytic}: the symmetrization $\frac1{n!}\sum_{\sigma\in S_n}\sigma$ has norm at most one on $T_n(P)$ and, by polarization, is a projection onto $\Delta_n(P)$.

For comparison, in the category of nonarchimedean normed spaces with the \emph{nonarchimedean} projective
tensor norm, the corresponding tensor power of $c_0(I,K)$ is $c_0(I^n,K)$ and the same sorted-orbit
projection has norm at most one. This uses the infimum of maxima in the tensor seminorm, rather than the
ordinary projective seminorm used in Theorem~\ref{fam:thm:tensor-analytic}.
It is the tensor counterpart of the norm-preserving coefficient construction in
Theorem~\ref{thm:summable-parameters}(1); the two tensor norms must not be interchanged.

Theorem~\ref{fam:thm:tensor-analytic} thus gives a complete answer for universal Banach-valued analytic
reflection. It supplies a sufficient parameter criterion for alternating families. Indeed, if the final
fiber is complete, the ambient multilinear operator space is Banach, and universal reflection through the
closed inclusion of alternating operators makes every alternating-map family analytic as an operator family.
Over a complete nonarchimedean field, Proposition~\ref{fam:prop:l1-optimal} shows that the criterion is not a necessary one.
An intrinsic characterization of the parameter spaces for which every alternating-map family is analytic
as an operator family
remains a separate question.


\section{Why a largest full analytic domain need not exist}\label{app:no-largest}\label{dom:section}\label{sec:domains}
The positive bundle theorems restrict either the fibers or the parameter spaces. Another possible
restriction keeps arbitrary analytic bases but chooses a full subcategory of fiber pairs on whose
hom spaces the alternating construction is analytic. There need not be a largest such subcategory:
endomorphism actions at two objects can each be analytic while their cross-action is not.
As in Proposition~\ref{prop:homogeneous-lift}, a bounded $d$-linear \emph{lift} of a
homogeneous map $f$ means a bounded multilinear map $B$ satisfying $B(x,\ldots,x)=f(x)$.

\begin{definition}\label{dom:definition}
A full subcategory $\mathcal D\subseteq\mathsf{Vec}_K^{\mathrm{op}}\times\mathsf{Vec}_K$ is an
\emph{analytic domain} for $\Alt^k$ if the restriction of $\Alt^k$ to $\mathcal D$ is analytic on every
hom space.
\end{definition}

\Needspace{5\baselineskip}
\begin{definition}\label{dom:split-pair}
A pair $(E,F)$ is \emph{split in degree $k$} if the isometric inclusion
\[
j_{E,F}\colon\Alt^k(E;F)\longrightarrow\Mult^k(E;F)
\]
has a bounded linear retraction.
\end{definition}

\begin{lemma}\label{dom:split-incoming}
Every morphism-space action whose destination is a split pair is analytic. In particular, the full
subcategory of split pairs is an analytic domain.
\end{lemma}

\begin{proof}
Let $(E,F)$ be split, with retraction $R$, and let the other object be $(D,G)$. For $u\colon E\to D$
and $v\colon G\to F$, the action is
\[
m\longmapsto R\bigl(v\circ m\circ(u,\ldots,u)\bigr).
\]
Before applying $R$, replace the $k$ copies of $u$ by independent variables. This gives a bounded
$(k+1)$\nobreakdash-linear lift of the joint action, with values in
$\cL(\Alt^k(D;G),\Mult^k(E;F))$. Postcomposition with $R$ gives a lift with the required alternating
codomain. No compatibility between the retractions chosen for different objects is needed.
\end{proof}

\begin{theorem}[no largest full analytic domain]\label{dom:no-largest}
Let $p$ be prime, let $k\ge p$, and equip $K=\F_p(t)$ with its $t$\nobreakdash-adic absolute value.
There is no largest full analytic domain for
\[
\Alt^k\colon\mathsf{Vec}_K^{\mathrm{op}}\times\mathsf{Vec}_K\longrightarrow\mathsf{Vec}_K.
\]
The assertion remains true if the domains must be closed under isomorphisms, and if their objects are
restricted to pairs of finite-dimensional, nonarchimedean normed $K$\nobreakdash-spaces. In the latter case,
isomorphism closure is understood inside that restricted ambient category.

More precisely, there are objects $X_k=(E_k,G)$ and $Y_k=(D_k,G)$ whose singleton full subcategories
are analytic, such that $Y_k$ is split and the action on morphisms from $Y_k$ to $X_k$ is not analytic.
\end{theorem}

Here finite-dimensional means finite \emph{algebraic} dimension over the incomplete field $K$. These spaces
need not admit continuous coordinates over $K$. The theorem does not give a counterexample over a complete
field, or a counterexample with Banach objects.


We give the entire construction. Put $L=\F_p((t))$, the completion of $K$.
We first work in degree $p$, then extend to every $k\ge p$ by adjoining coordinate summands.

\subsection{The rigid source and the split pair}

We use the following elementary completion principle. Let $K$ be dense in a complete valued extension $L$,
let $(E_i)$ be a finite family of dense $K$\nobreakdash-subspaces of Banach $L$\nobreakdash-spaces $H_i$,
and let $Z$ be Banach over $L$.
A bounded $K$\nobreakdash-multilinear map $\prod_iE_i\to Z$ extends uniquely to a bounded
$L$\nobreakdash-multilinear map $\prod_iH_i\to Z$, with the same norm. Indeed, approximate each input by
vectors in $E_i$. A telescoping expansion and the multilinear bound make the resulting values Cauchy,
independently of the approximations. The bound persists in the limit, and density of $K$ in $L$ gives
$L$\nobreakdash-linearity in each slot. In the applications below the extended codomain is the complete
space $L$ or $A$, while membership of values on the original inputs in the smaller spaces is retained.

Put $L=\F_p((t))$, the completion of $K$, and consider the normed $L$\nobreakdash-algebra
\[
A=L[\epsilon]/(\epsilon^p),\qquad e_i=\epsilon^i\quad(0\le i<p),
\]
with the maximum norm in this basis. Multiplication satisfies $\norm{xy}\le\norm x\,\norm y$.
Choose elements $a_0,\ldots,a_{p-1}\in L$, together with the finitely many auxiliary scalars
$\tau_w$ introduced below, algebraically independent over $K$. Such a choice is possible because $K$ is
countable and $L$ is uncountable: the elements of $L$ algebraic over any finitely generated extension of
$K$ form a countable set. Set
\[
a=\sum_{i=0}^{p-1}a_i e_i,\qquad E=K^p+Ka\subseteq A.
\]
Every subspace used below carries the norm induced by $A$ or $L$. In particular, $E$ has $K$\nobreakdash-dimension
$p+1$ and completion $A$.

\begin{lemma}\label{dom:rigid}
Every bounded endomorphism of $E$ is multiplication by a scalar in $K$, and every bounded linear map
$E\to K$ is zero.
\end{lemma}

\begin{proof}
A bounded $K$\nobreakdash-linear map extends to completions, and the extension is $L$\nobreakdash-linear
by density of $K$ in $L$. If it preserves $E$, its columns, the images of the $e_i$, give a matrix
$M+ac^{\mathsf T}$, with $M\in M_p(K)$ and $c\in K^p$. Its value at $a$ lies in $K^p+Ka$.
Comparison of homogeneous degrees in the algebraically independent $a_i$ shows first that
$a(c^{\mathsf T}a)=0$, hence $c=0$, and then that $M=sI$ for some $s\in K$.

Similarly, the extension of a bounded map $E\to K$ is an $L$\nobreakdash-linear functional on $A$.
Its coefficients belong to $K$, since its values on the $e_i$ do. Its value on $a$ also lies in $K$,
so algebraic independence makes all its coefficients zero.
\end{proof}

Multiplication by $a$ has lower triangular Toeplitz matrix
\[
(M_a)_{ij}=\begin{cases}a_{i-j},&i\ge j,\\0,&i<j.\end{cases}
\]
Define
\[
D_0=\spn_K\{e_i,\epsilon^i a:0\le i<p\}+Ka^2,
\qquad
\mathcal W=\{e_i:0\le i<p\}\cup\{e_i+e_j:0\le i<j<p\},
\]
and put
\[
D=D_0+\sum_{w\in\mathcal W}K\tau_w w,\qquad
G=\spn_K\{\det(d_1,\ldots,d_p):d_i\in D\}\subseteq L.
\]
The determinant is taken in the basis $(e_0,\ldots,e_{p-1})$. These are finite-dimensional
$K$\nobreakdash-spaces, with completions $A$ and $L$: $D$ contains $K^p$, and $G$ contains $1$.
By multilinearity, $G$ is spanned by determinants of the displayed finite generating set for $D$.
We may therefore identify $G$ with a finite-dimensional subspace of
$K[a_0,\ldots,a_{p-1},(\tau_w)_w]$. Every element has degree at most one in each individual $\tau_w$,
since a nonzero determinant cannot repeat the generator $\tau_w w$. Algebraic independence gives
\begin{equation}\label{dom:tau-gap}
\tau_w^2G\cap G=\{0\}\qquad(w\in\mathcal W).
\end{equation}

\begin{lemma}\label{dom:all-alt}
Every bounded $p$\nobreakdash-linear map $D^p\to G$ is alternating. Thus $(D,G)$ is split in degree $p$.
\end{lemma}

\begin{proof}
Compose such a map $m$ with $G\hookrightarrow L$ and extend to a bounded $L$\nobreakdash-multilinear map
on $A^p$. Fix two argument positions. Put the same $w\in\mathcal W$ in those positions and arbitrary
standard basis vectors in the others, and call the value $c\in G$. Replacing both copies of $w$ by
$\tau_w w\in D$ shows that $\tau_w^2c\in G$, so~\eqref{dom:tau-gap} gives $c=0$.
Taking $w=e_i$ gives vanishing on repeated basis vectors. Taking $w=e_i+e_j$ then gives the signed
interchange identity in the two positions. Expanding two equal arbitrary arguments in the basis proves
that the completed map, and hence $m$, is alternating, including when $p=2$. The inclusion of alternating
maps into multilinear maps is therefore an isometric bijection, whose inverse is the desired retraction.
\end{proof}

\subsection{A missing coefficient}

The coefficient space of alternating maps on $E$ is
\[
C=\{\lambda\in G:a_i\lambda\in G\text{ for every }i\}.
\]
Indeed, an alternating map extends to $\lambda\det$ on the $p$\nobreakdash-dimensional completion $A$.
Determinants of tuples from $E$ span $\spn_K\{1,a_0,\ldots,a_{p-1}\}$: a determinant can contain at
most one copy of $a$, and replacing one basis vector by $a$ realizes each $a_i$ up to sign. Consequently
\begin{equation}\label{dom:coefficient-space}
\Alt^p(E;G)=C\det.
\end{equation}
This identification is isometric, since the determinant bound gives one inequality and evaluation on the
standard basis gives the reverse inequality.

Algebraically set every $\tau_w$ to zero. The image of $G$ is
\[
G_0=\spn_K\{\det(d_1,\ldots,d_p):d_i\in D_0\}.
\]
Put $C_0=\{\lambda\in G_0:a_i\lambda\in G_0\text{ for every }i\}$. Specialization sends $C$ into $C_0$;
no continuity of specialization is asserted. The displayed generators of $D_0$ are homogeneous polynomial
vectors of degrees zero, one, and two. Their determinants are homogeneous, so $G_0$ is graded, and the
defining conditions show that $C_0$ is graded as well.

\begin{lemma}\label{dom:degree-gap}
The homogeneous degree-one part of $C_0$ is zero.
\end{lemma}

\begin{proof}
The degree-two part of $G_0$ is spanned by the $2\times2$ minors of $M_a$ and the coordinates of $a^2$.
These correspond respectively to choosing two degree-one columns, or one degree-two column, in a
determinant of generators. On homogeneous quadratic polynomials define the $K$\nobreakdash-linear map
$\Phi\colon K[a_0,\ldots,a_{p-1}]_2\to K[T]$ by
\[
\Phi(a_i a_j)=
\begin{cases}
T^{i+j},&i+j\ge p-1,\\
0,&i+j<p-1.
\end{cases}
\]
Every $2\times2$ minor of $M_a$ is killed by $\Phi$. When both product terms are nonzero, their index
sums coincide. If only one is nonzero, write the rows as $r<s$ and columns as $i<j$. The surviving term
has index sum $r+s-i-j$, and the missing term implies $r<j$. Hence
$r+s-i-j\le s-i-1\le p-2$, so its image is again zero.

The coefficient of $\epsilon^r$ in $a^2$ is $\sum_{i+j=r}a_i a_j$. Its image is zero for $r<p-1$;
for $r=p-1$ there are exactly $p$ summands and its image is $pT^{p-1}=0$. Thus $\Phi$ kills the entire
degree-two part of $G_0$. If $\ell=\sum_jc_j a_j$ lies in the degree-one part of $C_0$, then
$a_{p-1}\ell\in G_0$, whereas
\[
\Phi(a_{p-1}\ell)=\sum_jc_jT^{p-1+j}.
\]
The powers are distinct, so all $c_j$ vanish.
\end{proof}

\subsection{The nonanalytic cross-action}

Set $X=(E,G)$ and $Y=(D,G)$. Lemma~\ref{dom:rigid} identifies $\cL(E,E)$ isometrically with $K$,
and the self-action of $X$ is
\[
(s,v)\longmapsto s^p v_*.
\]
This is a bounded polynomial, so the singleton full subcategory on $X$ is analytic. The singleton on $Y$
is analytic by Lemmas~\ref{dom:all-alt} and~\ref{dom:split-incoming}.

For $x\in E$, multiplication by $x$ in $A$ restricts to a bounded map $M_x\colon E\to D$.
Indeed, writing $x=b+ha$ with $b\in K^p$ and $h\in K$, its values on the generators $e_j,a$ lie in
\[
K^p+\sum_iK\epsilon^ia+Ka^2\subseteq D.
\]
Submultiplicativity gives a bounded linear family $x\mapsto M_x$ from $E$ to $\cL(E,D)$.
Its matrix is lower triangular with diagonal entry $x_0$, the $e_0$\nobreakdash-coordinate of $x$ in $A$, so $\det M_x=x_0^p$. In the cross-action from
$Y$ to $X$, fix the target morphism to $\id_G$, evaluate at $\det\in\Alt^p(D;G)$, and use
\eqref{dom:coefficient-space}. The resulting map is
\begin{equation}\label{dom:bad-coordinate}
f\colon E\longrightarrow C,\qquad f(x)=x_0^p.
\end{equation}
Its values lie in $C$ because it was obtained by pulling back an alternating $G$\nobreakdash-valued map.

Suppose that $f$ were analytic at zero. Since $f(sx)=s^pf(x)$, uniqueness of power-series coefficients
on scalar lines gives a bounded $p$\nobreakdash-linear map $B\colon E^p\to C$ with diagonal $f$.
This is the homogeneous coefficient argument used in Proposition~\ref{prop:homogeneous-lift}; it requires no factorial
division or completeness. Compose $B$ with $C\hookrightarrow L$ and extend to an $L$\nobreakdash-multilinear
map on $A^p$. Put
\[
b_i=B(e_0,\ldots,e_0,e_i)\in C.
\]
There are $p-1$ initial copies of $e_0$, and $b_0=f(e_0)=1$. Evaluation at $a$ in the last position gives
$\sum_i a_i b_i\in C$. Specializing all $\tau_w$ to zero yields elements $d_i\in C_0$ with
\[
d_0=1,\qquad \sum_i a_i d_i\in C_0.
\]
The degree-one part of this last expression is $\sum_i a_i(d_i)_0$, where $(d_i)_0\in K$ is the
constant term. Its $a_0$\nobreakdash-coefficient equals one. This contradicts Lemma~\ref{dom:degree-gap}.
Thus~\eqref{dom:bad-coordinate}, and therefore the cross-action from $Y$ to $X$, is not analytic.

\subsection{Higher degrees and the categorical conclusion}

For $k=p+r$, let
\[
E_k=E\oplus K^r,\qquad D_k=D\oplus K^r
\]
with maximum norms. Their completions have dimension $k$ over $L$. Determinants of tuples from $E_k$
still span $\spn_K\{1,a_0,\ldots,a_{p-1}\}$, so extraction of the top-degree coefficient identifies
\begin{equation}\label{dom:padded-coefficients}
\Alt^k(E_k;G)=C\det_k
\end{equation}
isometrically. In particular, the coefficient takes values in $C$, a fact needed to retain the preceding
obstruction after increasing the degree.

By Lemma~\ref{dom:rigid}, the extension of every bounded endomorphism of $E_k$ to its
$L$\nobreakdash-completion $A\oplus L^r$ has block form
\[
\begin{pmatrix}sI_p&B\\0&T\end{pmatrix},\qquad s\in K,\quad T\in M_r(K).
\]
The scalar $s$ and entries of $T$ are bounded linear functions of the endomorphism, as is seen by
evaluation on the standard basis. Its determinant is $s^p\det T$, so its action on
\eqref{dom:padded-coefficients} is multiplication by this bounded polynomial. Including postcomposition
shows that the singleton full subcategory on $X_k=(E_k,G)$ is analytic.

The pair $Y_k=(D_k,G)$ is split. Given $m\in\Mult^k(D_k;G)$, restrict the first $p$ arguments to
$D$ and fix the last $r$ arguments at the standard auxiliary basis $e'_1,\ldots,e'_r$ of $K^r$.
The resulting $p$-linear map $\alpha_m$ is alternating by Lemma~\ref{dom:all-alt}. Write
$\pi_D\colon D_k\to D$ and $\pi_r\colon D_k\to K^r$ for the two projections. Explicitly, set
\[
 (\mathcal Rm)(v_1,\ldots,v_k)=
 \sum_{\substack{S\subseteq\{1,\ldots,k\}\\|S|=p}}
 \varepsilon(S)\,
 \alpha_m(\pi_Dv_{s_1},\ldots,\pi_Dv_{s_p})
 \det(\pi_rv_{t_1},\ldots,\pi_rv_{t_r}),
\]
where $s_1<\cdots<s_p$ list $S$, $t_1<\cdots<t_r$ list its complement, and $\varepsilon(S)$
is the sign of the permutation $(s_1,\ldots,s_p,t_1,\ldots,t_r)$. The determinant and the sign
are scalars multiplying the $G$-valued factor. In degree $r=0$ the determinant is $1$.
This shuffle expression is alternating by the same repeated-argument cancellation as the wedge
product, involves no factorial division, and gives a bounded linear map
$\mathcal R\colon\Mult^k(D_k;G)\to\Alt^k(D_k;G)$ with norm at most $\binom{k}{p}$.
It fixes every alternating map: after extension to the $k$\nobreakdash-dimensional completion, both maps
are determined by their value on the standard basis, where the shuffle formula has exactly one nonzero
term. This is the required retraction. For $r>0$, we do not need or assert that every $k$\nobreakdash-linear
map on $D_k$ is alternating.

Finally, restrict the cross-action along the bounded affine family
\[
E\longrightarrow\cL(E_k,D_k),\qquad x\longmapsto M_x\oplus I_r.
\]
Evaluation at $\det_k$, which is $G$\nobreakdash-valued on $D_k$ by Laplace expansion along the first $p$ rows, and coefficient extraction through~\eqref{dom:padded-coefficients} give exactly
the nonanalytic map~\eqref{dom:bad-coordinate}. Thus the cross-action from $Y_k$ to $X_k$ is not analytic.

A largest analytic full subcategory would have to contain both individually analytic singleton
subcategories, and fullness would then require this cross-action to be analytic. This is impossible.
If closure under isomorphisms is required, replace each singleton by its full isomorphism closure.
Every action between two copies is transported from the self-action by bounded linear changes of
coordinates, so both closures remain analytic. The same contradiction proves Theorem~\ref{dom:no-largest}.


The scope of this obstruction is precise. Finite algebraic dimension over an incomplete field does
not ensure continuous coordinates. The same construction does not answer the largest-domain question
over complete fields or for Banach pairs: completion replaces the exotic finite-dimensional
$K$-spaces by ordinary finite-coordinate $L$-spaces and removes the missing-coefficient obstruction.

\paragraph{\textit{Finite-dimensional bases.}}
The same spaces show that Theorem~\ref{thm:finite-parameters} needs continuous coordinates on the base,
not only finite algebraic dimension. Let $k=p+r$, and regard $E$, of dimension $p+1$, as an analytic
manifold with one chart. It has no continuous coordinates, since every bounded linear map $E\to K$ is
zero by Lemma~\ref{dom:rigid}. Over $E$ take the trivial bundle with fiber $D_k\oplus E_k$ and the two
frames $\id$ and
\[
 S_x(d,e)=\bigl(d+(M_x\oplus I_r)e,e\bigr),\qquad S_x^{-1}(d,e)=\bigl(d-(M_x\oplus I_r)e,e\bigr),
\]
and the trivial bundle with fiber $G$. Both $S_x$ and $S_x^{-1}$ are affine in $x$, so these are analytic bundles with
finite-dimensional fibers. By~\eqref{eq:shear-compression}, analyticity of the induced alternating
transition at $x=0$ would make $x\mapsto A^k_{E_k,D_k;G}(M_x\oplus I_r)$ analytic at zero. Evaluation at
$\det_k$ and~\eqref{dom:padded-coefficients} would then make~\eqref{dom:bad-coordinate} analytic at zero,
which was excluded above. So the induced alternating atlas is not analytic.


\section{Further coordinate methods}\label{app:coordinate-methods}
\subsection{Countable-type scalar sources}
The obstruction requires large source spaces in a stronger sense than infinite dimension.
A normed space is \emph{of countable type} if a countable subset has dense linear span.

\begin{proposition}[countable-type scalar sources]\label{prop:scalar-countable}
Let $K$ be complete and nonarchimedean. If either $E$ or $D$ is of countable type, then
$A^k_{E,D;K}$ is analytic in every degree. Neither completeness nor a nonarchimedean norm
is required of $E,D$. In particular, the two spaces in
Theorem~\ref{thm:operator-obstructions}(2) must both fail to be of countable type.
\end{proposition}

\begin{proof}
The case $k=0$ is constant. For a normed space $V$, set
\[
 q_V(x)=\inf_{x=\sum_i x_i}\max_i\norm{x_i},
\]
where the decompositions use nonempty finite families, allowing a zero term. Scaling and concatenating them show that
$q_V$ is a nonarchimedean seminorm. Let $V^u$ be the completion of its separated quotient;
the canonical map $j_V\colon V\to V^u$ is contractive and has dense range.
Expanding a bounded scalar multilinear form $m$ along decompositions of each argument gives
\[
 \abs{m(x_1,\ldots,x_k)}\le\norm m\prod_a q_V(x_a).
\]
Thus scalar multilinear forms extend uniquely and isometrically to $V^u$, as do
alternating forms. Extension preserves alternation by density of tuples with repeated
entries, and commutes with permutations. Every bounded $u\colon V\to W$ induces
$u^u\colon V^u\to W^u$ with $\norm{u^u}\le\norm u$. Uniqueness on the dense image makes
$u\mapsto u^u$ bounded linear, and extension of forms intertwines precomposition by $u$ and $u^u$.

If $V$ is of countable type, so is $V^u$. By van der Put's basis theorem
\cite[Proposition~1.4]{VdP}, it has a $t$\nobreakdash-orthogonal basis $(b_i)$ for any
fixed $0<t<1$. Writing $f_i$ for its coordinate functionals, define
\[
 S(m)(x_1,\ldots,x_k)
 =\sum_{i_1<\cdots<i_k}m(b_{i_1},\ldots,b_{i_k})
                        \prod_{a=1}^k f_{i_a}(x_a).
\]
The weighted coordinates $\abs{f_i(x)}\norm{b_i}$ form a null family bounded by
$t^{-1}\norm x$. Products of finitely many bounded null families are null, so the
displayed sum converges unconditionally in $K$ and $\norm{S(m)}\le t^{-k}\norm m$.
Let $\mathcal A_V$ be unnormalized antisymmetrization, the signed sum over permutations
of the arguments. It is bounded and alternating-valued in every characteristic.
On finite-coordinate vectors, sorting distinct indices shows
$\mathcal A_V S(m)=m$ for alternating $m$; repeated indices contribute zero.
Density proves the same identity everywhere. Hence $R=\mathcal A_V S$ retracts all
multilinear forms onto alternating forms, while $S|_{\Alt}$ is a bounded right inverse
of $\mathcal A_V$. Transport through $j_V$ gives both operators on the original $V$.

If $E$ is of countable type, a lift of $A$ is
$R_E\bigl(m\circ(u_1,\ldots,u_k)\bigr)$.
If $D$ is of countable type, use instead
$\mathcal A_E\bigl(S_D(m)\circ(u_1,\ldots,u_k)\bigr)$.
Both are bounded multilinear in the $u_i$, linear in $m$, and have diagonal $A(u)(m)$;
the second identity uses that antisymmetrization commutes with simultaneous precomposition.
Proposition~\ref{prop:homogeneous-lift} concludes the proof.
\end{proof}


\subsection{Continuous algebraic polynomials on \texorpdfstring{$c_0$}{c0}}
The reflection theorem starts with bounded multilinear coefficients in an ambient space. For one
homogeneous polynomial, interpolation supplies a lift even when only the diagonal is known to be
continuous. The resulting bound is degreewise; it gives no assertion about a common radius across
an infinite sequence of degrees.

\begin{proposition}[continuous algebraic polynomials on $c_0$]\label{fam:prop:c0-algebraic}
Let $K$ be nontrivially normed and nonarchimedean, let $Z$ be complete and nonarchimedean, and let
$n\ge1$. Suppose $p\colon c_0(I,K)\to Z$ is continuous and is an algebraic homogeneous polynomial:
there exists a $K$\nobreakdash-multilinear map $b\colon c_0(I,K)^n\to Z$, not assumed continuous,
with $p(x)=b(x,\ldots,x)$. Then $p$ has a bounded $n$\nobreakdash-linear lift $q$. More precisely,
there is $C_{K,n}<\infty$, independent of $I,Z,p$, such that
\[
 \norm q\le C_{K,n}\norm p_{\mathrm{diag}},\qquad
 \norm p_{\mathrm{diag}}:=\inf\{D\ge0:\norm{p(x)}\le D\norm x^n\text{ for all }x\}.
\]
No completeness assumption on $K$ is needed.
\end{proposition}
\begin{proof}
Continuity at zero and $p(ax)=a^np(x)$ imply $\norm p_{\mathrm{diag}}<\infty$.
Explicitly, choose $t\in K$ with $0<\abs t<1$ and $\delta>0$ such that $\norm{p(x)}\le1$
when $\norm x<\delta$. For $x\ne0$, choose an integer $m$ with
$\abs t\delta\le\abs{t^m}\norm x<\delta$. Homogeneity gives
$\norm{p(x)}\le(\abs t\delta)^{-n}\norm x^n$.

Fix a total order on $I$. On each finite coordinate subspace, the algebraic formula for $p$ is an
ordinary homogeneous polynomial of degree $n$. Polynomial uniqueness over the infinite field $K$
makes its coefficients compatible as the finite coordinate set varies. Denote the coefficient of
$x^\alpha$ by $c_\alpha$, where $\alpha\colon I\to\N$ is finitely supported and $\sum_i\alpha_i=n$.
Its support has at most $n$ elements.

Choose distinct $a_0,\ldots,a_n$ in the valuation ring of $K$, and let $L_n\ge1$ bound the absolute
values of all entries of the inverse Vandermonde matrix $(a_i^j)_{0\le i,j\le n}$.
If the support of $\alpha$ has $d\le n$ elements, interpolate the restriction of $p$ to those
coordinates on the grid $\{a_0,\ldots,a_n\}^d$. Each grid vector has norm at most one. Applying the
inverse Vandermonde matrix in each coordinate and using the nonarchimedean inequality in $Z$ yields
\[
 \norm{c_\alpha}\le L_n^d\norm p_{\mathrm{diag}}
 \le L_n^n\norm p_{\mathrm{diag}}.
\]
For a nondecreasing word $i_1\le\cdots\le i_n$, write $c_{i_1,\ldots,i_n}$ for the coefficient
with those multiplicities. Define
\[
 q(x_1,\ldots,x_n)=
 \sum_{i_1\le\cdots\le i_n}
 c_{i_1,\ldots,i_n}(x_1)_{i_1}\cdots(x_n)_{i_n}.
\]
The coefficients are uniformly bounded and each input is a null family, so the summands tend to zero
outside finite sets of words. Completeness and the nonarchimedean inequality give convergence,
multilinearity, and the bound $\norm q\le L_n^n\norm p_{\mathrm{diag}}$.
On finite supports its diagonal is $p$; density and continuity extend the identity to all of $c_0(I,K)$.
Thus $C_{K,n}=L_n^n$ suffices.
\end{proof}

This bound is uniform in the number of coordinates, but the interpolation nodes and $L_n$ depend on
the degree. No exponential growth bound for $C_{K,n}$ is asserted. Consequently this proposition does
not justify replacing lifting norms by diagonal norms in an arbitrary infinite analytic expansion.
Theorem~\ref{thm:summable-parameters}(1) obtains its stronger, radius-preserving conclusion from the
already bounded ambient coefficients.


\subsection{A nonspherical target norm with analytic precomposition}
\begin{remark}\label{rem:nonspherical-analytic-target}
Spherical completeness of the given target norm is not necessary for analyticity, even in degrees
at least the characteristic. Let $K=\F_p((t))$ with $\abs t=1/2$ and let $F=c_0(\N,K)$, the space
of sequences tending to zero. Define
\[
 w_j=1+\frac1{j+2},\qquad \norm{x}_w=\sup_{j\ge0}w_j\abs{x_j}.
\]
This is a nonarchimedean norm satisfying
\[
 \norm{x}_\infty\le\norm{x}_w\le\frac32\norm{x}_\infty.
\]
Thus $F$ is Banach for either norm. The usual supremum norm has nonzero values in $2^\Z$,
since every nonzero null sequence attains its maximum absolute value. The discrete-distance
argument in the proof of Corollary~\ref{cor:discrete-targets} therefore makes the supremum norm
spherically complete. Theorem~\ref{thm:spherical} gives analytic precomposition in every finite
degree for that norm, for all normed $E,E'$. Equivalent norms preserve the bounded maps and their
analytic power-series expansions: the norms of coefficients change by a fixed factor and convergence
is unchanged. Thus precomposition remains analytic for $\norm\cdot_w$, also directly by
Corollary~\ref{cor:discrete-targets}.

Nevertheless, the weighted norm is not spherically complete. Write $e_j$ for the coordinate vectors,
set $a_n=\sum_{j<n}e_j$, and consider the closed balls
\[
 B_n=\overline B_w(a_n,w_n).
\]
They are nested because $w_{n+1}<w_n$ and $\norm{a_{n+1}-a_n}_w=w_n$; the nonarchimedean
inequality gives $B_{n+1}\subseteq B_n$. If $x$ belonged to all of them, then, for each $j$ and
any $n>j$,
\[
 \abs{x_j-1}\le\frac{w_n}{w_j}<1.
\]
Hence $\abs{x_j}=1$ for every $j$, contradicting $x_j\to0$. Their intersection is empty.
This example separates spherical completeness of a particular norm from analyticity, which is
invariant under equivalent norms. It does not exclude a characterization using an equivalent
spherically complete norm: this target has exactly such a norm.
\end{remark}


\begin{thebibliography}{99}
\bibitem{BGN} W.~Bertram, H.~Gl\"ockner and K.-H.~Neeb,
\emph{Differential calculus, manifolds and Lie groups over arbitrary infinite fields},
arXiv:math/0303300 (2003); published as \emph{Differential calculus over general base fields and rings}, Expo.\ Math.\ \textbf{22} (2004), 213--282,\\
{\small\url{https://arxiv.org/abs/math/0303300}}.

\bibitem{Bou} N.~Bourbaki, \emph{Vari\'et\'es diff\'erentielles et analytiques. Fascicule de r\'esultats}, Hermann, Paris,
1967 (\S\S1--7) and 1971 (\S\S8--15); reprinted by Springer, 2007.

\bibitem{G1} S.~Gou\"ezel, \emph{In characteristic $0$, precomposition by a linear map is analytic on the space of
alternating maps}, pull request \#43338 to Mathlib, merged 30 September 2026,\\ {\small\url{https://github.com/leanprover-community/mathlib4/pull/43338}}.

\bibitem{G2} S.~Gou\"ezel, \emph{The bundle of continuous alternating maps is $C^n$}, draft pull request \#43234 to Mathlib,\\
{\small\url{https://github.com/leanprover-community/mathlib4/pull/43234}}.

\bibitem{In} A.~W.~Ingleton, \emph{The Hahn--Banach theorem for non-Archimedean-valued fields}, Proc.\ Cambridge Philos.\
Soc.\ \textbf{48} (1952), 41--45.

\bibitem{La} S.~Lang, \emph{Fundamentals of Differential Geometry}, Graduate Texts in Mathematics 191, Springer, 1999.

\bibitem{LeanRef} The Lean Development Team, \emph{The Lean Language Reference},
\S4.2, Propositions, accessed 6 October 2026,\\
{\small\url{https://lean-lang.org/doc/reference/latest/The-Type-System/Propositions/}}.

\bibitem{docs} The mathlib Community, \emph{Mathlib documentation}, \texttt{docs\#ContDiffWithinAt},\\
{\footnotesize\url{https://leanprover-community.github.io/mathlib4_docs/find/?pattern=ContDiffWithinAt#doc}}.

\bibitem{LeanP} J.~McCarthy, \emph{Lean formalization of alternating precomposition in positive characteristic},
commit \texttt{ca36ba5},\\
{\small\url{https://github.com/Deicyde/alternating-analytic-characteristic-p}}.

\bibitem{PR} J.~McCarthy, \emph{Smoothness reflects along a linear isometry with closed range},\\ draft pull request \#43548
to Mathlib,\\ {\small\url{https://github.com/leanprover-community/mathlib4/pull/43548}}.

\bibitem{VdP} M.~van der Put, \emph{Espaces de Banach non archim\'ediens}, Bull.\ Soc.\ Math.\ France
\textbf{97} (1969), 309--320,\\
{\small\url{https://www.numdam.org/article/BSMF_1969__97__309_0.pdf}}.

\bibitem{Ra} F.~P.~Ramsey, \emph{On a problem of formal logic}, Proc.\ London Math.\ Soc.\ (2) \textbf{30} (1930), 264--286.

\bibitem{vR} A.~C.~M.~van Rooij, \emph{Non-Archimedean Functional Analysis}, Marcel Dekker, 1978.

\bibitem{Schi} W.~H.~Schikhof, \emph{Banach spaces over nonarchimedean valued fields}, Topology Proc.\
\textbf{24} (1999), 547--581,\\
{\small\url{https://topology.nipissingu.ca/tp/reprints/v24/tp24224.pdf}}.

\bibitem{Sch} P.~Schneider, \emph{Nonarchimedean Functional Analysis}, Springer Monographs in Mathematics,
Springer, Berlin, 2002, \href{https://doi.org/10.1007/978-3-662-04728-6}{doi:10.1007/978-3-662-04728-6}.
Online version of 25 October 2005,\\
{\small\url{https://ivv5hpp.uni-muenster.de/u/pschnei/publ/lectnotes/nfa.pdf}}.

\bibitem{Se} J.-P.~Serre, \emph{Endomorphismes compl\`etement continus des espaces de Banach $p$\nobreakdash-adiques}, Publ.\ Math.\
IH\'ES \textbf{12} (1962), 69--85.

\bibitem{Yuan} Y.~Yuan, \emph{Formalization of non-Archimedean functional analysis 1: spherically complete spaces},\\
arXiv:2601.21734v3 (2026), accepted for publication in the Journal of Symbolic Logic,\\
{\small\url{https://arxiv.org/abs/2601.21734v3}}.

\end{thebibliography}
\end{document}
